\subsection{An affirmative integral result for ring class fields}

\begin{result}{E.2.2}{One labelled CM elliptic curve separates the places of its ring class field}{Proved}
Let $E/\overline\Q$ have complex multiplication by an order $\mathcal O$ in an imaginary quadratic field $K$, and put $k=K(j(E))$. Choose an elliptic model of that $j$-invariant over $k$. The integral Betti diagram consisting of this elliptic curve and its labelled endomorphisms distinguishes all complex places of $k$. In fact, one nonrational endomorphism suffices. Thus the injectivity proposed in the number-field paragraph of 15:223 holds for every elliptic ring class field with this particular finite test.
\end{result}
\begin{proof}
First all the geometric endomorphisms of $E$ are defined over $k$. The action on a nonzero invariant differential gives an injective homomorphism
\[
 \End_{\overline\Q}(E)\longrightarrow\overline\Q,
 \qquad u^*\omega=d(u)\omega,
\]
whose image is the specified copy of $\mathcal O\subset K$. Injectivity follows because a nonzero endomorphism of an elliptic curve in characteristic zero is a separable isogeny and has nonzero differential. For $\sigma\in\Gal(\overline\Q/k)$ one has $d(u^\sigma)=\sigma(d(u))=d(u)$; hence $u^\sigma=u$. The coefficients defining $u$ lie in a finite algebraic extension, and this invariance descends the morphism to $k$.

Represent any two complex places by embeddings $v,w:k\hookrightarrow\C$ whose restrictions to $K$ are the same fixed embedding. This is possible by replacing either representative by its complex conjugate. Choose $\alpha\in\mathcal O-\Z$ nonreal, and label the corresponding endomorphism $u$. Suppose an integral Betti isomorphism on the displayed diagram exists. In covariant homology its component is
\[
 A:H_1(E_v,\Z)\xrightarrow{\sim}H_1(E_w,\Z),
 \qquad A u_{v,*}=u_{w,*}A.
\]
The exponential uniformization identifies $H_1(E_v,\R)$ with the real tangent space of $E_v$ by extending its period lattice; similarly for $w$. Write $T_v=u_{v,*}$ on that real vector space. Since $u$ acts on the complex tangent space as the scalar $v(\alpha)$, its complex structure is exactly
\[
 J_v=\frac{T_v-\operatorname{Re}(v(\alpha))I}
                 {\operatorname{Im}(v(\alpha))}.
\]
The denominator is nonzero, and $v(\alpha)=w(\alpha)$ because the embeddings agree on $K$. Thus the real extension of $A$ satisfies $A J_v=J_w A$. It is a complex-linear isomorphism of tangent spaces carrying the period lattice onto the period lattice. Quotienting gives an isomorphism of complex elliptic tori $E_v\simeq E_w$, hence an isomorphism of elliptic curves. Therefore $v(j(E))=w(j(E))$. As $k=K(j(E))$, the embeddings agree on all of $k$, and the two original complex places coincide. Dualizing the free lattices proves the same assertion for contravariant $H^1$.

This proof explicitly reconstructs the required complex structure from a labelled algebraic endomorphism and its real action. It does not add a Hodge filtration to the tested data. In particular, the exact map from a natural integral motivic comparison to the contradiction is its restriction to $E$ and $u$ followed by the displayed formula for $J$.
\end{proof}

For a completely explicit instance take
\[
 K=\Q(\sqrt{-5}),\qquad
 j_+=632000+282880\sqrt5,\qquad
 k=\Q(i,\sqrt5)=K(j_+),
\]
and the smooth elliptic curve
\[
 E:\quad y^2=x^3-\frac{3j_+}{j_+-1728}x
                         -\frac{2j_+}{j_+-1728}.
 \tag{E.2.2a}
\]
Its $j$-invariant is $j_+$. The CM identification
$\End(E_{\overline\Q})=\Z[\sqrt{-5}]$ and the equality of the Hilbert class field with $\Q(\sqrt{-5},i)$ are explicitly recorded in \cite[Section 5, preprint p.~11, immediately before (19)]{GenCampagnaPengo2022}. The two complex places have representatives
\[
 v(i)=i,\quad v(\sqrt5)=\sqrt5;\qquad
 w(i)=-i,\quad w(\sqrt5)=-\sqrt5.
\]
Both fix $\sqrt{-5}=i\sqrt5$, whereas $w(j_+)=j_-=632000-282880\sqrt5\ne j_+$. The proof above therefore excludes an integral comparison commuting with the labelled multiplication by $\sqrt{-5}$.

The ideal-lattice obstruction in this example can also be exhibited directly. The two classes of the maximal order $\mathcal O=\Z[\sqrt{-5}]$ are represented by
\[
 \mathcal O,\qquad I=(2,1+\sqrt{-5}).
\]
The norm of $I$ is $2$, $I^2=(2)$, and $I$ is nonprincipal, since an element generating an ideal of norm $2$ would give integers $a,b$ with $a^2+5b^2=2$. The usual quadratic Minkowski bound is $(2/\pi)\sqrt{20}<3$, so every ideal class contains an integral ideal of norm $1$ or $2$; the unique prime ideal of norm $2$ is $I$. This proves that there are precisely these two classes. Analytic uniformization identifies CM elliptic curves of the fixed type with homothety classes of these lattices, and an $\mathcal O$-linear integral lattice isomorphism is multiplication by an element of $K^\times$. The two distinct values $j_\pm$ therefore give the two nonisomorphic classes, in some order. Their real scalar extensions are both free rank-one modules over $K\otimes_\Q\R\simeq\C$. Consequently this example also verifies directly that a real comparison on this labelled endomorphism diagram need not preserve its integral lattice.

The scope of E.2.2 is $k=K(j(E))$. For a larger number field, two embeddings agreeing on $K(j(E))$ give exactly the same labelled elliptic test. This is an equality of the two base-change diagrams under the restriction map of embeddings; it records the precise remaining issue for extending this construction to all number fields. The theorem of Charles, \cite[Theorem 2]{Charles2009}, supplies another finite cohomological test when the automorphism acts trivially on one of two specified imaginary quadratic fields and nontrivially on the other. That hypothesis does not include every prescribed pair of places, and no such extension is asserted here.

\begin{result}{E.2.3}{Descent of a finite separating diagram and the exact transporter}{Proved}
Let $k'/k$ be a finite extension of number fields. A finite integral Betti correspondence diagram over $k'$ that distinguishes all complex places of $k'$ gives a finite such diagram over $k$ that distinguishes all complex places of $k$. Its objects lie in the tensor diagram generated by one finite labelled direct sum. In particular, every subfield of an elliptic ring class field has a finite separating test.

For the entire integral correspondence diagram, there is also an exact description. Fix the usual Betti functor $\omega_{\Q,\Z}$ over $\Q$. For a number field $k$, let $A_k=h^0(\operatorname{Spec}k)$, and denote by $e_u$ the primitive idempotent of $\omega_{\Q,\Z}(A_k)=\Z^{\operatorname{Hom}(k,\C)}$ indexed by $u$. For embeddings $v,w:k\hookrightarrow\C$, restriction from $k$ to the base-changed $\Q$-diagram gives the bijection
\[
 \Isom^\otimes(\omega_{v,\Z},\omega_{w,\Z})
 \xrightarrow{\ \sim\ }
 \bigl\{g\in\Aut^\otimes(\omega_{\Q,\Z}):
                  g_{A_k}(e_v)=e_w\bigr\}.
 \tag{E.2.3a}
\]
Here, as throughout this section, the integral source is the concrete graded cohomology diagram with integral algebraic correspondences and the displayed tensor operations. The statement does not assign a free integral fibre functor to a $\Q$-linear category.
\end{result}
\begin{proof}
Write $A=h^0(\operatorname{Spec}k')$ as an object over $k$. For a $k'$-variety $X$, write $X_{[k]}$ for the same scheme with structural morphism composed with $\operatorname{Spec}k'\to\operatorname{Spec}k$. This specifies the operation on schemes exactly. It preserves smoothness and projectivity because $k'/k$ is finite separable. The same restriction construction and its sum over embeddings appear in \cite[Example 0.1.1, p.~314, and (0.5.1), p.~316]{Deligne1979}; the integral maps needed here are checked below. At an embedding $v:k\hookrightarrow\C$ there is a canonical decomposition
\[
 H^*(X_{[k],v},\Z)/\mathrm{tors}
   =\bigoplus_{u:k'\hookrightarrow\C,\ u|_k=v}
                   H^*(X_u,\Z)/\mathrm{tors}.
 \tag{E.2.3b}
\]
The structure map $X_{[k]}\to\operatorname{Spec}k'$ makes this an $A$-module: pull back degree-zero cohomology and then use cup product. Its primitive idempotent $e_u$ projects exactly to the $u$-summand.

Given the finite diagram over $k'$, include over $k$ the object $A$, its multiplication and unit, the objects $X_{[k]}$, their $A$-module structures, and the specified correspondences. Each $k'$-correspondence is a cycle on $X\times_{k'}Y$. This product is the open and closed diagonal component of $X_{[k]}\times_kY_{[k]}$, so the same cycle, with the same codimension, is a correspondence over $k$. The relative product is selected on cohomology by the algebraic diagonal idempotent
\[
 \Delta_A=\sum_u e_u\otimes e_u.
 \tag{E.2.3c}
\]
In particular all original product operations are recovered on the corresponding diagonal summands.

Suppose a natural integral comparison exists between two embeddings $v,w$ of $k$ on this descended diagram. Its component on $A$ is a unital ring isomorphism between finite products of $\Z$, hence permutes their primitive idempotents. Fix one extension $u$ of $v$ to $k'$, and let $u'$ be the extension of $w$ selected by the image of $e_u$. Naturality for the $A$-actions carries every $u$-summand in (E.2.3b) onto the $u'$-summand. The correspondence maps restrict to exactly the original maps over $u,u'$. The permutation of idempotents preserves (E.2.3c), so these restrictions preserve the tensor operations as well. They yield a comparison on the original finite diagram. Its separating property gives $u'=u$ or $u'=\overline u$ as complex embeddings, whence $w=v$ or $w=\overline v$ by restriction. This proves descent.

Only finitely many varieties, cycles and maps were added. If the original objects are $h(X_i)(r_i)$, one may use the labelled direct sum
\[
 A\ \oplus\ \bigoplus_i h(X_{i,[k]})(r_i)
\]
and the finitely many Tate objects used by the original maps as one generator; the summand inclusions and projections, the multiplication of $A$, the module maps and the original correspondences specify the required diagram in its tensor constructions. Grading and the usual K\"unneth signs are retained. On torsion-free quotients the canonical K\"unneth map is an isomorphism: the integral K\"unneth exact sequence for these finite CW complexes splits as a sequence of abelian groups, and its $\operatorname{Tor}$ terms are torsion. Thus no change of coefficient ring is used in this construction. Applying it to E.2.2 proves the ring-class-subfield assertion.

The same argument also proves the following stronger form of the descent step, used below: it suffices that a comparison on the original diagram between extensions $u,u'$ forces their \emph{restrictions to $k$} to agree up to conjugation. No separation of two extensions with the same restriction to $k$ is needed. Indeed the preceding construction always supplies a comparison on one such pair of extension blocks, and the conclusion concerns exactly their restrictions $v,w$.

We prove (E.2.3a), including both inverses. A comparison $\eta$ over $k$ restricts along base change $M\mapsto M_k$ from $\Q$ to $k$, since the two complex realizations of $M_k$ are canonically those of $M$. This defines $g_M$. The diagonal section
\[
 \operatorname{Spec}k\longrightarrow
              \operatorname{Spec}k\times_{\Q}\operatorname{Spec}k
\]
is a morphism over $k$. On cohomology, at $v$ and $w$, it gives evaluation at the respective idempotents. Naturality says
\[
 \operatorname{ev}_w\,g_{A_k}=\operatorname{ev}_v,
\]
which is equivalent to $g_{A_k}(e_v)=e_w$, because $g_{A_k}$ permutes the primitive idempotents.

Conversely take such a $g$. For $X/k$, its action on the cohomology of $X_{[\Q]}$ commutes with the $A_k$-module structure and therefore carries the $v$-summand onto the $w$-summand. Define $\eta_X$ to be this restriction. The open and closed component construction just proved gives naturality for every $k$-correspondence. The fixed diagonal (E.2.3c) gives the tensor compatibility. The unit, duality, cup products and Tate maps restrict to their original maps on each summand.

To check the two inverses, $X/k$ is canonically the diagonal open and closed summand of $X_{[\Q]}\times_{\Q}k$; naturality for its inclusion and projection recovers the initial $\eta_X$. In the other direction, for an object $M/\Q$ the underlying $\Q$-object of $M_k$ is $M\otimes A_k$. Its action is $g_M\otimes g_{A_k}$, and restriction from the $e_v$ block to the $e_w$ block is exactly $g_M$. Thus both composites are the identity, proving the stated bijection.
\end{proof}

For example, E.2.3 applies to the real quadratic subfield $\Q(\sqrt5)$ of the explicit field $\Q(i,\sqrt5)$ in E.2.2. One takes the degree-two finite \emph{\'etale} scheme $\operatorname{Spec}\Q(i,\sqrt5)$ over $\Q(\sqrt5)$, the displayed CM elliptic curve regarded as a scheme over $\Q(\sqrt5)$, its module action and its labelled multiplication by $\sqrt{-5}$. Any comparison between the two real embeddings would select two complex components above them. Those components have the distinct real invariants $j_+,j_-$ and hence cannot correspond by either equality or complex conjugation in E.2.2. This gives a numerical finite test with all the descent maps specified.

The transporter also isolates the remaining universal assertion. Restrict a tensor automorphism of $\omega_{\Q,\Z}$ to the finite \emph{\'etale} schemes. On every finite Galois extension $L/\Q$, naturality for the deck maps makes its permutation of $\operatorname{Hom}(L,\C)$ a left translation; these translations are compatible under restriction. They define an element of $\Gal(\Qb/\Q)$. The identity and ordinary complex conjugation both occur, the latter by complex conjugation on every complex variety and its singular cohomology. Formula (E.2.3a) proves that injectivity for complex places of \emph{every} number field is equivalent to the exact equality
\[
 \operatorname{im}\!\left(
   \Aut^\otimes(\omega_{\Q,\Z})\longrightarrow\Gal(\Qb/\Q)
                 \right)=\{1,c\}.
 \tag{E.2.3d}
\]
Indeed a Galois element outside $\{1,c\}$ differs from both on some finite Galois subextension, and then (E.2.3a) produces two different complex places with a comparison. Conversely a comparison at different complex places gives an element outside $\{1,c\}$. The equality in (E.2.3d) is the unresolved step; the bijection and the equivalence have just been proved.

\begin{result}{E.2.4}{Generalized-dihedral fields covered by the finite CM test}{Proved}
Let $L/\Q$ be a finite Galois extension containing an imaginary quadratic field $K$, and suppose $B=\Gal(L/K)$ is abelian. Let $c$ be actual complex conjugation on $L$. If $cbc^{-1}=b^{-1}$ for every $b\in B$, then every subfield of $L$ has a finite integral Betti diagram distinguishing its complex places. This includes every cubic field of negative discriminant.

The ring-class-field construction has an exact restriction: a totally real Galois subfield of any elliptic ring class field has elementary abelian $2$-group as its Galois group. In particular, a totally real number field with nonabelian-dihedral normal closure is not a subfield of an elliptic ring class field. This latter assertion identifies the limit of the particular CM construction, and makes no negative claim about other motivic tests for those fields.
\end{result}
\begin{proof}
Since $K$ is imaginary quadratic, $c\notin B$ and $c^2=1$. Thus
\[
 G:=\Gal(L/\Q)=B\rtimes\langle c\rangle.
\]
The transfer $\operatorname{Ver}:G^{\mathrm{ab}}\to B$ is trivial. Indeed, using the two coset representatives $1,c$, its value at $b\in B$ is $b(cbc^{-1})=1$, and its value at $c$ is $c^2=1$. These elements generate $G$, and transfer is a homomorphism.

Choose a positive integer $f$ divisible by every rational prime ramified in $L/\Q$ and large enough that $f\mathcal O_K$ is a modulus for the Artin map of $L/K$. Such an integer exists: the conductor is an integral ideal of $K$, and its local valuations are bounded by those of $f\mathcal O_K$ after increasing the finitely many exponents in $f$. Write $I_K(f)$ for fractional $\mathcal O_K$-ideals prime to $f$. The Artin homomorphism
\[
 \theta:I_K(f)\twoheadrightarrow B
\]
kills the ray principal group $P_{K,1}(f\mathcal O_K)$. Functoriality under extension of ideals gives the commutative square with upper row the Artin map from rational ideals prime to $f$ into $G^{\mathrm{ab}}$, lower row $\theta$, left map $\mathfrak a\mapsto\mathfrak a\mathcal O_K$, and right map $\operatorname{Ver}$. This is the transfer compatibility in \cite[Chapter V, Section 3, Proposition 3.18 and the adjoining Artin-map diagram, p.~162]{GenMilneCFT2020}. Therefore
\[
 \theta(a\mathcal O_K)=1
       \qquad(a\in\Z,\ \gcd(a,f)=1).
 \tag{E.2.4a}
\]
This proves actual vanishing; it uses the order-two complex-conjugation lift, in addition to the inversion action.

Let $\mathcal O_f=\Z+f\mathcal O_K$. The ideal-class map for this order is
\[
 \operatorname{Pic}(\mathcal O_f)
    \simeq I_K(f)/P_{K,\Z}(f),
 \tag{E.2.4b}
\]
where $P_{K,\Z}(f)$ is generated by the ideals $\alpha\mathcal O_K$ with $\alpha\in\mathcal O_K$ and $\alpha\equiv a\pmod{f\mathcal O_K}$ for an integer $a$ prime to $f$. The isomorphism sends an invertible order ideal prime to the conductor to its extension to $\mathcal O_K$; contraction is its inverse. The principal-subgroup calculation is \cite[Proposition 7.22, pp.~145--146]{GenCox1989}. To check that $\theta$ factors through this quotient, choose $\gamma\in\mathcal O_K$ with $a\gamma\equiv1\pmod{f\mathcal O_K}$. Then $\alpha\gamma\equiv1$ as well. The ray-principal kernel gives $\theta((\alpha))=\theta((a))$, and (E.2.4a) makes this value $1$. Consequently (E.2.4b) gives the explicit surjection
\[
 \operatorname{Pic}(\mathcal O_f)\twoheadrightarrow\Gal(L/K),
 \qquad [\mathfrak a]\longmapsto
         \operatorname{Art}_{L/K}(\mathfrak a\mathcal O_K).
 \tag{E.2.4c}
\]
Under the class-field correspondence this is restriction
$\Gal(H_f/K)\twoheadrightarrow\Gal(L/K)$, and hence $L\subset H_f$, where $H_f$ is the ring class field of $\mathcal O_f$. This proves the conductor passage and the required field inclusion. It also recovers the equivalence in \cite[Theorem 9.18 and its proof, pp.~190--192]{GenCox1989}.

The complex multiplication theorem gives $H_f=K(j(E_f))$ for a CM elliptic curve with endomorphism order $\mathcal O_f$; the exact assertion is \cite[Theorem 11.1, p.~220]{GenCox1989}. Its $j$-invariant is algebraic, so an elliptic model with that invariant exists over $H_f$. Apply E.2.2 to this model and then E.2.3 along $H_f/k$ for each subfield $k\subset L$. This proves the first assertion with the explicit maps (E.2.3b)--(E.2.3c) and (E.2.4c).

For a cubic field of negative discriminant, its normal closure has Galois group $S_3$: irreducibility gives a transitive subgroup of $S_3$, and its negative discriminant is not a rational square. The fixed field of $A_3$ is the imaginary quadratic discriminant field. Actual complex conjugation is a transposition, since there is one real embedding and one conjugate pair. Its conjugation action on $A_3\simeq C_3$ is inversion. The construction therefore applies.

Finally $\Gal(H_f/\Q)=B_f\rtimes\langle c\rangle$ with $B_f$ abelian and $c$ acting by inversion. One may obtain this directly from the Artin description: complex conjugation preserves $P_{K,\Z}(f)$, and an ideal class satisfies $[\mathfrak a][\overline{\mathfrak a}]=[N\mathfrak a]=1$. This is also \cite[Lemma 9.3, pp.~180--181]{GenCox1989}. If $T/\Q$ is a totally real Galois subfield of $H_f$, restriction sends $c$ to $1$. The relation $cbc^{-1}=b^{-1}$ then forces the image of every $b$ to have square $1$, and these commuting images generate $\Gal(T/\Q)$. It is an elementary abelian $2$-group.

If a totally real nonnormal field $k$ were contained in $H_f$, its normal closure would also be contained, because $H_f/\Q$ is Galois; that closure is still totally real. The preceding quotient computation excludes a nonabelian-dihedral normal closure. For example $x^3-4x+1$ is irreducible over $\Q$ by the rational-root test and has positive nonsquare discriminant $229$. Its three roots are real, so the field generated by any root has a totally real $S_3$ normal closure and cannot be contained in any elliptic ring class field. No assertion about the existence of other finite separating diagrams for this field follows from that exclusion.
\end{proof}

\begin{result}{E.2.5}{A finite polarized CM test for every totally real field}{Proved}
Every totally real number field has a finite labelled algebraic-correspondence diagram whose Betti realizations at its distinct embeddings are not isomorphic, even over $\R$. Every CM number field has such a finite test for its complex places. In particular, together with E.2.4, every cubic number field has a finite integral test distinguishing its complex places. The totally real cubic generated by a root of $x^3-4x+1$ has the explicit CM threefold construction below.
\end{result}
\begin{proof}
We first give the real linear map that extracts a CM type from the labelled data. Let $B/\C$ be an abelian variety of dimension $n$ with multiplication by a CM field $K$ of degree $2n$, and let $Q$ be a polarization whose Rosati involution on $K$ is complex conjugation. Put $V=H_1(B,\R)$. This is free of rank one over
\[
 K\otimes_\Q\R=\prod_{\tau:K^+\hookrightarrow\R}\C.
\]
The labelled action gives real two-dimensional summands $V_\tau$. They are orthogonal for $Q$: the idempotents of $K^+\otimes\R$ are self-adjoint. Consequently the restriction of $Q$ to each summand is nondegenerate. Choose once for each factor one of the two identifications with $\C$, and let $I_\tau$ be multiplication by $i$ in that identification. The analytic complex structure $J$ commutes with the $K$-action, and therefore on each factor
\[
 J|_{V_\tau}=s(\tau)I_\tau,\qquad s(\tau)\in\{1,-1\}.
 \tag{E.2.5a}
\]
Indeed an operator commuting with the complex scalars is a complex scalar, and its square is $-1$. With the common homological convention $Q(x,Jx)>0$ for nonzero $x$, the sign of $Q(x,I_\tau x)$ is exactly $s(\tau)$. These signs are the choices of embeddings in the CM type.

Now a real comparison $T$ commuting with every labelled element of $K$ preserves the individual summands and the $I_\tau$. Preservation of the one labelled Tate-valued polarization means
\[
 Q'(Tx,Ty)=\lambda Q(x,y)
\]
for one scalar $\lambda\in\R^\times$ common to all factors. For an integral comparison this scalar is $1$ or $-1$. Substituting $y=I_\tau x$ shows
\[
 s'(\tau)=\operatorname{sign}(\lambda)s(\tau)
              \quad\hbox{for every }\tau.
 \tag{E.2.5b}
\]
Thus the CM types are identical or everywhere complex conjugate. This conclusion uses only the labelled endomorphisms and the algebraic polarization in the Betti diagram. The Hodge filtration is recovered in this calculation, rather than added to the tested data. Pointwise preservation of the labels matters: the comparison is not permitted to replace the specified action of $K$ by an arbitrary automorphism of that field.

For the contravariant cohomology convention, the passage is exact: $H^1(B,\Z)=\operatorname{Hom}(H_1(B,\Z),\Z)$, and an isomorphism on $H^1$ yields $T$ by taking the inverse of its dual. The first Chern class of the polarization is the algebraic map $\Z(-1)\to\bigwedge^2 H^1(B,\Z)$; dualizing gives the Tate-valued alternating form on $H_1$. Naturality gives the displayed similitude equation, with the inverse scalar on the dual Tate line. It is still one common nonzero real scalar, or one common sign integrally. Thus the proof applies to the precise cohomology diagram stipulated in this section.

Let now $F\subset\R$ be totally real of degree $n\ge3$, put $K=F(i)$, and enumerate the real embeddings $\tau_1,\ldots,\tau_n$ with $\tau_1$ the given inclusion. Take the CM type $\Phi$ that sends $i$ to $i$ at $\tau_1$ and to $-i$ at all other embeddings. To construct a polarized torus of this type, choose $b\in\mathcal O_F$ with
\[
 \tau_1(b)>0,\qquad \tau_j(b)<0\quad(j>1).
\]
Such an element exists by weak approximation, followed by multiplication by a positive integer. Put $\xi=ib$, choose the full lattice $R=\mathcal O_F[i]\subset K$, and set
\[
 B_\Phi(\C)=\C^\Phi/\Phi(R),\qquad
 Q(x,y)=\operatorname{Tr}_{K/\Q}(\xi x\overline y)
             \quad(x,y\in R).
 \tag{E.2.5c}
\]
The form is integral, alternating and nondegenerate. At every $\phi\in\Phi$ write $\phi(\xi)=i d_\phi$ with $d_\phi>0$. Its real extension is
\[
 Q(z,w)=2\sum_{\phi\in\Phi}
                d_\phi\operatorname{Im}(\overline{z_\phi}w_\phi),
 \qquad
 Q(z,Jz)=2\sum_\phi d_\phi|z_\phi|^2>0.
\]
These are the Riemann bilinear relations, so (E.2.5c) is a polarized abelian variety; multiplication by $R$ supplies the labelled endomorphisms and the formula for $Q$ gives the required Rosati involution. This is the explicit analytic construction in \cite[Example 2.11 and Proposition 3.17]{GenMilneCM2020}.

The CM abelian variety has a model over $\Qb$, and its endomorphisms and polarization descend there: \cite[Proposition 7.9 and Corollary 7.10, pp.~54--55]{GenMilneCM2020} prove the equivalence under algebraically closed base extension. Only finitely many endomorphisms are needed, namely a $\Z$-basis of $R$. Choose a finite number field $L$ containing their coefficients, those of the abelian variety and polarization, and the normal closure of $K$. This gives an actual finite algebraic diagram over $L$. The polarization can equivalently be represented by its homomorphism to the dual abelian variety and the algebraic Poincar\'e pairing, so its Tate-valued alternating form is among the algebraic maps of the diagram.

For an embedding $u:L\hookrightarrow\C$, the labelled type is $u\Phi$. This follows directly on the tangent space: the endomorphism matrices are conjugated coefficientwise by $u$, so each eigencharacter $\phi$ becomes $u\circ\phi$. Define the type traces
\[
 t_\Phi(a)=\sum_{\phi\in\Phi}\phi(a)\qquad(a\in K).
\]
For $a\in F$ the choices in $\Phi$ give the exact formulas
\[
 t_\Phi(i)=(2-n)i,
 \qquad
 t_\Phi(ia)=i\bigl(2a-\operatorname{Tr}_{F/\Q}(a)\bigr).
 \tag{E.2.5d}
\]
The first is nonzero because $n\ge3$, and the traces in (E.2.5d) generate $F(i)$ over $\Q$. Thus this is also a direct calculation of the reflex field, without an uncomputed field-of-moduli identification.

Suppose embeddings $u,u'$ of $L$ give a real Betti comparison on this diagram. Equation (E.2.5b) implies $u'\Phi=u\Phi$ or $u'\Phi=\overline{u\Phi}$. Summing the eigencharacters gives, in the two cases,
\[
 u'(t_\Phi(a))=u(t_\Phi(a))
 \quad\hbox{or}\quad
 u'(t_\Phi(a))=\overline{u(t_\Phi(a))}
                       \qquad(a\in K).
\]
By (E.2.5d), the restrictions to $F(i)$ agree or are conjugate. Both cases force $u'|_F=u|_F$, since every embedding of $F$ is real. Apply the restriction-of-place form of descent proved in E.2.3, using $h^0(\operatorname{Spec}L)$, the underlying $F$-scheme of $B_\Phi$, its labelled endomorphisms, and the polarization. The resulting finite diagram over $F$ separates its embeddings. This step does not require separation of two places of $L$ with the same restriction to $F$.

For $F=\Q$ there is one embedding. For a real quadratic field $F$, one can either descend from a ring class field as in E.2.3--E.2.4, or reduce to the preceding construction without the degree-two type degeneracy: choose a totally real cubic field $F_0$ disjoint from $F$ (a cubic field has no quadratic subfield), put $F'=FF_0$, and use the degree-six construction followed by E.2.3. This also yields a real separating test. For a CM field $K_0$ with maximal totally real subfield $F$, restriction gives a bijection from complex places of $K_0$ to real embeddings of $F$: the two extensions of each real embedding are complex conjugate. Base-changing the finite diagram over $F$ to $K_0$ therefore separates its complex places, because a comparison would restrict to the corresponding comparison over $F$. Finally every cubic field is either totally real or has negative discriminant, so the last assertion follows from this result and E.2.4.
\end{proof}

Here are explicit analytic and integral data for the cubic in E.2.4. Let $r_1<r_2<r_3$ be the roots of $x^3-4x+1$, take $\theta=r_3$, $F=\Q(\theta)$ and $K=F(i)$. Use the type with $i$-signs $(-1,-1,1)$ at those three roots, and the lattice $R=\Z[\theta,i]$. The choice $b=\theta-1$ has signs $(-1,-1,1)$, so $\xi=i(\theta-1)$ satisfies every positivity requirement in (E.2.5c). In the ordered lattice basis
\[
 1,\theta,\theta^2,\ i,i\theta,i\theta^2
\]
the polarization and the labelled endomorphisms are
\[
 [Q]=\begin{pmatrix}0&2T\\-2T&0\end{pmatrix},\qquad
 T=\begin{pmatrix}-3&8&-11\\8&-11&35\\-11&35&-52\end{pmatrix},
 \tag{E.2.5e}
\]
\[
 [\theta]=\operatorname{diag}(C,C),\qquad
 C=\begin{pmatrix}0&0&-1\\1&0&4\\0&1&0\end{pmatrix},
 \qquad
 [i]=\begin{pmatrix}0&-I_3\\I_3&0\end{pmatrix}.
\]
Indeed $T_{ab}=\operatorname{Tr}_{F/\Q}((\theta-1)\theta^{a+b})$ for $0\le a,b\le2$. The matrix is integral and alternating, with $\det T=458$ and $\det[Q]=3664^2$. The type traces are $t_\Phi(i)=-i$ and $t_\Phi(i\theta)=2i\theta$. They distinguish all three embeddings of $F$ by the proof. The program \texttt{code/verify\_numberfield\_descent.py} checks the displayed matrices, their adjoint relations, the traces, the discriminant, and the real root intervals. The algebraic number-field model is supplied by the CM descent theorem just invoked; these are explicit torus and polarization data, not a claim that projective equations for that model were numerically computed.

\begin{result}{E.2.6}{A global integral CM comparison for two quartic places}{Proved, with the stated CM scope}
Let $k=\Q(\alpha)$, where $\alpha$ is a real root of $x^4-x-1$. Its two real embeddings are distinct complex places, but their integral Betti functors are isomorphic on the following entire diagram: CM abelian varieties over $k$, finite \'etale $k$-schemes, their cohomology, products, Tate twists, tensor constructions and duals, algebraic correspondences, and compatible direct summands with their assigned intersection lattices. The comparison preserves every labelled polarization and is induced by actual algebraic isomorphisms on the CM abelian generators.

Thus this CM diagram, including all its isogenies and twists, does not separate those two places. The comparison respects ordinary Hodge structures and extension of scalars to $\R$. It does not intertwine the additional geometric real-Frobenius involutions at the two real places, as the explicit Artin witness in E.2.7 shows. No isomorphism on the diagram of all pure motives is asserted.
\end{result}
\begin{proof}
We first construct the global comparison. Inside the fixed $\Qb\subset\C$, define
\[
 \Q^{\mathrm{cm}}=\prod_{K\text{ a CM field}}K,
 \qquad
 F_{\mathrm{cm}}=\prod_{K\text{ a CM field}}K^{\mathrm{ab}},
 \tag{E.2.6a}
\]
where each product denotes the compositum of fields. Both fields are normal over $\Q$, because an automorphism of $\Qb$ permutes the CM fields and their maximal abelian extensions. We have $\Q^{\mathrm{cm}}\subset F_{\mathrm{cm}}$, and $F_{\mathrm{cm}}/\Q^{\mathrm{cm}}$ is abelian: every base extension $K^{\mathrm{ab}}\Q^{\mathrm{cm}}/\Q^{\mathrm{cm}}$ is abelian, and the compositum of such Galois extensions is abelian. Also
\[
 \Q^{\mathrm{ab}}\subset F_{\mathrm{cm}}.
 \tag{E.2.6b}
\]
Indeed, for an imaginary quadratic field $K$, the extension $K\Q^{\mathrm{ab}}/K$ is abelian and is contained in $K^{\mathrm{ab}}$.

Fix $\sigma\in\Gal(\Qb/F_{\mathrm{cm}})$ and a CM abelian variety $A/\Qb$. Choose a CM algebra $E$ acting on $A$ and its CM type $\Phi$. The reflex field $E^*$ is a CM field, also when $E$ is a CM algebra, by \cite[Proposition 1.18, p.~14]{GenMilneCM2020}. In particular, $\sigma$ fixes $E^{*\mathrm{ab}}$. The main theorem of complex multiplication gives, for a finite id\`ele $s$ with Artin image $\sigma|_{E^{*\mathrm{ab}}}$, a unique $E$-linear quasi-isogeny satisfying
\[
 \alpha_A:A\longrightarrow\sigma A,
 \qquad
 \alpha_A\bigl(N_\Phi(s)x\bigr)=\sigma x
          \quad(x\in V_f A).
\]
This is exactly \cite[Theorem 9.10, p.~78]{GenMilneCM2020}; $V_f A$ is the product of the rational Tate modules with its usual restricted product structure. Here we may take $s=1$, so $N_\Phi(s)=1$ and
\[
 V_f(\alpha_A)=\sigma:V_f A\longrightarrow V_f(\sigma A).
 \tag{E.2.6c}
\]

We now prove that this is an actual algebraic isomorphism. Write $\alpha_A=f/m$, where $m>0$ is an integer and $f:A\to\sigma A$ is an actual homomorphism. Galois transport is an isomorphism on every integral Tate module, so (E.2.6c) implies
\[
 f(T_\ell A)\subset mT_\ell(\sigma A)
                         \quad\hbox{for every prime }\ell.
\]
For each $\ell\mid m$, reduction modulo $\ell^{v_\ell(m)}$ says that $f$ kills $A[\ell^{v_\ell(m)}]$. Thus $f$ kills $A[m]$. The quotient property of $[m]:A\to A$ gives an actual homomorphism $\alpha_0$ with $f=\alpha_0\circ[m]$, and equality in $\operatorname{Hom}(A,\sigma A)\otimes\Q$ gives $\alpha_0=\alpha_A$. The same argument applies to $\alpha_A^{-1}$, since its Tate action is the inverse integral Galois transport. Hence $\alpha_A$ and its inverse are actual morphisms. In particular,
\[
 (\alpha_A)_*:H_1(A(\C),\Z)\xrightarrow{\sim}
                         H_1((\sigma A)(\C),\Z)
 \tag{E.2.6d}
\]
comes from one global algebraic map. This verifies the integral passage, rather than assuming that separate local comparisons glue.

A homomorphism of abelian varieties in characteristic zero is determined by its action on a rational Tate module: after the complex comparison, a homomorphism acting trivially has zero map on its homology lattice and hence zero map of complex tori. Consequently $\alpha_A$ is independent of the chosen CM algebra. For every homomorphism $f:A\to B$, the two maps
\[
 (\sigma f)\circ\alpha_A,\qquad \alpha_B\circ f
 \tag{E.2.6e}
\]
have the same Tate action and are equal. The same argument gives $\alpha_{A\times B}=\alpha_A\times\alpha_B$ and compatibility with dual abelian varieties. By (E.2.6b), $\sigma$ fixes every root of unity, so its cyclotomic character is $1$. The Weil pairing therefore shows that the one Tate object is fixed and every labelled polarization is preserved. Equivalently, the precise multiplier in \cite[Remark 9.11(c), p.~78]{GenMilneCM2020} is
\[
 \chi_{\mathrm{cyc}}(\sigma)/N_{E^*/\Q}(s)=1.
\]
The induced maps preserve the Hodge structures as well, because $\alpha_A$ is an algebraic isomorphism over $\Qb\subset\C$.

Here is the extension to arbitrary correspondences in the stated diagram. Integral Betti cohomology of an abelian variety is the exterior algebra of $H^1(A,\Z)$, and the maps (E.2.6d) induce the corresponding exterior-power and K\"unneth maps. The comparison on cohomology is the inverse pullback $(\alpha_A^{-1})^*$. After tensoring with $\Z_\ell$, it equals Galois transport on \'etale cohomology, by the finite-coefficient comparison and its inverse limit \cite[Expos\'e XVI, Th\'eor\`eme 4.1]{GenSGA4}. An algebraic correspondence $Z$ has a Galois-equivariant \'etale cycle class: transport sends its class to the class of $\sigma Z$. Thus the maps commute with the induced correspondence operators after tensoring with $\Q_\ell$. Faithfulness of $\Q\hookrightarrow\Q_\ell$ gives the same equality on rational Betti cohomology. On the assigned free integral groups the maps are integral, so the equality holds there too. This argument treats all algebraic correspondences, without assuming that they are generated by endomorphisms.

For a finite \'etale scheme $S/\Qb$, send the basis element of $H^0(S,\Z)$ indexed by a geometric point $x$ to the basis element indexed by $\sigma x$ in $H^0(\sigma S,\Z)$. These permutations commute with maps of finite schemes, cup products and the diagonal classes. On products with CM abelian varieties take these permutations of components together with the maps $\alpha_A$ on the corresponding components. This is again the \'etale Galois action after completion, so the same cycle-class argument proves naturality for correspondences involving Artin factors. Tensor products and duals use the tensor products and dual inverses of the maps already constructed; the Tate map is the identity. Finally, a compatible idempotent $p$ satisfies $T p=(\sigma p)T$, so $T$ restricts to its image and carries the assigned intersection lattice isomorphically onto the conjugate one. This completes a simultaneous integral tensor comparison on the entire stated diagram.

It remains to realize two distinct places with such a $\sigma$. Let $L$ be the splitting field of $x^4-x-1$, put
\[
 K=\Q(\sqrt{-283}),\qquad H=L^{V_4}.
\]
The polynomial is irreducible modulo $2$ and has discriminant $-283$. Its cubic resolvent $x^3+4x-1$ is irreducible by the rational-root test. Hence $\Gal(L/\Q)=S_4$: the resolvent makes its action on pair partitions transitive, excluding the transitive $2$-subgroups, and the nonsquare discriminant excludes $A_4$. These calculations are checked in \texttt{code/verify\_numberfield\_descent.py}. The polynomial has two real roots and one conjugate pair, so actual complex conjugation $c$ is a transposition. Moreover,
\[
 H\text{ is the Hilbert class field of }K,
 \qquad [H:K]=3;
 \tag{E.2.6f}
\]
this exact description is given in \cite[Notes 5.4, p.~438]{GenSerreJordan2003}.

Complex conjugation is central on $\Q^{\mathrm{cm}}$. To see the precise relation, for a CM field $E$ with intrinsic involution $c_E$ and any embedding $r:E\hookrightarrow\C$, the total reality of $E^+$ gives
\[
 c\circ r=r\circ c_E.
\]
These identities imply commutation with every automorphism on the compositum of all conjugates. Therefore $L\cap\Q^{\mathrm{cm}}$ is a Galois subfield of $L$ whose Galois group is a quotient of $S_4$ in which the image of $c$ is central. The normal subgroup generated by all commutators $[g,c]$ is $A_4$: each is even, and conjugating $c$ to a transposition sharing exactly one letter produces a $3$-cycle; its conjugates generate $A_4$. The quotient thus has order at most two. Since the imaginary quadratic field $K$ is contained in both fields,
\[
 L\cap\Q^{\mathrm{cm}}=K,
 \qquad
 \Gal(L\Q^{\mathrm{cm}}/\Q^{\mathrm{cm}})=A_4.
 \tag{E.2.6g}
\]

Because $F_{\mathrm{cm}}/\Q^{\mathrm{cm}}$ is abelian, the intersection of $L\Q^{\mathrm{cm}}$ with $F_{\mathrm{cm}}$ corresponds to an abelian quotient of $A_4$. Its commutator subgroup is $V_4$, so that intersection is contained in $H\Q^{\mathrm{cm}}$. Intersecting with $L$ and using (E.2.6g) gives $L\cap F_{\mathrm{cm}}\subset H$. Conversely (E.2.6f) gives $H\subset K^{\mathrm{ab}}\subset F_{\mathrm{cm}}$. We have proved
\[
 L\cap F_{\mathrm{cm}}=H.
 \tag{E.2.6h}
\]

Consequently the restriction map
\[
 \Gal(\Qb/F_{\mathrm{cm}})\longrightarrow\Gal(L/H)=V_4
\]
is surjective. Explicitly, an automorphism of $L$ fixing $H$ extends to $LF_{\mathrm{cm}}$ by the identity on $F_{\mathrm{cm}}$, because of (E.2.6h), and then extends to $\Qb$. Label the roots so that $c=(12)$ and roots $3,4$ are real. Choose $\sigma$ restricting to $(12)(34)$, and let $\alpha$ be root $3$. The embeddings
\[
 u:k\hookrightarrow\R\subset\C,\qquad w=\sigma\circ u
\]
send $\alpha$ to the two different real roots. They are distinct and are not related by complex conjugation. For every CM abelian variety over $k$, the construction above applied after $u$ compares $uA$ with $\sigma(uA)=wA$. Equations (E.2.6d)--(E.2.6e), the Artin permutations, and the correspondence argument give the required single natural tensor isomorphism.
\end{proof}

The remaining full-category question has an exact form. Let $G_{\Z}$ denote the tensor automorphisms of the entire concrete integral Betti diagram over $\Q$, with the integral conventions of E.2.3. Its Artin restriction is a homomorphism to $\Gal(\Qb/\Q)$. Universal separation of number-field complex places is the equality of its image with $\{1,c\}$ proved equivalent in E.2.3. The automorphisms constructed in E.2.6 belong to the restricted CM-abelian/Artin diagram. They have not been extended to $G_{\Z}$. In particular, this result gives neither a full-category comparison between the quartic places nor a finite full-category diagram separating them.

\begin{result}{E.2.7}{Labelled real Frobenius separates every real embedding}{Proved}
Let $k$ be a number field. Distinct real embeddings of $k$ are separated by a finite Artin diagram equipped with the geometric real-Frobenius involutions. If $k$ has $r_1\ge2$ real embeddings, $r_1-1$ explicitly chosen quadratic Artin objects suffice. In the quartic example of E.2.6, the single object $\operatorname{Spec}k(\sqrt{\alpha})$ separates the two real embeddings by Frobenius traces $0$ and $2$.
\end{result}
\begin{proof}
Choose a primitive element $\theta$ of $k/\Q$ and order its values at the real embeddings as
\[
 u_1(\theta)<\cdots<u_{r_1}(\theta).
\]
Choose rational numbers $q_j$ with
\[
 u_j(\theta)<q_j<u_{j+1}(\theta)
                         \qquad(1\le j<r_1),
\]
and put
\[
 S_j=\operatorname{Spec}\bigl(k[t]/(t^2-(\theta-q_j))\bigr).
\]
Each is finite \'etale of degree two. At a real embedding $u$, its two complex geometric points are both real when $u(\theta-q_j)>0$, and are exchanged by complex conjugation when $u(\theta-q_j)<0$. Hence geometric real Frobenius on the free rank-two group $H^0(S_{j,u}(\C),\Z)$ has trace
\[
 \operatorname{Tr}(f_{\infty,u}|H^0(S_j))=
 \begin{cases}
 2,&u(\theta-q_j)>0,\\
 0,&u(\theta-q_j)<0.
 \end{cases}
 \tag{E.2.7a}
\]
For $u_i\ne u_{i'}$, one of the chosen thresholds lies between their values of $\theta$, giving different traces. An isomorphism intertwining the involutions would preserve trace, so no such comparison exists on that finite diagram, even over $\Q$ or $\R$.

For $x^4-x-1$, the two certified real-root intervals in the computation are negative and positive respectively. Taking $\theta=\alpha$ and $q=0$ gives $S=\operatorname{Spec}k(\sqrt{\alpha})$. Its Frobenius matrix is a transposition at the negative root and the identity at the positive root, with the displayed traces. The plain comparison from E.2.6 still permutes the two basis elements of $H^0$, whose ordinary Hodge structures both have type $(0,0)$, but it cannot conjugate a transposition to the identity. Thus the two interpretations of \emph{real forms} have precise different tests: scalar extension of the Betti groups to $\R$ is preserved in E.2.6, whereas adjoining the actual real-place involutions invokes (E.2.7a). This conclusion leaves the plain integral question unchanged.
\end{proof}

\begin{result}{E.2.8}{A finite comparison skeleton in one fixed-field isogeny class}{Proved}
Let $A$ be an abelian variety over a number field $k$, and let
$\mathcal I_k(A)$ be the category of abelian varieties over $k$ that are
$k$-isogenous to $A$, with $k$-homomorphisms. For embeddings
$v,w:k\hookrightarrow\C$, there is a finite labelled diagram
$\mathcal D\subset\mathcal I_k(A)$ such that restriction gives a bijection
between natural integral $H_1$ comparisons on $\mathcal I_k(A)$ and
comparisons on $\mathcal D$.
\end{result}
\begin{proof}
Finiteness I gives representatives $B_1,\ldots,B_s$ of the finitely many
$k$-isomorphism classes in $\mathcal I_k(A)$. The precise fixed-field
assertion is stated in \cite[Introduction, p.~1]{GenMW1991}; that primary
paper gives its isogeny-estimate proof in Section~7, pp.~10--11.
For every ordered pair $(i,j)$, the group
$\operatorname{Hom}_k(B_i,B_j)$ is free of finite rank. Indeed, one
complex embedding injects it into
\[
 \operatorname{Hom}_{\Z}
       \bigl(H_1(B_{i,v},\Z),H_1(B_{j,v},\Z)\bigr).
\]
Injectivity follows because a homomorphism of complex tori is determined
by its linear map on the period lattice, and a subgroup of a finite free
abelian group is finite free. Choose a $\Z$-basis for each such Hom
group. The finite diagram $\mathcal D$ consists of the $B_i$ and these
labelled basis homomorphisms.

Write $F_u(B)=H_1(B_u,\Z)$ for $u=v,w$. A comparison on $\mathcal D$
consists of integral isomorphisms $T_i:F_v(B_i)\to F_w(B_i)$ commuting
with those basis maps. For each $B$ choose a $k$-isomorphism
$b_B:B\to B_i$, and define
\[
 T_B=F_w(b_B)^{-1}\,T_i\,F_v(b_B):
               F_v(B)\longrightarrow F_w(B).
 \tag{E.2.8a}
\]
Two choices of $b_B$ differ by a $k$-automorphism of $B_i$.
That automorphism is a $\Z$-linear combination of the labelled basis
of $\End_k(B_i)$, so naturality for the basis makes (E.2.8a)
independent of the choice. For a $k$-homomorphism $f:B\to C$,
the conjugate $b_Cfb_B^{-1}$ is a $\Z$-linear combination of the
basis of $\operatorname{Hom}_k(B_i,B_j)$. It follows that
\[
 T_C F_v(f)=F_w(f)T_B.
\]
Thus (E.2.8a) extends the comparison to the whole category. Conversely,
naturality for $b_B$ forces (E.2.8a), proving uniqueness and the
restriction bijection.

For the contravariant cohomology convention, a comparison
$\eta_B:H^1(B_v,\Z)\to H^1(B_w,\Z)$ gives the homology map
$T_B=(\eta_B^{-1})^\vee$. The universal coefficient pairing is perfect
because these groups are free. This is the precise passage between the
two conventions used in this result.

The finite skeleton concerns the degree-one isogeny category just
specified. The full integral correspondence diagram also contains
varieties obtained from arbitrarily large finite extensions, and its
additional constraints are calculated next.
\end{proof}

\begin{result}{E.2.9}{Geometric cyclic quotients impose projective Tate compatibility over a fixed number field}{Proved necessary condition}
Let $A/k$ be an abelian variety of dimension $g>0$, and let
$\eta$ be a tensor comparison between the full graded integral Betti
correspondence diagrams at embeddings $v,w:k\hookrightarrow\C$.
Let $\sigma\in\Gal(\Qb/\Q)$ be the Artin permutation selected by the
exact transporter E.2.3, so $\sigma\circ v=w$, and put
\[
 \Lambda_u=H_1(A_u,\Z),\qquad
 T_A=(\eta_A^{-1})^\vee:\Lambda_v\xrightarrow{\sim}\Lambda_w.
\]
For every $n\ge2$, define
\[
 \kappa_{u,n}:\Lambda_u/n\Lambda_u\xrightarrow{\sim}A_u[n](\C),
 \qquad x\bmod n\Lambda_u\longmapsto x/n\bmod\Lambda_u
\]
through exponential uniformization. Torsion points are algebraic, so
Galois transport gives the isomorphism
\[
 \rho_{A,n}(\sigma)
   =\kappa_{w,n}^{-1}\circ\sigma\circ\kappa_{v,n}:
       \Lambda_v/n\Lambda_v\longrightarrow\Lambda_w/n\Lambda_w.
\]
These maps are compatible under reduction, and their inverse limit is
$\rho_A(\sigma):\Lambda_v\otimes\widehat\Z
 \to\Lambda_w\otimes\widehat\Z$. There is a unique
$\lambda_A\in\widehat\Z^\times$ such that
\[
 T_A\otimes\widehat\Z=\rho_A(\sigma)\,\lambda_A.
 \tag{E.2.9a}
\]
In particular, the projective Tate transport of the selected
$\sigma$ must be represented by one integral unimodular Betti matrix.
\end{result}
\begin{proof}
The action on all finite \emph{\'etale} schemes supplies one compatible
Artin permutation, as proved in E.2.3. Its compatibility with every
finite extension, rather than an independent choice for each quotient,
is essential here.

Fix $n\ge2$ and a cyclic subgroup $H\subset A_v[n]$. Choose a finite
extension $L/k$ and an embedding $u:L\hookrightarrow\Qb$ extending
$v$ over which all the points of $H$ are defined. The finite subgroup
then gives an abelian quotient and an isogeny over $L$,
\[
 q:A_L\longrightarrow B=A_L/H.
 \tag{E.2.9b}
\]
Existence over the field of definition of the subgroup is the quotient
and descent statement in \cite[Chapter IV, Lemma~2.1, p.~135]{GenMilneAV2008}.
It can also be checked by the universal property of the finite quotient:
the invariant subgroup makes every Galois conjugate quotient canonically
the same quotient, and the resulting descent maps satisfy the cocycle
identity.

Regard both schemes in (E.2.9b) as schemes over $k$ by composing their
structural maps with $\operatorname{Spec}L\to\operatorname{Spec}k$.
Include the Artin algebra $h^0(\operatorname{Spec}L)$, its
multiplication and the module actions. The primitive idempotent
indexed by $u$ selects exactly the groups $\Lambda_v$ and
$H_1(B_u,\Z)$. Its image under the comparison is the idempotent
indexed by $\sigma u$. On $A_L=A\times_k\operatorname{Spec}L$,
tensor compatibility makes the map on these selected blocks exactly
$T_A$. The graph of $q$ is a correspondence on the diagonal open and
closed component of the product over $k$, as checked in E.2.3.
Naturality therefore gives the integral square
\[
 T_B q_*=(q_{\sigma H})_*T_A,
 \qquad
 T_B:H_1(B_u,\Z)\xrightarrow{\sim}H_1(B_{\sigma u},\Z).
 \tag{E.2.9c}
\]

Under $\kappa_{v,n}$ the subgroup $H$ is exactly
\[
 \ker\!\left(
    q_*\bmod n:\Lambda_v/n\Lambda_v
         \longrightarrow H_1(B_u,\Z)/nH_1(B_u,\Z)\right).
 \tag{E.2.9d}
\]
Indeed $q(x/n)$ is zero on the quotient torus precisely when
$q_*(x)$ belongs to $nH_1(B_u,\Z)$. This proves (E.2.9d) for every
cyclic subgroup of $A_v[n]$, including subgroups of order properly
dividing $n$. Since $T_B$ and its reduction modulo $n$ are invertible,
(E.2.9c) carries this kernel onto the kernel indexed by $\sigma H$.
Consequently $T_A\bmod n$ and $\rho_{A,n}(\sigma)$ have the same
action on every cyclic subgroup of $\Lambda_v/n\Lambda_v$.

We use the following elementary computation over the ring $\Z/n$,
with no assumption that $n$ is prime. An invertible matrix on
$(\Z/n)^r$, $r\ge2$, preserving every cyclic subgroup must be scalar.
The coordinate lines force all off-diagonal entries to vanish.
The diagonal entries are units because the matrix is invertible.
Preserving the line generated by $e_i+e_j$ forces the $i$-th and
$j$-th diagonal entries to agree. Thus the matrix is
$\lambda_n I$ for a unique $\lambda_n\in(\Z/n)^\times$.
This argument includes $n=2$ and every composite $n$.

Apply it to
$\rho_{A,n}(\sigma)^{-1}(T_A\bmod n)$, using $r=2g\ge2$.
For $n\mid m$, reduction of the scalar at level $m$ is the scalar
at level $n$ by uniqueness. The $\lambda_n$ therefore define one
$\lambda_A\in\widehat\Z^\times$, and passage to the inverse limit
gives (E.2.9a), including uniqueness. The compatibility of the
$\kappa_{u,n}$ uses multiplication by $m/n$ on $m$-torsion and
reduction $\Lambda_u/m\Lambda_u\to\Lambda_u/n\Lambda_u$.

These are actual objects and maps in the full diagram over the
\emph{fixed} field $k$. They do not replace $k$ by $\Qb$ in the
claim. The schemes $B_{[k]}$ have $[L:k]$ geometric components and
their degree-one Betti rank is $2g[L:k]$. The rational isogeny (E.2.9b)
identifies this realization with
$H_1(A)\otimes H_0(\operatorname{Spec}L)$.
As the construction ranges through finite extensions, these objects
are outside the fixed rank-$2g$ isogeny category of E.2.8.
This explicit inclusion of diagrams explains the two finiteness
scopes. It does not assert that any selected Galois element fails
(E.2.9a); proving such a failure or constructing all the compatible
integral matrices remains part of (E.2.3d).
\end{proof}

\begin{result}{E.2.10}{Finite congruence solvability does not supply a simultaneous rational lift}{Proved algebraic-group counterexample}
There are nested compact open cosets in
$\mathrm{SL}_2(\widehat\Z)$, each containing an element of
$\mathrm{SL}_2(\Z)$, whose common intersection contains no element
of $\mathrm{SL}_2(\Q)$.
\end{result}
\begin{proof}
In $\widehat\Z=\prod_\ell\Z_\ell$, choose the explicit element
$a_2=0$ and $a_\ell=1$ for every odd prime $\ell$, and set
\[
 u(a)=\begin{pmatrix}1&a\\0&1\end{pmatrix}.
 \]
This element is not rational: its $2$-adic component would force a
rational representative to be zero, contradicting each odd component.
For $N\ge2$, let $K_N$ be the subgroup of
$\mathrm{SL}_2(\widehat\Z)$ preserving every cyclic subgroup of
$(\Z/N!\Z)^2$. It is the inverse image of the scalar subgroup
under reduction modulo $N!$, by the all-lines computation in
E.2.9; hence it is compact open, and $K_{N+1}\subset K_N$.
Choose an integer $a_N$ with $a_N\equiv a\pmod{N!}$.
Then $u(a_N)\in u(a)K_N\cap\mathrm{SL}_2(\Z)$.

On the other hand, an element $B$ in every coset satisfies
$u(a)^{-1}B\equiv\lambda_N I\pmod{N!}$ for compatible units
$\lambda_N$. Therefore $B=\lambda u(a)$ for
$\lambda\in\widehat\Z^\times$. If $B$ had rational entries, its
nonzero diagonal entry $\lambda=B_{11}$ and its upper-right
entry $\lambda a=B_{12}$ would be rational. Their quotient
would make $a$ rational, a contradiction. Thus
\[
 \mathrm{SL}_2(\Q)\ \cap\
       \bigcap_{N\ge2}u(a)K_N=\varnothing.
 \tag{E.2.10a}
\]
This proves the failure of the indicated finite-to-infinite inference
even for the simply connected group $\mathrm{SL}_2$. The cosets in
this proof are not asserted to be the motivic transporter of the
quartic example. They specify the additional simultaneous-integrality
step that a strong-approximation approach to (E.2.3d) must actually
establish.
\end{proof}

\begin{result}{E.2.11}{The prescribed CM comparison has real and rational abelian lifts}{Proved auxiliary result; full integral gate retained}
Let $F_{\rm CM}\subset\Qb$ be the compositum of $K^{\rm ab}$ for all
finite CM fields $K$, and take $\sigma\in\Gal(\Qb/F_{\rm CM})$.
Let $\mathrm{AV}(\Q)$ be the category of pure abelian motives with
absolute Hodge correspondences, including Artin motives and Tate
objects. Write
\[
 \pi:G_{\mathrm{AV},B}\longrightarrow G_{\mathrm{CM},B},
 \qquad D=\ker\pi
\]
for restriction of its rational Betti group to the CM/Artin subcategory.
The prescribed CM comparison $g_\sigma$ of E.2.6 defines a point of
$G_{\mathrm{CM},B}(\Q)$. Its fibre
\[
 X_\sigma=\pi^{-1}(g_\sigma)
 \tag{E.2.11a}
\]
has both a real point and a rational point.

The real point can be chosen to preserve Hodge structures and labelled
polarizations. The rational point is a tensor comparison in the
absolute-Hodge abelian category. This assertion supplies neither
integrality of that rational point nor its equality with the chosen
real Hodge point. The full pure integral question (E.2.3d) remains open
at this step.
\end{result}
\begin{proof}
\emph{The prescribed CM point and its morphisms.}
Theorem~3.15 of
\cite[p.~41 in the corrected author version]{GenMilneShimura1994}
gives, for a CM motive $M$ and a representative $\widetilde r(\sigma)$
of its reciprocity class, a unique isomorphism
$\alpha_M:M\to{}^\sigma M$ carrying every labelled Hodge tensor $s$
to $\sigma(s)$, with
\[
 \sigma x=(\alpha_M)_f\bigl(\widetilde r(\sigma)x\bigr)
                  \qquad(x\in\omega_f(M)).
 \tag{E.2.11b}
\]
Every relevant reflex field $E$ is CM or totally real. In the totally
real case $E(i)$ is CM and $E^{\rm ab}E(i)/E(i)$ is abelian, so
$E^{\rm ab}\subset E(i)^{\rm ab}\subset F_{\rm CM}$.
The containment is immediate when $E$ is CM. Thus the reciprocity
class is trivial, and the representative $1$ in (E.2.11b) gives the
precise maps of E.2.6. This verifies their naturality for all
absolute Hodge correspondences as well as algebraic ones.
Their Artin components are $e_x\mapsto e_{\sigma x}$ and their
component on the one Tate line is the identity.

\emph{A finite comparison with the prescribed CM restriction.}
Consider finitely many abelian varieties over $\Qb$, finitely many
labelled absolute Hodge tensor maps, chosen polarizations, and a
finite set of CM objects on which the comparison is prescribed.
Take one product $A$ of the abelian generators, including a product
$B$ of CM abelian generators for those CM objects. Include the factor
projectors, the polarization tensors, and the one Tate object.
Let $V=H_1(A(\C),\Q)$, $\Lambda=H_1(A(\C),\Z)$, and
$G=\operatorname{MT}(V,h)$. Enlarge the indexed Hodge tensors to
cut out $G$, retaining all the original labels.

The $G(\R)$-conjugacy class $X$ of $h$ gives its Hodge-type Shimura
datum. The exact moduli statements used here are
\cite[Example~3.1, p.~36; Lemma~3.29 and Theorem~3.31,
pp.~45--46]{GenMilneShimura1994}: the canonical model parametrizes
the indexed absolute Hodge tensors and finite-adic level structure,
and conjugating over the reflex field conjugates those labelled
data. The integral local system is explicit in the same source:
conditions (2.7.1)--(2.7.3), Proposition~2.8, Example~2.11 and
Theorem~2.12, pp.~23--25, construct the constant lattice on $X$
and its quotient on $\Gamma(N)\backslash X$. Take $N$ sufficiently
divisible, with $3\mid N$. Example~3.7, p.~38, identifies this with
the $K(N)$-level description. The level orbit determines
$\eta(\Lambda\otimes\widehat\Z)$; intersection with the rational
realization recovers the actual integral lattice. Thus no unspecified
lattice on an isogeny-class representative is substituted.

The reflex field $E(G,X)$ is a subfield of a CM field by the proof
of \cite[Theorem~3.33, p.~46]{GenMilneShimura1994}. Such a subfield
is CM or totally real: in a CM Galois closure complex conjugation is
central, hence preserves each intermediate field, acting either
trivially or as its CM involution. The preceding calculation gives
$E(G,X)^{\rm ab}\subset F_{\rm CM}$.
Deligne's component reciprocity theorem
\cite[Section~2.6.1 and Theorem~2.6.3, p.~46 of the translation]{GenDeligneShimura1979}
therefore puts the point defined by $A$ and its $\sigma$-conjugate
in the same connected component at this level.
Parallel transport along a path between them supplies an integral
map
\[
 T_0:H_1(A(\C),\Z)\xrightarrow{\sim}
                     H_1({}^\sigma A(\C),\Z)
 \tag{E.2.11c}
\]
preserving every labelled tensor and compatible with $\sigma$ on
level $N$. Its Tate action is the identity: on the integral rank-one
Tate local system it is a unit congruent to $1$ modulo $3$.

The CM restriction of (E.2.11c) is exactly $\alpha_B$.
Indeed the representation of $G$ on $H_1(B)$ has torus image;
the corresponding Hodge homomorphism is constant under conjugation
along the connected domain $X^+$. Hence $T_{0,B}$ is an integral
Hodge isometry. Both $T_{0,B}$ and $\alpha_B$ induce $\sigma$ on
$B[N]$. Their quotient is a polarized integral Hodge automorphism
that is the identity modulo $3$, and such an automorphism is the
identity. To check the last assertion, it preserves the positive
definite form obtained from the polarization and Weil operator,
so its integral matrix lies in a finite group. But
$1+3M_r(\Z_3)$ is torsion-free. A matrix $U=I+3^sV$, $s\ge1$,
with $V\not\equiv0\pmod3$ cannot have prime order $p\ne3$ by
binomial expansion modulo $3^{s+1}$; for $p=3$,
\[
 U^3=I+3^{s+1}\bigl(V+3^sV^2+3^{2s-1}V^3\bigr)
\]
has the parenthesis nonzero modulo $3$. Taking a prime-order power
excludes every nontrivial finite-order element.

The map $T_0$ need not preserve the Hodge structure of $A$.
After trivializing along the path, $h$ and
$T_0^{-1}h_{{}^\sigma A}T_0$ belong to the same $X^+$.
The map
$G^{\rm der}(\R)^+\to G^{\rm ad}(\R)^+$
is surjective: its differential is an isomorphism, so its image is
open in the connected target. Choose $d$ in the derived group
carrying the first Hodge homomorphism to the second. Then
\[
 T=T_0d
 \tag{E.2.11d}
\]
is a real Hodge isometry preserving all the indexed tensors.
The derived group acts trivially on CM representations and on
Tate, so $T_B=\alpha_B$. It preserves polarizations because their
similitude characters factor through the abelian quotient.
This proves finite real solvability with the prescribed CM map.

\emph{Coherence and Artin blocks.}
For every abelian variety $A/\Qb$, choose a polarization and let
$C_A$ be the set of real Hodge isometries from $H_1(A)$ to
$H_1({}^\sigma A)$ for that polarization and its conjugate.
Each $C_A$ is nonempty by the preceding construction and compact:
after choosing one isometry it is a torsor under a closed subgroup
of the orthogonal group of the associated positive definite form.
In $\prod_A C_A$, impose all naturality, tensor, unit, duality,
projector and labelled-polarization equations, together with
$T_B=\alpha_B$ for every CM generator. These are closed conditions,
each involving finitely many coordinates. Every finite collection
has a solution by (E.2.11d), applied to its finite product of
generators. Compactness and the finite intersection property give
a simultaneous solution.

Extend it by tensor operations and exterior powers; restrict to
the image of each projector. Naturality for the transition
isomorphisms between presentations makes the resulting comparison
well-defined on every abelian motive over $\Qb$.
For descent to $\Q$, send $e_u$ to $e_{\sigma u}$ on each finite
Artin embedding set. On the underlying $\Q$-scheme of an object
over a finite field $L$, send its $u$-summand to its $\sigma u$-summand
by the constructed map. The Artin module action selects these
blocks, and $\sum_u e_u\otimes e_u$ selects the relative tensor
product, exactly as in E.2.3. Each correspondence equation holds
on every geometric block and hence on their direct sum. Restriction
of embedding sets gives compatibility under changing $L$.
On $\Q$-defined objects, the canonical identification with their
$\sigma$-conjugates yields
$g_{\R}\in G_{\mathrm{AV},B}(\R)$ with
$\pi(g_{\R})=g_\sigma$. This proves $X_\sigma(\R)\ne\varnothing$.
No geometric real-Frobenius intertwining has been imposed;
E.2.7 gives its separate obstruction for the quartic places.

\emph{The rational point.}
For every finite prime $\ell$, Galois transport by $\sigma$ on
the etale realization, carried through Betti--etale comparison,
gives $g_\ell\in G_{\mathrm{AV},B}(\Q_\ell)$.
Naturality includes absolute Hodge correspondences through their
etale realizations. Equation (E.2.11b) with representative $1$
proves $\pi(g_\ell)=g_\sigma\otimes\Q_\ell$.
Thus the actual fibre (E.2.11a) has a point over every completion.

The proof of \cite[Proposition~3.12 and Theorem~3.13,
p.~26 of the author version]{GenMilnePeriods2004}
identifies $D$ as a countable product of semisimple algebraic
factors whose geometric simple factors have covering index $1$
or $2$. Proposition~3.5, pp.~23--24, proves the all-local Hasse
injection for those factors, and Corollary~1.14, p.~9, identifies
$H^1$ of their countable product with the product of the $H^1$
sets. Consequently
\[
 H^1(\Q,D)\longrightarrow\prod_v H^1(\Q_v,D)
\]
is injective. The class of $X_\sigma$ is zero in every target
factor by the points just constructed. It is therefore zero
over $\Q$, giving $X_\sigma(\Q)\ne\varnothing$.
This uses the proved product structure; it does not infer a
compatible lift from unrelated points of an inverse system.

The rational point obtained at this last step is not required to
equal $g_{\R}$ after scalar extension. In particular the construction
has proved a rational abelian comparison with the prescribed CM
restriction, while the simultaneous lattice conditions of E.2.9
and the full pure-motive requirement in (E.2.3d) still require their
own proofs.
\end{proof}

\begin{result}{E.2.12}{Finite equations give a comparison in an admissibly generated integral category}{Proved with explicit lattice hypotheses}
Let $R$ be noetherian, let $\mathcal C$ be a small rigid symmetric
$R$-linear exact tensor category, and let
\[
 \omega_1,\omega_2:\mathcal C\longrightarrow\operatorname{Proj}_f(R)
\]
be faithful exact tensor functors. Assume every object is obtained,
up to isomorphism, from one object $M$ by finite sums, tensor products,
duals, and admissible subobjects and quotients, and that
$\omega_1(M),\omega_2(M)$ are finite free of the same rank.
Then $\Isom^\otimes(\omega_1,\omega_2)$ is represented by a closed
subscheme of $\Isom_R(\omega_1(M),\omega_2(M))$. A finite diagram of
original objects, maps and admissible sequences determines each
comparison uniquely.

For $R=\Z$, consider such a category of abelian motives over the
quartic field of E.2.6 with its two integral Betti realizations.
Require that the generator's lattices are given by a specified
admissible expression from the actual Betti lattices of finitely
many abelian and Artin generators. Then these two realizations
are integrally tensor isomorphic on $\mathcal C$. In this
specified integral model, a single admissible tensor generator
does not distinguish those quartic places.
\end{result}
\begin{proof}
After bases, the ambient isomorphism scheme has coordinate ring
$B=R[u_{ij},\det(u)^{-1}]$, with universal matrix $u$ and inverse
$v$. Finite sums, tensor products and duals give the corresponding
maps $u_P,v_P$. For later objects whose realizations are projective
but not free, realize each as an idempotent summand of a finite
free module and write the equations there. Equivalently, evaluate
each module-valued equation against the dual of its finite
projective Hom module to obtain its coefficient ideal. The
equations are therefore global over $R$, without bases on every
projective module.

For an admissible monomorphism $i:N\to P$, with $P$ already
constructed, put $i_a=\omega_a(i)$ and choose retractions $r_a$
with $r_ai_a=1$. They exist because the cokernel is projective.
Impose
\[
 (1-i_2r_2)u_Pi_1=0,\qquad
 (1-i_1r_1)v_Pi_2=0,
 \tag{E.2.12a}
\]
and define
\[
 u_N=r_2u_Pi_1,\qquad v_N=r_1v_Pi_2.
 \tag{E.2.12b}
\]
The equations preserve the first image under $u_P$ and the
second under $v_P$. Split identities persist after extension to
every $R$-algebra $S$. Substitution therefore gives
$u_Nv_N=1=v_Nu_N$ over every $S$, including nonflat base changes.
For an admissible quotient $q:P\to N$, preserve its kernel in
both directions and choose sections $s_a$ of $q_a=\omega_a(q)$.
The induced inverse quotient maps are
\[
 u_N=q_2u_Ps_1,\qquad v_N=q_1v_Ps_2.
 \tag{E.2.12c}
\]
They are independent of the sections because the differences
lie in the preserved kernels; their composites are identities
because $s_aq_a-1$ has image in those kernels.

Choose one finite admissible expression for every object, apply
these formulas, and impose every original naturality equation
\[
 u_Y\omega_1(f)=\omega_2(f)u_X\qquad(f:X\to Y),
 \tag{E.2.12d}
\]
together with the tensor, unit, symmetry and duality equations.
The original isomorphisms between presentations are included, so
the resulting maps are independent of the chosen expressions.
Every solution is a tensor natural isomorphism, and conversely
every such isomorphism satisfies the equations. Its component on
$M$ determines all the other components. This proves the stated
closed-subscheme representation.

Let $I\subset B$ be the ideal of all coefficient equations.
Noetherianness makes $I$ generated by a finite subset of the
\emph{original} equations: choose finitely many ideal generators,
express each using finitely many original coefficients, and take
their union. Include the finite admissible presentations of the
involved objects. This gives the finite original diagram.
Prescribed components on a fixed Artin or CM object add ordinary
matrix equations and are covered by the same argument.

For the geometric application, use the element
$\sigma\in\Gal(\Qb/F_{\rm CM})$ of E.2.6 exchanging the specified
real roots. The finite-diagram construction (E.2.11c), before its
real Hodge correction, gives an integral map $T_0$ on the finite
defining diagram. It preserves all the requested rational absolute
Hodge tensors and has the $\sigma$-selected Artin permutation.
The morphisms have integral Betti matrices because they are
morphisms of $\mathcal C$. The admissible expressions recover
saturated image lattices and free quotient lattices: an integral
ambient isomorphism preserving a rational image subspace preserves
its intersection with the ambient lattice, and then induces the
isomorphism on the free quotient. Thus $T_0$ satisfies the finite
original equations. Its component on $M$ satisfies all of $I$,
and (E.2.12a)--(E.2.12d) extend it uniquely to the asserted
global integral comparison.

The lattice hypothesis is part of the assertion. The injection
on $H_1$ induced by a nontrivial isogeny $q:E\to E/H$ has finite
torsion cokernel, so that injection is not an admissible subobject
construction in the exact category of finite free $\Z$-modules.
This does not exclude $q$ as an ordinary morphism when both
objects have independently supplied admissible expressions:
then (E.2.12d) includes its naturality. Nor does it assert that
every isogenous target lacks such an expression. It states the
additional presentation required before applying the result to
that lattice. The entire geometric diagram in E.2.9 has not been
identified with this admissibly generated category, and the full
pure integral assertion (E.2.3d) is not concluded.
\end{proof}

\subsection{An actual quotient lattice and the prescribed CM lifting fibre}

\begin{result}{E.2.13}{A fixed-\(\Q\) elliptic quotient outside the admissible closure}{Proved}
Let
\[
 E:\ y^2=x(x-1)(x-3),\qquad
 B:\ Y^2=X^3+8X^2+4X .
\]
Let \(\mathcal C_{\mathrm{adm}}(E)\) be the full rigid exact tensor
category generated by the actual integral Hodge structure
\(H_1(E(\C),\Z)\), using finite sums, tensor products, duals,
saturated Hodge subobjects and free Hodge quotients. Its morphisms
are all integral Hodge morphisms between the objects so obtained;
direct sums of pure weights carry their weight grading.
The actual \(H_1(B(\C),\Z)\) has no admissible presentation in this
category. An explicit integral tensor automorphism on
\(\mathcal C_{\mathrm{adm}}(E)\) fails to extend naturally across
the \(\Q\)-defined isogeny \(q:E\to B\).
\end{result}
\begin{proof}
The discriminants of the displayed equations are \(576\) and
\(12288\), respectively. Define
\[
 q(x,y)=
 \left(x-4+\frac3x,\ y\left(1-\frac3{x^2}\right)\right).
 \tag{E.2.13a}
\]
Substitution verifies the equation of \(B\). The rational map extends
between the smooth projective curves, takes the origin to the origin,
and satisfies
\[
 q^*(dX/Y)=dx/y.
\]
In the corresponding normalized complex universal covers its lift
therefore has derivative one and takes zero to zero, so it is the
map \(z\mapsto z\). This proves directly that \(q\) is a group
homomorphism. The nontrivial involution
\[
 (x,y)\longmapsto(3/x,-3y/x^2)
\]
fixes both coordinates in (E.2.13a); it is translation by
\(P=(0,0)\). Conversely \(x\) satisfies
\(x^2-(X+4)x+3=0\), and the displayed formula for \(Y\) then
recovers \(y\) generically. Thus \(q\) has degree exactly two and
kernel \(\{O,P\}\).

The curve \(E\) has
\[
 j(E)=\frac{21952}{9}.
\]
This rational number is not an algebraic integer; CM
\(j\)-invariants are algebraic integers by
\cite[Theorem 12.14 and Corollaries 12.15--12.16,
pp.119--120]{EllMilneMF}. Hence \(E\) has no CM. Its
Mumford--Tate group on \(H_1(E,\Q)\) is \(\GL_2\), by the group
calculation in A.4.9 and
\cite[Proposition 5.14]{GenKSV2026}.

Every \(T\in\mathrm{SL}(H_1(E,\Z))\) consequently defines an
integral tensor automorphism of the forgetful functor on
\(\mathcal C_{\mathrm{adm}}(E)\). To verify every lattice operation,
first act by the usual tensor representation on each integral
tensor word. Every rational Hodge subspace is
Mumford--Tate stable, and every rational Hodge morphism commutes
with this action. A saturated subobject \(N\subset P\) has
\[
 N=N_{\Q}\cap P.
\]
Therefore both \(T\) and \(T^{-1}\) preserve \(N\), and induce
mutually inverse maps on every free quotient. Induction constructs
the maps on all admissible expressions. The rational
Mumford--Tate action proves naturality for all original integral
Hodge morphisms and for isomorphisms between expressions, so the
construction is independent of the presentation. Tensor, unit and
dual compatibilities follow from the same representation.
The determinant-one condition fixes the labelled Tate object
and the elliptic polarization.

Choose a lattice basis \(e_1,e_2\) of
\(\Lambda_E=H_1(E(\C),\Z)\) in which \(P\) is represented by
\(e_1/2\). A nonzero vector modulo two has a primitive integral
lift, and such a lift extends to a basis, so this choice is possible.
The analytic quotient and the normalization \(dq=1\) identify
\[
 \Lambda_B=\Lambda_E+\Z(e_1/2).
\]
In the basis \(\mu_1=e_1/2,\ \mu_2=e_2\) of \(\Lambda_B\), the
matrix of \(q_*\) is
\[
 Q=\begin{pmatrix}2&0\\0&1\end{pmatrix}.
\]
Use the tensor automorphism constructed above with
\[
 T_E=J=\begin{pmatrix}0&-1\\1&0\end{pmatrix}.
\]
Naturality for \(q_*\) would require \(T_BQ=QJ\), and therefore
\[
 T_B=QJQ^{-1}
    =\begin{pmatrix}0&-2\\1/2&0\end{pmatrix}.
 \tag{E.2.13b}
\]
This matrix is not integral. Thus the stated extension does not
exist. An admissible presentation of the actual \(H_1(B,\Z)\)
would produce such an extension: fullness includes the integral
Hodge morphism \(q_*\), to which the constructed tensor
automorphism must be natural. This contradiction proves the
nonmembership assertion.

The same obstruction applies to a motivic admissible presentation
with these specified actual Betti lattices. Its Betti--Hodge
realization is exactly an admissible presentation of the sort just
excluded, and the isogeny supplies the algebraic correspondence
\(q_*\). Passing to cohomology uses the inverse dual
\(T=(\eta^{-1})^\vee\), so it does not change the obstruction.
The identities of curves, differentials, discriminants and matrices
are checked exactly by \texttt{verify\_numberfield\_descent.py}.

This supplies an actual failure of the additional presentation
step needed to identify E.2.12 with the full geometric diagram,
over the one fixed field \(\Q\). Enlarging the generator to
\(H_1(E)\oplus H_1(B)\) includes \(B\) and imposes the new
congruence constraint in (E.2.13b), consistently with E.2.8.
No failure of comparison between places of \(\Q\), or resolution
of the quartic full-category transporter, follows from this
automorphism obstruction.
\end{proof}

\begin{result}{E.2.14}{The common scalar and uniqueness in the prescribed CM fibre}{Proved}
For a full geometric abelian integral comparison as in E.2.9,
the units \(\lambda_A\) are one common
\(\lambda\in\widehat\Z^\times\), over all positive-dimensional
abelian varieties and all Artin-selected embedding blocks.
Fix \(\sigma\in\Gal(\Qb/F_{\mathrm{cm}})\) and prescribe its
particular CM comparison \(\alpha\) from E.2.6. There is at most
one extension of this prescribed comparison to the full geometric
abelian integral diagram. Every such extension satisfies
\[
 T_A\otimes\widehat\Z=\rho_A(\sigma)
 \qquad\hbox{for every positive-dimensional abelian variety }A.
 \tag{E.2.14a}
\]
\end{result}
\begin{proof}
Apply E.2.9 to \(A\times B\) and include the two factor
injections and projections in the diagram. Naturality gives
\[
 T_{A\times B}=T_A\oplus T_B,\qquad
 \rho_{A\times B}(\sigma)=\rho_A(\sigma)\oplus\rho_B(\sigma).
\]
Restricting the scalar equality for \(A\times B\) to each
positive-rank direct summand forces
\(\lambda_A=\lambda_{A\times B}=\lambda_B\).
For varieties over a finite extension of the original field, use
the Artin-module decomposition and selected primitive idempotents
of E.2.3/E.2.9; these preserve precisely the same factor maps.
This proves the common-scalar assertion.

Choose one CM elliptic curve \(C\) in the prescribed subdiagram.
The exact \(s=1\) formula in the main theorem of CM, used in
E.2.6, gives
\[
 \alpha_C\otimes\widehat\Z=\rho_C(\sigma).
\]
The common scalar on this nonzero factor is therefore one.
This proves (E.2.14a) on every factor.

Integral matrices inject into their profinite completions, so
(E.2.14a) determines every \(H_1\) component uniquely. The
cohomology of each abelian variety is the exterior algebra on
\(H^1\). Tensor operations, the prescribed Artin permutation,
and labelled direct-summand projectors then determine all
components in the geometric abelian diagram. This proves
the asserted uniqueness.

In particular, the rational point constructed in E.2.11 must
satisfy the additional equalities (E.2.14a) at every finite place
to yield this integral extension. The Hasse argument constructs
a point in the prescribed CM fibre; it does not establish its
equality to all the actual \(\ell\)-adic Galois comparison
points. For the main place-pair problem, other extensions
\(\sigma\) of the same field embeddings and other compatible
CM comparisons remain to be examined. The further objects of
the full pure-motive diagram remain to be included as well.
\end{proof}

\subsection{The exact integral abelian image}

Use the concrete full geometric integral diagram of E.2.3 and
E.2.9: all actual Betti lattices of abelian varieties over finite
number fields, their underlying-\(\Q\) schemes, Artin algebra
and module maps, correspondences, tensor operations and specified
direct-summand lattices. Write
\[
 \mathcal G_{\mathrm{ab},\Z}
   =\Aut^\otimes(\omega_{B,\Z}),\qquad
 \operatorname{art}:\mathcal G_{\mathrm{ab},\Z}\longrightarrow
       \Gamma_{\Q}=\Gal(\Qb/\Q),\qquad
 I_{\mathrm{ab}}=\operatorname{im}(\operatorname{art}).
\]
The map is the restriction to finite Artin algebras, with the
primitive-idempotent description proved in E.2.3.
These are statements about that specified integral diagram.
Restriction from the full pure diagram gives an inclusion of its
Artin image in \(I_{\mathrm{ab}}\).

\begin{result}{E.2.15}{The Artin kernel and the uniform CM-field restriction}{Proved}
There is an exact central extension of abstract groups
\[
 1\longrightarrow\{1,w\}\longrightarrow
 \mathcal G_{\mathrm{ab},\Z}\longrightarrow I_{\mathrm{ab}}
 \longrightarrow1,
 \tag{E.2.15a}
\]
where \(w\) acts by \((-1)^m\) on every pure weight-\(m\)
summand. Thus every Artin fibre has at most two elements.
Let \(\Q^{\mathrm{tr}}\subset\R\) be the compositum of all
finite totally real fields. Then
\[
 \Q^{\mathrm{cm}}=\Q^{\mathrm{tr}}(i),\qquad
 I_{\mathrm{ab}}\subset\Gal(\Qb/\Q^{\mathrm{tr}}).
 \tag{E.2.15b}
\]
Consequently each element of \(I_{\mathrm{ab}}\) restricts to
either the identity or actual complex conjugation on the entire
field \(\Q^{\mathrm{cm}}\), with one choice for all CM fields.
\end{result}
\begin{proof}
Let \(g\) have trivial Artin restriction. The selected-block
construction of E.2.3 permits E.2.9 on every abelian variety over
\(\Qb\): each geometric cyclic quotient is defined over some
finite extension, and its underlying-\(\Q\) object and Artin
module maps belong to the diagram. Hence
\[
 T_A\otimes\widehat\Z=\lambda_A I .
\]
Products and the factor inclusions/projections force one common
\(\lambda\), by E.2.14. On any nonzero factor, injectivity
\(\Z\hookrightarrow\widehat\Z\) makes \(T_A=nI\) for an integer
\(n\); its inverse is integral, so \(n=1\) or \(-1\).
This same sign acts on every geometric embedding block.
Integral cohomology of an abelian variety is the exterior algebra
on \(H^1\), so the action on \(H^r\) is \(n^r\).
The Artin action is the identity, and tensor operations and
labelled projectors determine all remaining components.
Thus \(g=1\) or \(w\).

Conversely the displayed weight sign is a tensor automorphism:
weights add under tensor product, negate under duality, and all
morphisms respect them. It preserves the integral lattices and
fixes Artin and Tate objects. It is nontrivial on an elliptic
\(H^1\), and it is central because the weight grading is
functorial. This proves (E.2.15a).

For a finite totally real field \(F\), apply the finite separating
diagram of E.2.5 to its inclusion in \(\Qb\) and to
\(\sigma|_F\), where \(\sigma=\operatorname{art}(g)\).
The underlying-\(\Q\) objects and the Artin algebra \(A_F\)
belong to the present diagram. The exact module-block map of
E.2.3 gives a comparison between these two \(F\)-realizations.
E.2.5 forces the embeddings to agree. Thus \(\sigma\) fixes
\(\Q^{\mathrm{tr}}\).

Every totally real \(F\) has CM extension \(F(i)\), proving one
inclusion in the field equality. Conversely, for a CM field
\(K\), choose a nonzero element \(b\) with \(\bar b=-b\) and
put \(\delta=-b^2\) in its maximal totally real subfield \(K^+\).
Then \(\delta\) is totally positive and
\(K=K^+(\sqrt{-\delta})\).
All conjugates of \(\sqrt{\delta}\) are real, so
\(\sqrt{\delta}\in\Q^{\mathrm{tr}}\) and
\(K\subset\Q^{\mathrm{tr}}(i)\).
This proves (E.2.15b). An automorphism fixing \(\Q^{\mathrm{tr}}\)
sends \(i\) to \(i\) or \(-i\), and these two maps are precisely
the claimed uniform restrictions.
\end{proof}

\begin{result}{E.2.16}{An exact lattice criterion for the abelian integral image}{Proved}
Let \(F_{\mathrm{cm}}\) be the field of E.2.6 and
\(\Gamma_{F_{\mathrm{cm}}}=\Gal(\Qb/F_{\mathrm{cm}})\).
For an abelian variety \(A/\Qb\), put
\(\Lambda_A=H_1(A(\C),\Z)\), included in its profinite
completion. Using the torsion comparison \(\kappa_{A,n}\)
of E.2.9, define
\[
 \rho_A(\sigma):
 \Lambda_A\otimes\widehat\Z
 \longrightarrow\Lambda_{\sigma A}\otimes\widehat\Z .
\]
Then
\[
 I_{\mathrm{ab}}\cap\Gamma_{F_{\mathrm{cm}}}
 =
 \left\{\sigma\in\Gamma_{F_{\mathrm{cm}}}:
   \rho_A(\sigma)(\Lambda_A)=\Lambda_{\sigma A}
   \text{ for every }A/\Qb\right\}.
 \tag{E.2.16a}
\]
Inclusion of lattices in place of equality gives the same set.
For each element of this set there are exactly two integral
abelian lifts, differing by \(w\); the one whose CM restriction
is the canonical comparison of E.2.6 is uniquely determined.
\end{result}
\begin{proof}
For an integral comparison with Artin restriction \(\sigma\),
E.2.9 and the product maps give
\[
 T_A\otimes\widehat\Z=\rho_A(\sigma)\lambda
 \tag{E.2.16b}
\]
with one common unit \(\lambda\).
On a CM elliptic curve \(C\), E.2.6 supplies the actual
integral comparison \(\alpha_C\) satisfying
\(\alpha_C\otimes\widehat\Z=\rho_C(\sigma)\).
Thus \(\alpha_C^{-1}T_C\) is an integral matrix whose
completion is \(\lambda I\); its integral inverse forces
\(\lambda=1\) or \(-1\).
This proves the necessary lattice equality without having
prescribed the CM restriction beforehand.

Conversely, suppose \(\sigma\) belongs to the right-hand side.
Define \(T_A\) by restricting \(\rho_A(\sigma)\) to
\(\Lambda_A\). It is the unique integral isomorphism with
that completion. Even an inclusion suffices: its integer
matrix has determinant a unit at every prime, so its determinant
is \(1\) or \(-1\).

For every homomorphism \(f:A\to B\) over \(\Qb\), the
actual étale identity
\[
 \rho_B(\sigma)f_*=(\sigma f)_*\rho_A(\sigma)
\]
restricts to the corresponding equation of integral Betti
matrices, since the finite-free Hom group injects into its
completion. The same argument treats product maps and all
tensor compatibilities. For cohomology use
\(\eta_A=(T_A^{-1})^\vee\), and on \(H^r\) its exterior
power. Their completions are the actual Galois maps in every
degree by compatibility with cup products.

Every algebraic correspondence has its functorial étale cycle
class, so its naturality equation holds after completion at
every prime and hence over \(\Q\). The absolute-Hodge
abelian version uses the same defining compatibility of its
étale tensor components. Therefore all labelled projectors
commute with the constructed maps. For a saturated summand
lattice use \(p(V_{\Q})\cap\Lambda\), and for an image
convention use \(p(\Lambda)\); the equation with \(p\) maps
either precisely to its target lattice. Tensor products and
duals retain integrality.

For a finite Artin algebra use \(e_u\mapsto e_{\sigma u}\).
If \(A/L\) is a finite model, its underlying-\(\Q\) object
has Betti realization the direct sum over \(u:L\to\Qb\).
Use the maps \(T_{uA}\) on these blocks. They preserve the
Artin-module equations and
\(\sum_u e_u\otimes e_u\).
Enlarging \(L\) or changing the finite model gives the same
maps, by the naturality already verified. Thus E.2.3 descends
this construction to an element of
\(\mathcal G_{\mathrm{ab},\Z}\) with Artin restriction
\(\sigma\).

Since \(\Q^{\mathrm{ab}}\subset F_{\mathrm{cm}}\), its Tate
map is the identity. Equality of completions shows that its
CM restriction is \(\alpha\).
Multiplication by \(w\) gives the second lift; E.2.15 excludes
any others. This proves (E.2.16a).

The criterion removes a separate tensor-coherence obligation
once actual all-\(A\) lattice integrality has been established.
It does not establish that integrality for an extension of the
quartic pair, or extend an abelian comparison to the further
objects of the full pure diagram.
\end{proof}

\begin{result}{E.2.17}{An explicit excluded Galois element with the quartic swap}{Proved}
There is a specified \(\sigma\in\Gamma_{F_{\mathrm{cm}}}\)
whose action on the quartic splitting field of E.2.6 is
\((12)(34)\), and which is outside \(I_{\mathrm{ab}}\).
For the explicit curve
\[
 E:y^2=x(x-1)(x-3)
\]
and a fixed integral Betti basis it satisfies
\[
 \rho_{E,5}(\sigma)=
 U(a):=\begin{pmatrix}1&a\\0&1\end{pmatrix},
 \qquad a=\sum_{n\ge1}5^{n!}\in\Z_5\setminus\Q .
 \tag{E.2.17a}
\]
In particular this element is also excluded from the Artin
image of the full pure integral diagram.
\end{result}
\begin{proof}
Fix the Betti basis geometrically as follows. Let \(\gamma_1\)
traverse the real interval \(0\le x\le1\) on
\(y=+\sqrt{x(x-1)(x-3)}\), then return on the negative
sheet. Let \(\gamma_2\) traverse \(1\le x\le3\) on
\(y=-i\sqrt{-x(x-1)(x-3)}\), then return on the positive
imaginary sheet. These are oriented closed cycles on the
branched double cover. They meet only above \(x=1\); in
the complex local coordinate \(y\), their oriented tangents
are respectively the negative real and negative imaginary
axes. Their intersection is therefore \(+1\).
The double cover has four simple branch points and genus
one; the unimodular intersection form on its rank-two
homology shows that these two classes form an integral
basis. Use this ordered basis, and its images under the
canonical torsion comparison, throughout (E.2.17a).

We first prove
\(\rho_{E,5}(\Gamma_{\Q})=\GL_2(\Z_5)\).
At \(3\), the equation has \(c_4=112\) a unit and
\(v_3(\Delta)=v_3(576)=2\).
The reduction \(y^2=x^2(x-1)\) has an ordinary node whose
tangent lines split over the unramified quadratic extension.
The equation is minimal, since its discriminant valuation is
less than twelve. Over that unramified extension the curve
is a Tate curve with parameter \(q\) of valuation two:
\(j(q)=q^{-1}+744+\cdots\) and \(v_3(j(E))=-2\).
The precise Galois-equivariant uniformization and this
expansion are given in \cite[Section19, p.44]{GenConradTate}.

Its \(5^n\)-torsion is generated by the classes of
\(\zeta_{5^n}\) and \(q^{1/5^n}\) in
\(\overline{\Q}_3^\times/q^\Z\).
Over the maximal unramified extension, roots of unity of
order prime to three and \(5^n\)-th roots of units are
available. Thus the Kummer character of \(q\) is twice
the tame \(5\)-adic character of a uniformizer.
Since two is a \(5\)-adic unit, the inertia image contains
the entire group \(U(\Z_5)\).

At \(7\), direct counting gives
\[
 E(\mathbf F_7)=\{O,(0,0),(1,0),(3,0)\}.
\]
The Frobenius polynomial \(T^2-4T+7\) has discriminant
\(3\) modulo five, a nonsquare. Hence the representation
modulo five is irreducible.
There is therefore an element \(g\) of the image carrying
the primitive \(U\)-fixed vector to a vector independent
modulo five. These vectors form a \(\Z_5\)-basis.
In it \(U(\Z_5)\) and \(gU(\Z_5)g^{-1}\) are the full
upper and lower root groups: their parameters change by
units. This is the elementary argument of
\cite[pp.265--266]{GenConradWong1999}, with the
valuation-two Kummer calculation supplied above.

The two root groups generate \(\mathrm{SL}_2(\Z_5)\).
Indeed elementary row operations use a unit in the first
column to clear both off-diagonal entries; the remaining
diagonal follows from
\[
 U(u)L(-u^{-1})U(u)
   =\begin{pmatrix}0&u\\-u^{-1}&0\end{pmatrix}
 \quad(u\in\Z_5^\times)
\]
and the same formula at \(u=1\).
The determinant is the surjective cyclotomic character, so
the full image is \(\GL_2(\Z_5)\).

Put \(N=\Gamma_{\Q^{\mathrm{cm}}}\).
Actual complex conjugation \(c\) is central on
\(\Q^{\mathrm{cm}}\), hence
\([\gamma,c]\in N\) for every \(\gamma\in\Gamma_{\Q}\).
Its action on \(T_5E\) is an involution of determinant
\(-1\). The projectors \((1+c)/2,(1-c)/2\) give a
\(\Z_5\)-basis with \(c=\operatorname{diag}(1,-1)\).
The full image just proved supplies each \(U(t)\), and
\[
 [U(t),c]=U(2t).
\]
Thus \(\rho_{E,5}(N)\) contains the upper root group.
It is normal in \(\GL_2(\Z_5)\), so it contains the lower
one and hence \(\mathrm{SL}_2(\Z_5)\).

Let \(N'\) be the closure of the commutator subgroup of
\(N\). Since \(F_{\mathrm{cm}}/\Q^{\mathrm{cm}}\) is
abelian, \(N'\subset\Gamma_{F_{\mathrm{cm}}}\).
Compactness and continuity give the exact identity
\[
 \rho_{E,5}(N')
  =\overline{[\rho_{E,5}(N),\rho_{E,5}(N)]}.
\]
For one inclusion the image is compact and contains all
image commutators; for the other, take limits of products
of commutators before applying \(\rho_{E,5}\).
The group \(\mathrm{SL}_2(\Z_5)\) is perfect:
with \(D=\operatorname{diag}(2,1/2)\),
\[
 [D,U(t)]=U(3t),
\]
so every upper root element is a commutator, and conjugation
gives the lower root group. Their generation proves the
claim. Since the determinant is one on
\(\Gamma_{F_{\mathrm{cm}}}\), we obtain
\[
 \rho_{E,5}(\Gamma_{F_{\mathrm{cm}}})
       =\mathrm{SL}_2(\Z_5).
 \tag{E.2.17b}
\]

The number \(a\) in (E.2.17a) is irrational. For its partial
sum \(a_N\), one has \(v_5(a-a_N)=(N+1)!\).
An equality \(a=r/s\) with \(5\nmid s\) would make the
nonzero integer \(r-sa_N\) divisible by \(5^{(N+1)!}\),
while \(|r-sa_N|\le C5^{N!}\) for all large \(N\).
This is impossible.
Equation (E.2.17b) therefore gives an element with the
specified nonintegral matrix. E.2.16 excludes that element
from \(I_{\mathrm{ab}}\).

It remains to impose the quartic action simultaneously.
Let \(L\) and \(H=L^{V_4}\) be as in E.2.6.
The proved intersection \(L\cap F_{\mathrm{cm}}=H\)
gives a surjection \(\Gamma_{F_{\mathrm{cm}}}\to V_4\).
The joint image with \(\rho_{E,5}\) is a closed subdirect
product of \(V_4\) and \(\mathrm{SL}_2(\Z_5)\).
Goursat's quotient construction identifies any proper
subdirect product with a common nontrivial quotient of
these groups. A quotient of \(V_4\) is abelian, whereas
a quotient of the perfect group
\(\mathrm{SL}_2(\Z_5)\) is perfect; an abelian perfect
group is trivial. The joint image is therefore the full
product. We prescribe \(((12)(34),U(a))\).

To specify one resulting \(\sigma\) completely, enumerate
the nonconstant integer polynomials by increasing maximum
of degree and coefficient height, with lexicographic ties.
Let \(R_n\) contain all roots of the first \(n\) polynomials,
all affine coordinates of \(E[5^n]\), and all roots of
\(x^4-x-1\), and put \(\Omega_n=\Q(R_n)\).
These form a nested finite Galois exhaustion of \(\Qb\).
Order \(R_n\) lexicographically by real and imaginary parts
in the fixed complex embedding, thereby ordering
\(\Gal(\Omega_n/\Q)\) by its permutation tuples.
The set of Galois elements fixing \(F_{\mathrm{cm}}\)
and having the prescribed pair of images is nonempty
and compact. At each stage choose the first permutation
extending the preceding choice whose cylinder meets this
compact set. The resulting nested cylinders have exactly
one common element because the \(\Omega_n\) exhaust
\(\Qb\). This defines \(\sigma\), without a computability
assertion.

The finite point count, valuation inputs and root-group
identities are checked exactly in
\texttt{verify\_numberfield\_descent.py}.
The construction excludes an actual extension of the
quartic embedding pair. Other extensions remain to be
tested, so the full place-pair question is still open.
\end{proof}

\begin{result}{E.2.18}{The abelian image for all remaining Galois candidates}{Proved}
Put
\[
 D=\mu_2(\widehat\Z)
   =\{\lambda\in\widehat\Z^\times:\lambda^2=1\}
   =\prod_\ell\{1,-1\},
\]
and define
\[
 \Sigma=\left\{\sigma\in\Gamma_{\Q^{\mathrm{cm}}}:
 \begin{array}{l}
  \text{there exists }\lambda\in D\text{ such that}\\
  \lambda\rho_A(\sigma)(\Lambda_A)=\Lambda_{\sigma A}
       \text{ for every abelian variety }A/\Qb
 \end{array}\right\}.
 \tag{E.2.18a}
\]
For the concrete full geometric integral abelian diagram,
\[
 I_{\mathrm{ab}}=\Sigma\ \sqcup\ c\Sigma,
 \tag{E.2.18b}
\]
where \(c\) is actual complex conjugation. Inclusion of lattices
in (E.2.18a) gives the same set. The sign pattern \(\lambda\)
is unique modulo the one global sign \(\{1,-1\}\), and its class
defines an abstract group homomorphism
\[
 \delta:\Sigma\longrightarrow
       \mu_2(\widehat\Z)/\{1,-1\}.
 \tag{E.2.18c}
\]
There are precisely two integral abelian lifts of each element
of \(I_{\mathrm{ab}}\). On
\(\Sigma\cap\Gamma_{F_{\mathrm{cm}}}\), the class \(\delta\)
is trivial and the criterion is exactly E.2.16.
\end{result}
\begin{proof}
Actual analytic conjugation on all underlying-\(\Q\) schemes
gives an integral tensor comparison with Artin restriction \(c\):
it carries the \(u\)-component to the \(cu\)-component and
respects every \(\Q\)-defined correspondence and product.
Its square is the identity. E.2.15 therefore reduces necessity
to \(\sigma\in I_{\mathrm{ab}}\cap\Gamma_{\Q^{\mathrm{cm}}}\).
For such an element, E.2.9 and the product maps of E.2.14 give
one \(\lambda\in\widehat\Z^\times\) satisfying
\[
 T_A\otimes\widehat\Z=\rho_A(\sigma)\lambda
 \quad\text{for every }A/\Qb.
 \tag{E.2.18d}
\]

Use the fixed curve \(C:y^2=x^3-x\) over \(\Q\), with labelled
\([i]:(x,y)\mapsto(-x,iy)\) over \(\Q(i)\).
The Artin-selected block is unchanged because \(\sigma\)
fixes \(\Q(i)\). Thus \(T_C\) is an integral unimodular
endomorphism of \(\Lambda_C\) commuting with \([i]_*\).
The lattice is a rank-one torsion-free \(\Z[i]\)-module and
is free because \(\Z[i]\) is Euclidean. In a basis \((e,ie)\),
commuting with \(J=\left(\begin{smallmatrix}0&-1\\1&0\end{smallmatrix}\right)\)
means
\[
 T_C=\begin{pmatrix}r&-s\\s&r\end{pmatrix},
 \quad r,s\in\Z,\qquad
 \det T_C=r^2+s^2=1.
\]
Every root of unity belongs to \(\Q^{\mathrm{cm}}\), so
the cyclotomic character \(\chi(\sigma)\) is one.
The Weil pairing gives
\(\det\rho_C(\sigma)=\chi(\sigma)=1\).
Taking determinants in (E.2.18d) yields \(\lambda^2=1\).
The exterior-square polarization identification with the
integral Tate line also shows that the comparison fixes
that line. This proves necessity without adding a Hodge
filtration or a real-positivity assumption.

Conversely let \(\sigma,\lambda\) satisfy (E.2.18a).
Restrict \(\lambda\rho_A(\sigma)\) to the actual lattice
to obtain \(T_A\). An inclusion already gives an integral
isomorphism: its determinant is an integer invertible at
every prime, hence \(1\) or \(-1\).
For a homomorphism \(f:A\to B\), the completed equation
\[
 (\lambda\rho_B(\sigma))f_*
       =(\sigma f)_*(\lambda\rho_A(\sigma))
\]
descends to integral Betti matrices. Product maps satisfy
the same equation.

In contravariant cohomology put
\(\eta_A=(T_A^{-1})^\vee\). Its completion is the actual
Galois cohomology map multiplied by
\(\lambda^{-1}=\lambda\). On \(H^r(A)\), its exterior power
therefore contributes \(\lambda^r\).
Declare the Tate component to be the identity.
A correspondence of codimension \(d+j\) from a
\(d\)-dimensional variety carries \(H^r\) to \(H^{r+2j}\),
with the corresponding Tate twist. Its two scalar factors
coincide since
\[
 \lambda^{r+2j}=\lambda^r.
\]
The actual étale cycle-class equation consequently proves
the required Betti naturality by faithful completion.
For tensor words and duals, the parity of the number of
\(H^1\) factors is the weight modulo two; Tate twists change
it by an even integer. Thus the same equality verifies
every homogeneous correspondence and polarization.
In the absolute-Hodge abelian version, its defining étale
tensor compatibility gives the same verification for
weight-zero tensors.

Commutation with each labelled projector maps both the
saturated convention \(p(V_\Q)\cap\Lambda\) and the image
convention \(p(\Lambda)\) precisely to their target lattices.
Tensor products and duals preserve integrality.
For an Artin algebra put \(e_u\mapsto e_{\sigma u}\);
on a finite model \(A/L\), use \(T_{uA}\) on each
embedding block. The Artin-module equations and the
diagonal tensor projector \(\sum e_u\otimes e_u\)
are preserved. Naturality makes the construction
independent of the finite model, and the exact descent
of E.2.3 gives an element of
\(\mathcal G_{\mathrm{ab},\Z}\).
This proves sufficiency and (E.2.18b).

If \(\lambda,\mu\) work for the same \(\sigma\), the ratio
of the constructed comparisons has trivial Artin
restriction. E.2.15 makes it the identity or global
weight sign, so \(\mu/\lambda\) is one global \(1\)
or \(-1\). The composition identity
\[
 \rho_{\tau A}(\sigma)\rho_A(\tau)=\rho_A(\sigma\tau)
\]
shows that these sign classes multiply, proving
(E.2.18c). E.2.15 also gives the assertion about the
two lifts. For \(\sigma\in\Gamma_{F_{\mathrm{cm}}}\),
the integral CM anchor in E.2.16 forces the sign
pattern to be global, as claimed.

This describes every abelian Artin candidate.
E.2.20 below proves that \(\delta\) is trivial and
removes the local-sign variable from the criterion.
E.2.21 then determines the exact CM-field restriction.
It does not prove \(\Sigma=\{1\}\), produce a swapping
quartic element of \(\Sigma\), or extend one to the
full pure diagram. The main full-pure separation
question is therefore retained.
\end{proof}

\begin{result}{E.2.19}{Ring-class restrictions and a fixed-field local-sign example}{Proved}
Every \(\sigma\in\Sigma\) of E.2.18 fixes every imaginary-quadratic
ring class field pointwise. There nevertheless exists an actual
\(\sigma\in\Gamma_{\Q^{\mathrm{cm}}}\), nontrivial on
\(\Q(i)^{\mathrm{ab}}\), with an integral tensor comparison on the
following precise diagram: all CM elliptic curves over \(\Qb\)
with geometric endomorphism algebra \(\Q(i)\), all actual abelian
variety lattices isogenous to products of those curves, their
tensor constructions and specified summands, algebraic
correspondences, and all Artin objects with every finite-model
module map. The same construction respects labelled
absolute-Hodge tensors in the abelian version.
Its common scalar is
\[
 \lambda_5=-1,\qquad \lambda_\ell=1\quad(\ell\ne5),
 \tag{E.2.19a}
\]
which is not one global sign. No membership of this \(\sigma\)
in the all-abelian group \(\Sigma\) is asserted.
\end{result}
\begin{proof}
Let \(R=K(j(E))\) be any imaginary-quadratic ring class field.
Apply the finite separating diagram of E.2.2 over \(R\), regard
its objects as underlying-\(\Q\) schemes, and include its
Artin algebra \(A_R\) and module maps. E.2.3 sends the identity
embedding block to the \(\sigma|_R\) block. E.2.2 forces this
restriction to be the identity or actual conjugation. Since
\(\sigma\) fixes \(K\subset\Q^{\mathrm{cm}}\), the latter is
excluded. This proves the first assertion by the explicit
Artin-block maps.

For every imaginary-quadratic \(K\), one has
\[
 K^{\mathrm{ab}}\cap\Q^{\mathrm{cm}}=\Q^{\mathrm{ab}}.
 \tag{E.2.19b}
\]
The inclusion from right to left follows because \(K/\Q\)
is abelian and cyclotomic fields are CM. For the reverse
inclusion, the intersection is Galois over \(\Q\).
The group \(\Gal(K^{\mathrm{ab}}/\Q)\) is generated by its
abelian index-two subgroup \(\Gal(K^{\mathrm{ab}}/K)\) and
actual complex conjugation \(c\), which has order two.
On the intersection the image of \(c\) is central, because
it is central on \(\Q^{\mathrm{cm}}\). The resulting quotient
group is therefore abelian, proving (E.2.19b).

Now take \(K=\Q(i)\), and let \(s\) be the scalar finite
\(K\)-idele (E.2.19a), with infinite component \(1\).
Its norm to \(\Q\) is exactly \(1\), since each local
scalar norm is \(\lambda_\ell^2\). Put
\[
 g_K=\operatorname{art}_K(s)\in\Gal(K^{\mathrm{ab}}/K).
\]
The Artin norm/restriction diagram makes \(g_K\) trivial
on \(\Q^{\mathrm{ab}}\)
\cite[Chapter V, Proposition 5.2 and its following diagram,
pp.~178--179]{GenMilneCFT2020}.
By (E.2.19b), it is compatible with the identity on
\(\Q^{\mathrm{cm}}\). These two maps define an automorphism
of their compositum and extend to \(\Qb\).
One particular extension is specified by the fixed finite
Galois exhaustion of E.2.17: recursively take the first
restriction whose cylinder meets the nonempty compact set
of these compatible extensions. This uses no computability
claim.

The restriction \(g_K\) is nontrivial. Modulo the ray
modulus \(15\), the idele \(s\) represents the rational
residue \(4\): it is \(1\) modulo \(3\) and \(-1\) modulo
\(5\). The ray-class exact sequence has kernel the image
of \(\Z[i]^\times=\{1,-1,i,-i\}\). Only \(1\) among these
units is \(1\) modulo \(3\), and it is not \(-1\) modulo
\(5\). Thus this ray class is nontrivial
\cite[Chapter V, Theorem 1.7 and proof,
pp.~150--151]{GenMilneCFT2020}.
In particular \(\sigma\notin\Gamma_{F_{\mathrm{cm}}}\).
For every order \(\mathcal O_f=\Z+f\Z[i]\), however,
\(s\in\widehat{\mathcal O_f}^{\,\times}\); hence its ring
class is trivial. This is also the rational-congruence
subgroup in the ideal quotient of
\cite[Theorem 9.18, pp.~190--192]{GenCox1989}.

Let \(A/\Qb\) be any elliptic curve of this fixed CM field.
The two CM types have reflex field \(K\); their reflex
norms are the identity and conjugation, both of which
send \(s\) to the same \(\lambda\).
The exact main theorem
\cite[Theorem 9.10, p.~78]{GenMilneCM2020} therefore gives
a \(K\)-linear quasi-isogeny
\[
 \alpha_A:A\longrightarrow\sigma A,\qquad
 \alpha_A(\lambda x)=\sigma x,
 \qquad
 \widehat\alpha_A=\lambda\rho_A(\sigma).
 \tag{E.2.19c}
\]
It is an actual isomorphism. Indeed, write
\(\alpha_A=f/m\), with \(f\) an honest homomorphism.
Integrality on every Tate lattice makes \(f\) vanish
on \(A[m]\); it factors through \([m]\), proving that
\(\alpha_A\) is honest. Apply the same argument to
its inverse. The formula is independent of the labelling
of \(K\), because both types give the same scalar and
completed homology is faithful on homomorphisms.

For every homomorphism \(f:A\to B\), the equality
\(\alpha_Bf=(\sigma f)\alpha_A\) follows after completion
from Galois naturality and the common scalar, and hence
holds exactly. Products and isogenies therefore extend
the construction to every actual lattice isogenous to
a product of these elliptic curves. More explicitly,
transport the product quasi-isogeny through any chosen
isogeny; its completion is again \(\lambda\rho\).
The preceding integrality argument makes it an actual
isomorphism and proves independence of that choice.

On cohomology use inverse duals and exterior powers.
The completed factor in degree \(r\) is \(\lambda^r\).
Correspondences change degree by an even integer, so
their two factors coincide. Their étale cycle-class
equations descend faithfully to the required Betti
equations. This also treats tensors, duals, the stated
projector lattices and labelled absolute-Hodge tensors.
The polarization multiplier is \(1\), since both the
cyclotomic character and the idele norm are \(1\)
\cite[Remark 9.11(c), p.~78]{GenMilneCM2020}.
The Tate component is the identity.

Finally an object defined over an arbitrary finite
field \(L\) has underlying-\(\Q\) Betti realization
\(\bigoplus_u H_B(uA)\). Every \(uA\) still has the
same imaginary-quadratic CM field, though its type
may be conjugated. Both types have already been
treated. Use \(\alpha_{uA}\) from the \(u\)-block to
the \(\sigma u\)-block and put \(e_u\mapsto e_{\sigma u}\)
on every Artin algebra. Naturality proves the module
equations, preserves \(\sum e_u\otimes e_u\), and
commutes with every change of finite model.
The exact descent of E.2.3 gives the asserted global
comparison on this specified fixed-field diagram.
\end{proof}

\begin{result}{E.2.20}{Two CM types eliminate the local-sign ambiguity}{Proved}
The homomorphism \(\delta\) in E.2.18 is trivial. More precisely,
an actual comparison on the full geometric CM/Artin diagram
with Artin restriction in \(\Gamma_{\Q^{\mathrm{cm}}}\)
and common scalar \(\lambda\in\prod_\ell\{1,-1\}\)
forces \(\lambda\) to be one global \(1\) or \(-1\).
Consequently E.2.18 becomes the exact criterion
\[
 \begin{split}
 \Sigma
  &=\{\sigma\in\Gamma_{\Q^{\mathrm{cm}}}:
       \rho_A(\sigma)(\Lambda_A)=\Lambda_{\sigma A}
       \text{ for every abelian variety }A/\Qb\},\\
 I_{\mathrm{ab}}&=\Sigma\sqcup c\Sigma.
 \end{split}
 \tag{E.2.20a}
\]
The sign in E.2.19 is already excluded by two CM types
of the explicit degree-eight field
\[
 L_{15}=\Q(\zeta_{15}+\zeta_{15}^{-1},i).
 \tag{E.2.20b}
\]
\end{result}
\begin{proof}
Fix \(m\ge3\), and write
\[
 F=\Q(\zeta_m+\zeta_m^{-1}),\qquad L=F(i).
\]
Then \(L/\Q\) is Galois and CM, with \(L^{\langle c\rangle}=F\).
Choose a CM type \(\Phi\) containing the identity embedding,
and let \(\Phi'\) replace that embedding by \(c\), leaving
the other conjugate pairs unchanged. The CM construction
gives abelian varieties \(A_\Phi,A_{\Phi'}\) with labelled
\(L\)-action and rational homology of rank one over \(L\).
This includes imprimitive types: the polarized torus
\(\C^\Phi/\Phi(\mathcal O_L)\) is supplied by
\cite[Example 2.11]{GenMilneCM2020}, and
\cite[Proposition 7.9 and Corollary 7.10,
pp.~54--55]{GenMilneCM2020} gives a model and its morphisms
over \(\Qb\). No equality of its full endomorphism algebra
with \(L\) is required. The labelled action and every
map are defined over a finite field, so their actual
Artin-selected blocks are available by E.2.3.

Choose one finite \(L\)-idele \(s\) whose Artin image is
\(\sigma|_{L^{\mathrm{ab}}}\). The global Artin map is
surjective, and the infinite ideles of this totally
imaginary field are connected and have trivial Artin
image, so finite ideles suffice.
For a CM type \(\Psi\) of \(L\), its reflex field
\(E_\Psi\) is contained in \(L\). Norm compatibility
identifies the Artin image of
\(\operatorname{Nm}_{L/E_\Psi}(s)\) with
\(\sigma|_{E_\Psi^{\mathrm{ab}}}\); here
\(E_\Psi^{\mathrm{ab}}L/L\) is abelian and contained
in \(L^{\mathrm{ab}}\), which is the precise
restriction map being used.

The exact norm formulas are
\[
 N_{L,\Psi}
  =N_\Psi\circ\operatorname{Nm}_{L/E_\Psi},
 \qquad
 N_{L,\Psi}(s)=\prod_{\psi\in\Psi}\psi^{-1}(s).
 \tag{E.2.20c}
\]
The first is \cite[Proposition 1.23, p.~16]{GenMilneCM2020};
the second is \cite[Proposition 1.26, p.~17]{GenMilneCM2020}
with \(k=E=L\). Remark 1.25 makes these morphisms of
tori and hence identities on finite ideles.
In particular the norm to the reflex field is already
included; there is no extra stabilizer multiplicity.

The main theorem gives an \(L\)-linear quasi-isogeny
\(\alpha_\Psi:A_\Psi\to\sigma A_\Psi\) with
\[
 \widehat\alpha_\Psi
     =\rho_{A_\Psi}(\sigma)N_{L,\Psi}(s)^{-1}.
\]
The Betti comparison \(T_{A_\Psi}\) is also \(L\)-linear
relative to the two labelled actions, because its
Artin-block naturality is
\[
 T_{A_\Psi}\,\iota(a)=(\sigma\iota(a))\,T_{A_\Psi}
 \quad(a\in L).
\]
Thus \(\alpha_\Psi^{-1}T_{A_\Psi}\) is an \(L\)-linear
automorphism of a rank-one \(L\)-vector space, hence
multiplication by an element \(b_\Psi\in L^\times\).
After completion this says
\[
 \lambda N_{L,\Psi}(s)=b_\Psi\in L^\times.
 \tag{E.2.20d}
\]
This centralizer calculation is valid for imprimitive
types and uses no Hodge preservation of the original
integral comparison.

Take the quotient of (E.2.20d) for \(\Phi,\Phi'\).
Only the identity/conjugate pair changes, so
\[
 r=s/c(s)\in L^\times,\qquad r\,c(r)=1.
\]
An explicit Hilbert-90 solution is
\(b=1+r\) when \(r\ne-1\), and \(b=i\) when \(r=-1\);
in both cases \(b/c(b)=r\).
Consequently \(t=s/b\) is fixed by \(c\), and therefore
\(t\in\mathbb A_{F,f}^{\times}\).
Each CM type contains one extension of every embedding
of \(F\), so (E.2.20c) gives
\[
 N_{L,\Psi}(t)=\operatorname{Nm}_{F/\Q}(t).
\]
Dividing (E.2.20d) by the principal element
\(N_{L,\Psi}(b)\) shows that
\(\lambda\operatorname{Nm}_{F/\Q}(t)\in L^\times\).
It is also a scalar rational finite idele. A diagonal
element of \(L\) with this property is fixed by every
element of \(\Gal(L/\Q)\), hence lies in \(\Q^\times\).
For some \(a\in\Q^\times\), we obtain
\[
 \operatorname{Nm}_{F/\Q}(t)=a\lambda.
 \tag{E.2.20e}
\]

The Artin image of the left side restricts trivially
to \(F\), by the norm/restriction diagram.
The finite principal idele \(a\), with infinite
component \(1\), acts on \(\Q^{\mathrm{ab}}\) as
the identity or actual complex conjugation:
insert its missing infinite component and use
global reciprocity. Both actions are trivial on
the totally real field \(F\).
Equation (E.2.20e) therefore makes
\(\operatorname{art}_{\Q}(\lambda)\) trivial on
\(\Q(\zeta_m+\zeta_m^{-1})\). Equivalently,
\[
 \lambda\bmod m\in\{1,-1\}.
\]
This holds for every \(m\). Their real cyclotomic
fields exhaust \(\Q^{\mathrm{ab},+}\), and
\(\Gal(\Q^{\mathrm{ab}}/\Q^{\mathrm{ab},+})
=\{1,c\}\). Cyclotomic reciprocity identifies these
two elements with the two global units \(1,-1\).
Thus \(\lambda\) is global, as asserted.

For E.2.19 the residue modulo \(15\) is \(4\), which
is neither \(1\) nor \(-1\). The finite obstruction
can be specified completely. Put
\[
 t=\zeta_{15}+\zeta_{15}^{-1}
   =2\cos(2\pi/15),\qquad L=F_{15}(i),\qquad
 \Lambda=\mathcal O_L.
\]
Let \(\Phi_0\) be all embeddings fixing \(i\), and let
\(\Phi_1\) replace its identity embedding by \(c\).
The first torus is isogenous to a fourth power of
the Gaussian CM elliptic curve: an \(L\)-lattice,
viewed as a rank-four \(\Q(i)\)-lattice with this
constant type, is commensurable with \(\Z[i]^4\).
It therefore already belongs to E.2.19's diagram.
The additional dimension-four CM variety is
\[
 A_1(\C)=\C^{\Phi_1}/\Phi_1(\mathcal O_L),\qquad
 \psi(x,y)=\operatorname{Tr}_{L/\Q}
       \bigl(i(3-2t)x\,c(y)\bigr).
 \tag{E.2.20f}
\]
Here \(t\) is the largest root of
\(X^4-X^3-4X^2+4X+1\). Its four roots lie respectively
in \((-2,-19/10)\), \((-1/4,-1/5)\),
\((13/10,7/5)\), and \((9/5,19/10)\).
Thus \(3-2t\) is negative only at the identity
embedding. Under \(\Phi_1\), every value of
\(i(3-2t)\) has positive imaginary part.
The element is integral and purely imaginary, so
\(\psi\) is integral and alternating; the sign
criterion of \cite[Example 2.9, pp.~25--26]{GenMilneCM2020}
gives \(\psi(x,Jx)>0\). This verifies the displayed
polarization with its actual sign.
The stabilizer of \(\Phi_1\) in \(C_4\times C_2\)
is trivial: a translation preserving its unique
negative choice fixes that choice, whereas
conjugation changes the count from one to three.
No simplicity assertion is needed.

Include a common finite field of definition, a finite
basis of labelled integral endomorphisms, and the
Artin module maps of E.2.3. One additional labelled
CM variety (E.2.20f) therefore obstructs extension
of E.2.19's specified comparison. This uses actual
objects and maps, without assigning an arbitrary
idele to a Galois element fixing \(\Q^{\mathrm{cm}}\).

Finally a global sign preserves every actual lattice.
Substitution into the proved necessity and sufficiency
of E.2.18 gives (E.2.20a), with exactly the two lifts
of E.2.15. E.2.21 further proves that every element
of \(\Sigma\) fixes \(F_{\mathrm{cm}}\).
Whether \(\Sigma\) contains a nonidentity
element, whether it connects the quartic place-pair,
and whether an abelian comparison extends to the full
pure diagram remain the precise unresolved questions.
\end{proof}

\begin{result}{E.2.21}{The exact Artin image of the full plain integral CM diagram}{Proved}
Let \(I_{\mathrm{CM}}\) be the Artin image for the concrete
geometric integral diagram generated by all CM abelian varieties
over \(\Qb\), their actual isogeny lattices, specified tensor
and summand lattices, labelled algebraic correspondences and
all finite-model Artin algebra/module maps. The same assertion
holds with the labelled absolute-Hodge tensors of the abelian
version. Then
\[
 I_{\mathrm{CM}}
   =\Gamma_{F_{\mathrm{cm}}}\sqcup
        c\Gamma_{F_{\mathrm{cm}}},
 \qquad
 F_{\mathrm{cm}}=\prod_{K\ {\rm CM}}K^{\mathrm{ab}},
 \tag{E.2.21a}
\]
where the product denotes the compositum in \(\Qb\).
Its Artin kernel consists of the identity and the global
weight sign; each displayed fibre has exactly two elements.
In particular every \(\sigma\in\Sigma\) fixes the full
abelian extension of every imaginary-quadratic field,
including all its ray class fields.

The all-abelian criterion is consequently
\[
 \begin{split}
 \Sigma_F
  &=\{\sigma\in\Gamma_{F_{\mathrm{cm}}}:
       \rho_A(\sigma)(\Lambda_A)=\Lambda_{\sigma A}
       \text{ for every abelian variety }A/\Qb\},\\
 I_{\mathrm{ab}}&=\Sigma_F\sqcup c\Sigma_F .
 \end{split}
 \tag{E.2.21b}
\]
This identifies the CM restriction of every remaining
abelian Artin candidate without identifying a field
generated by CM division points with \(K^{\mathrm{ab}}\).
\end{result}
\begin{proof}
The CM tests used in E.2.5 and E.2.15 already belong
to the present CM/Artin diagram. They force the
restriction of an Artin candidate to
\(\Q^{\mathrm{cm}}\) to be the identity or actual
conjugation. Compose with conjugation and assume
\(\sigma\in\Gamma_{\Q^{\mathrm{cm}}}\).
The all-cyclic-quotient argument of E.2.9 also applies
inside this diagram, because a quotient of a CM
abelian variety is CM. Products give one common
scalar, and the Gaussian anchor makes its square
one. The two-type equations of E.2.20 are therefore
available for every Galois CM field \(F(i)\).

First, the imaginary-quadratic conclusion follows
directly from E.2.20. For a finite \(K\)-idele \(s\)
representing \(\sigma|_{K^{\mathrm{ab}}}\), the
identity CM type gives
\(\lambda s\in K^\times\). Since \(\lambda\) is now
one global sign, \(s\) is principal. Its finite
Artin image is trivial: the missing infinite
component belongs to \(\C^\times\), whose image
is trivial by connectedness. This proves fixation
of \(K^{\mathrm{ab}}\) through the actual reciprocity
map.

For the general assertion, let \(F/\Q\) be finite,
totally real and Galois, and put \(L=F(i)\).
For a finite \(L\)-idele \(s\) representing
\(\sigma|_{L^{\mathrm{ab}}}\), the quotient of the
two CM-type equations proves
\[
 s=b\,t,\qquad b\in L^\times,\quad
 t\in\mathbb A_{F,f}^{\times}.
 \tag{E.2.21c}
\]
This is the explicit Hilbert-90 calculation in
E.2.20. It cancels the common scalar, so this
part of the proof does not require its globality.
The principal factor has trivial finite Artin
image because \(L\) is CM.

Fix an arbitrary finite abelian extension \(M/L\).
Define the continuous homomorphism
\[
 v_M:\Gamma_F^{\mathrm{ab}}
   \xrightarrow{\ \operatorname{Ver}_{F\to L}\ }
   \Gamma_L^{\mathrm{ab}}
   \xrightarrow{\ \operatorname{res}_M\ }
   \Gal(M/L).
 \tag{E.2.21d}
\]
The first arrow is the actual transfer for the
index-two subgroup \(\Gamma_L\subset\Gamma_F\).
It is related to extension \(j_F^L\) of finite
ideles by
\[
 \operatorname{art}_L(j_F^L(u))
  =\operatorname{Ver}_{F\to L}
       \bigl(\operatorname{art}_F(u)\bigr).
 \tag{E.2.21e}
\]
The finite Galois transfer and Artin square are
\cite[Chapter V, Proposition 3.18 and adjoining
diagram, p.~162]{GenMilneCFT2020}.
For a finite target, choose an ideal representative
prime to its modulus to pass from that square to
finite ideles; weak approximation and the
principal-idele relation justify the replacement.
Taking inverse limits gives (E.2.21e).

For each real embedding \(\tau:F\hookrightarrow\R\),
let \(c_\tau\in\Gamma_F\) be its actual complex
conjugation. It has order two and lies outside
\(\Gamma_L\), since \(i\) is nonreal at every real
place. Using the two coset representatives
\(1,c_\tau\) gives the exact formula
\[
 \operatorname{Ver}_{F\to L}(c_\tau)=c_\tau^2=1.
 \tag{E.2.21f}
\]
No Galois assumption on \(M/F\), or inversion
assumption on \(\Gal(M/L)\), enters this calculation.

The finite quotient \(v_M\) defines an abelian
extension \(N/F\): its fixing subgroup in
\(\Gamma_F\) is the inverse image of
\(\ker(v_M)\). Equation (E.2.21f) kills the
complex conjugation at every real place.
Consequently every embedding of \(N\) is real,
so \(N\) is totally real.

Choose \(F'/\Q\) to be the normal closure of
\(N/\Q\). It is totally real and contains \(F\).
Moreover \(F'/F\) is abelian: since \(F/\Q\)
is Galois, each conjugate of \(N\) is an abelian
extension of \(F\), and their compositum remains
abelian over \(F\). Put \(L'=F'(i)\).
Then \(L'/\Q\) is Galois and CM, and \(L'/L\)
is abelian. Apply (E.2.21c) at this larger
field to obtain
\[
 s'=b't',\qquad
 \operatorname{art}_{L'}(s')
      =\sigma|_{L'^{\mathrm{ab}}},
 \quad b'\in L'^\times,\quad
 t'\in\mathbb A_{F',f}^{\times}.
\]

The norm/restriction Artin diagram gives
\[
 \begin{split}
 \sigma|_{L^{\mathrm{ab}}}
  &=\operatorname{art}_L
       \bigl(\operatorname{Nm}_{L'/L}(s')\bigr)\\
  &=\operatorname{art}_L
       \bigl(j_F^L(\operatorname{Nm}_{F'/F}(t'))\bigr).
 \end{split}
 \tag{E.2.21g}
\]
For the first equality, the actual field map is
\(L^{\mathrm{ab}}L'\subset L'^{\mathrm{ab}}\),
obtained by abelian base change.
For the second, the principal norm of \(b'\)
has trivial Artin image in the totally imaginary
field \(L\). The remaining norm identity follows
from
\[
 F'\otimes_F L=L',
\]
because \(L=F(i)\) and \(F'\) is totally real:
the determinant defining the norm commutes with
this extension of scalars.

Set
\[
 a=\operatorname{art}_F
       \bigl(\operatorname{Nm}_{F'/F}(t')\bigr).
\]
The same norm/restriction diagram makes \(a\)
trivial on \(N\), since \(N\subset F'\) and
\(\operatorname{art}_{F'}(t')\) fixes \(F'\).
Thus \(a\in\ker(v_M)\). Combining (E.2.21d),
(E.2.21e), and (E.2.21g) yields
\[
 \sigma|_M=v_M(a)=1.
\]
As \(M/L\) was arbitrary finite abelian, \(\sigma\)
fixes \(L^{\mathrm{ab}}\) pointwise.

Every finite CM field \(K\) is contained in
a Galois CM field \(L\) containing \(i\), obtained
by adjoining \(i\) to its normal closure.
For \(F=L^+\), one has \(L=F(i)\), with \(F/\Q\)
Galois and totally real. The preceding argument
applies. Since
\[
 K^{\mathrm{ab}}L\subset L^{\mathrm{ab}},
\]
\(\sigma\) fixes \(K^{\mathrm{ab}}\).
Varying \(K\) proves
\(\sigma\in\Gamma_{F_{\mathrm{cm}}}\).

Conversely E.2.6 constructs an actual comparison
for every element of \(\Gamma_{F_{\mathrm{cm}}}\):
the main CM theorem with idele \(1\), the
Tate-integral quasi-isogeny argument, exterior
cohomology, all labelled correspondences and
projector lattices, and the finite-model
Artin descent give every required map.
Complex conjugation supplies the other coset.
This proves (E.2.21a).
The proof of the Artin kernel in E.2.15 applies
verbatim to CM quotients and products, giving
the two global weight choices and exactly
two lifts in each fibre.

Restriction from the full geometric abelian
diagram gives the necessary field constraint
in (E.2.21b), and E.2.16 proves its exact
lattice necessity and sufficiency.
The known quartic construction already has
swapping elements in \(\Gamma_{F_{\mathrm{cm}}}\).
Thus the remaining all-abelian lattice condition,
the quartic place-pair, and extension to all
pure motives remain unresolved.
\end{proof}

\begin{result}{E.2.22}{An exact reduction to one tower of covers}{Proved}
Put \(X=\mathbf P^1_{\Q}\setminus\{0,1,\infty\}\), and retain
the actual Betti lattices and comparison maps of E.2.21.
For \(\sigma\in\Gamma_{F_{\mathrm{cm}}}\), the following
conditions are equivalent:
\begin{enumerate}
\item \(\rho_A(\sigma)\Lambda_A=\Lambda_{{}^{\sigma}A}\)
for every abelian variety \(A/\Qb\).
\item For every connected finite \(\acute{e}\)tale cover
\(Y\to X_{\Qb}\), with smooth projective compactification
\(C_Y\), the actual Galois comparison on the Jacobian satisfies
\[
 \rho_{J(C_Y)}(\sigma)\Lambda_{J(C_Y)}
       =\Lambda_{J({}^{\sigma}C_Y)}.
 \tag{E.2.22a}
\]
\end{enumerate}
There is a precise topological formulation of condition 2.
Fix \(x=2\), let \(\Pi=\pi_1(X(\C),x)\), and identify
\(\widehat\Pi\) with \(\pi_1^{\mathrm{et}}(X_{\Qb},x)\).
The rational point gives the Galois action
\(\alpha_\sigma\in\Aut(\widehat\Pi)\).
For every open subgroup \(U\subset\widehat\Pi\), set
\[
 \Pi_U=\Pi\cap U,\qquad
 \Lambda_U=H_1(C_U(\C),\Z),
 \tag{E.2.22b}
\]
where \(C_U\) compactifies the pointed cover associated to \(U\).
Then condition 2 says exactly that the induced map
\[
 \bar\alpha_{\sigma,U}:
 \Lambda_U\otimes\widehat\Z
       \longrightarrow
 \Lambda_{\alpha_\sigma(U)}\otimes\widehat\Z
 \tag{E.2.22c}
\]
carries \(\Lambda_U\) onto \(\Lambda_{\alpha_\sigma(U)}\)
for every \(U\). This property depends only on the outer
Galois action. Consequently these conditions describe exactly
the subgroup \(\Sigma_F\) in E.2.21.
\end{result}
\begin{proof}
We first supply the integral quotient statement needed for the
converse. Every nonzero abelian variety \(A/\Qb\) admits a
smooth projective curve \(C/\Qb\) and a homomorphism
\[
 q:J(C)\longrightarrow A,\qquad
 q_*H_1(J(C)(\C),\Z)=H_1(A(\C),\Z).
 \tag{E.2.22d}
\]
For dimension one take \(C=A\). In higher dimension take the
successive smooth ample sections in the proof of
\cite[Theorem 10.1, Lemma 10.3, pp.33--35]{GenMilneJVs2021}.
For every finite surjective \(\acute{e}\)tale isogeny
\(A'\to A\), their successive inverse images are smooth and
connected: pullback preserves ampleness, an ample effective
divisor on a projective variety of dimension at least two is
connected, and smoothness is preserved by an \(\acute{e}\)tale
base change. Thus the resulting curve has connected inverse
image under every such isogeny. Its inclusion induces the
surjective \(q\) of that theorem.

Here is why this particular construction gives the asserted
integral surjectivity. Write
\(A(\C)=V/\Lambda_A\), and put
\(\Lambda'=q_*H_1(J(C)(\C),\Z)\).
Surjectivity of \(q\) makes \(\Lambda'\) a subgroup of finite
index in \(\Lambda_A\). Suppose that index exceeds one.
Choose \(n\geq1\) with \(n\Lambda_A\subset\Lambda'\).
The complex torus \(V/\Lambda'\), with its map to \(A\),
comes from an actual isogeny over \(\Qb\): form
\[
 H=(n^{-1}\Lambda')/\Lambda_A\subset A[n],
 \qquad A/H\longrightarrow A,
 \tag{E.2.22e}
\]
where the last map is induced by \([n]\).
Multiplication by \(n\) identifies its analytic source with
\(V/\Lambda'\). All points of \(A[n]\) and hence this finite
quotient are defined over \(\Qb\).
The image of \(\pi_1(C(\C))\) in
\(\pi_1(A(\C))=\Lambda_A\) is precisely \(\Lambda'\),
because the Abel--Jacobi map is an isomorphism on integral
first homology. The pullback of (E.2.22e) to \(C\) therefore
has trivial monodromy on \(\Lambda_A/\Lambda'\); it is a
disjoint union of more than one copy of \(C(\C)\).
This contradicts the connected inverse-image property.
It proves (E.2.22d). The same argument applies to
\({}^{\sigma}q\).

Belyi's construction, in the primary version
\cite[Theorem 1 and old proof, pp.1--2]{GenBelyi1997}, gives a
finite morphism \(C\to\mathbf P^1_{\Qb}\) whose branch
values lie in \(\{0,1,\infty\}\). Removing the inverse image
of these three points makes \(C\) the compactification of a
cover occurring in condition 2.
The relevant actual homology/Tate maps commute:
\[
\begin{array}{ccc}
 \Lambda_{J(C)}\otimes\widehat\Z
   &\xrightarrow{\rho_{J(C)}(\sigma)}&
 \Lambda_{J({}^{\sigma}C)}\otimes\widehat\Z\\
 \big\downarrow q_* && \big\downarrow({}^{\sigma}q)_*\\
 \Lambda_A\otimes\widehat\Z
   &\xrightarrow{\rho_A(\sigma)}&
 \Lambda_{{}^{\sigma}A}\otimes\widehat\Z.
\end{array}
 \tag{E.2.22f}
\]
In these completed modules, (E.2.22d) and condition 2 imply
the exact equality
\[
\begin{split}
 \rho_A(\sigma)\Lambda_A
  &=\rho_A(\sigma)q_*\Lambda_{J(C)}\\
  &=({}^{\sigma}q)_*\rho_{J(C)}(\sigma)\Lambda_{J(C)}\\
  &=({}^{\sigma}q)_*\Lambda_{J({}^{\sigma}C)}
    =\Lambda_{{}^{\sigma}A}.
\end{split}
 \tag{E.2.22g}
\]
The zero abelian variety contributes no constraint.
This proves 2 implies 1; 1 implies 2 by taking the indicated
Jacobians.

For completeness, the map in (E.2.22c) uses the compactified
homology quotient, not merely the homology of the affine
cover. The finite-index subgroup \(\Pi_U\) has profinite
completion \(U\). In its abelianization, quotient by the
classes of positively oriented loops around the punctures;
this gives \(\Lambda_U\). Equivalently, quotient
\(U^{\mathrm{ab}}\) by the closed subgroup generated by its
peripheral inertia images; the result is
\(\Lambda_U\otimes\widehat\Z\).
The Galois action carries those inertia groups to the inertia
groups of the conjugate cover, and hence induces
\(\bar\alpha_{\sigma,U}\).

The Abel--Jacobi map induces the abelianized fundamental-group
isomorphism of \cite[Proposition 9.1 and Corollary 9.6,
pp.30--32]{GenMilneJVs2021}. Together with finite-coefficient
Betti--\(\acute{e}\)tale comparison
\cite[Expos\'e XVI, Th\'eor\`eme 4.1]{GenSGA4}, it gives
the commutative diagram
\[
\begin{array}{ccc}
 \Lambda_U\otimes\widehat\Z
   &\xrightarrow{\bar\alpha_{\sigma,U}}&
 \Lambda_{\alpha_\sigma(U)}\otimes\widehat\Z\\
 \big\downarrow\simeq && \big\downarrow\simeq\\
 T J(C_U)&\xrightarrow{\rho_{J(C_U)}(\sigma)}&
 T J(C_{\alpha_\sigma(U)}).
\end{array}
 \tag{E.2.22h}
\]
Here \(TJ=\prod_\ell T_\ell J\). On the Jacobian the
comparison sends a lattice vector \(v\bmod n\) to
\(v/n\bmod\Lambda_J\); on the curve it is induced by
Abel--Jacobi. Changing its base point is a translation on
the Jacobian and acts trivially on first homology.
Thus (E.2.22h) identifies the actual integral lattices as
well as their completions. Algebraic cover maps, their
Jacobian maps and deck actions commute with this diagram;
no independent choices of transport are introduced.

Finally, inner automorphisms do not change this lattice
condition. Given \(h\in\widehat\Pi\) and open \(U\),
density supplies \(p\in\Pi\) with \(h\in pU\).
Writing \(h=pu\), \(u\in U\), gives
\(hUh^{-1}=pUp^{-1}\), and conjugation by \(h\) on
\(U^{\mathrm{ab}}\) equals conjugation by \(p\), because
conjugation by \(u\) is inner. The latter is induced by an
actual topological change of base path, so it carries the
compactified integral homology lattice onto the corresponding
lattice. This proves the asserted outer independence.
\end{proof}

\paragraph{The remaining map.}
E.2.22 is a proved equivalence of the existing all-abelian
lattice condition with the stated condition on one cover
tower. It does not establish that its outer automorphisms lie
in the image of a profinite mapping-class group.
The opened primary theorem
\cite[Theorem 1.1, p.2, proof p.3]{GenMatsumotoTamagawa1997}
identifies the intersection of the Galois image and the
orientation-preserving profinite mapping-class image as the
identity for an affine hyperbolic curve with rational
punctures. The map from the integral condition (E.2.22c) to
that mapping-class image remains unconstructed here.
Completed cup or Reidemeister pairing preservation alone
does not supply it: the actual \(\acute{e}\)tale Galois maps
already preserve those pairings up to the cyclotomic
character, which is one on \(\Gamma_{F_{\mathrm{cm}}}\).
The quartic place-pair and the full pure-motive problem
therefore remain open.

\begin{result}{E.2.23}{One integral comparison for all rational elliptic isogeny classes}{Proved}
Let \(\mathcal A_{\mathrm{ratell}}\) be the class of abelian
varieties over \(\Qb\) that are geometrically isogenous to a
product of non-CM elliptic curves having models over \(\Q\)
and an arbitrary CM abelian variety. Include every actual
isogeny-class lattice, all products, all finite models and their
underlying-\(\Q\) Artin module maps, all Artin objects, and
the tensor, correspondence and specified projector constructions
of E.2.16.

Let \(P\) be the splitting field of \(x^4-x-1\), with roots
labelled \(\alpha_1>0\), \(\alpha_2<0\),
\(\operatorname{Im}\alpha_3>0\),
\(\operatorname{Im}\alpha_4<0\). There is a specified element
\(\sigma\in\Gamma_{F_{\mathrm{cm}}}\) such that
\[
 \sigma|_P=(12)(34),\qquad
 \rho_A(\sigma)\Lambda_A=\Lambda_{{}^{\sigma}A}
       \quad(A\in\mathcal A_{\mathrm{ratell}}).
 \tag{E.2.23a}
\]
The actual completed transports therefore restrict to a single
natural integral tensor comparison on the full stated diagram.
For \(k=\Q(\alpha_1)\), they give a comparison between its
two real embeddings sending \(\alpha_1\) to
\(\alpha_1\) and \(\alpha_2\), on this entire diagram.
The assertion is about the plain integral diagram; the
place-labelled real Frobenius of E.2.7 is not imposed.
\end{result}
\begin{proof}
First we prove the finite-family image statement that makes
the simultaneous construction possible. Let
\(E_1,\ldots,E_r\) be non-CM elliptic curves over \(\Q\)
that are pairwise nonisogenous over \(\Qb\), and write
\[
 \rho:\Gamma_{\Q}\longrightarrow
 \Delta=\{(g_i)\in\GL_2(\widehat\Z)^r:
                  \det g_i=\det g_j\},\qquad
 S=\operatorname{SL}_2(\widehat\Z)^r.
\]
The product open-image theorem
\cite[Theorem 1.1 and Remark 1.2, author PDF p.1]{GenLombardo}
makes \(G=\rho(\Gamma_{\Q})\) open in \(\Delta\).
For one factor use \cite[Theorem 3, p.299]{GenSerre1972}.
Thus \(G\) contains a congruence product
\[
 K(M)=\prod_{i,\ell}B_\ell(v_\ell(M)),\qquad
 B_\ell(a)=\ker\bigl(\operatorname{SL}_2(\Z_\ell)
                 \to\operatorname{SL}_2(\Z/\ell^a\Z)\bigr),
 \tag{E.2.23b}
\]
where \(B_\ell(0)=\operatorname{SL}_2(\Z_\ell)\).
In particular its elements may be independently supported
at one prime and one elliptic factor.

Let \(c\) be actual complex conjugation and \(C=\rho(c)\).
It is central on \(\Q^{\mathrm{cm}}\), so
\[
 [\gamma,c]\in\Gamma_{\Q^{\mathrm{cm}}}
       \quad(\gamma\in\Gamma_\Q).
 \tag{E.2.23c}
\]
Each component of \(C\) is an involution of determinant
\(-1\). At an odd prime it is conjugate to
\(\operatorname{diag}(1,-1)\) over \(\GL_2(\Z_\ell)\),
by the two integral projectors \((1\pm C)/2\).
At \(2\) it is conjugate over \(\GL_2(\Q_2)\);
sufficiently deep root subgroups absorb the denominators.
Put
\[
 U(t)=\begin{pmatrix}1&t\\0&1\end{pmatrix},\qquad
 L(t)=\begin{pmatrix}1&0\\t&1\end{pmatrix}.
\]
For the commutator convention \([a,b]=aba^{-1}b^{-1}\),
\[
 [U(t),\operatorname{diag}(1,-1)]=U(2t),\qquad
 [L(t),\operatorname{diag}(1,-1)]=L(2t).
 \tag{E.2.23d}
\]
Equations (E.2.23b)--(E.2.23d) put independently supported
upper and lower roots of some depth in
\(N=\rho(\Gamma_{\Q^{\mathrm{cm}}})\).
For \(\ell\nmid2M\) they put both full root groups there,
and hence the full supported \(\operatorname{SL}_2(\Z_\ell)\).

Here is an explicit argument at the remaining primes.
Roots with parameters in \(\ell^a\Z_\ell\), \(a\geq1\),
generate a group containing \(B_\ell(2a)\), since
\[
 U(x)L(y)U\!\left(\frac{-x}{1+xy}\right)
 L\!\left(-y(1+xy)\right)
       =\operatorname{diag}(1+xy,(1+xy)^{-1}).
 \tag{E.2.23e}
\]
Take \(x=\ell^a\) to obtain every diagonal with first
entry in \(1+\ell^{2a}\Z_\ell\), and then use
\[
 \begin{pmatrix}a_0&b_0\\c_0&d_0\end{pmatrix}
 =L(c_0/a_0)\operatorname{diag}(a_0,a_0^{-1})U(b_0/a_0)
 \quad(a_0\text{ a unit}).
 \tag{E.2.23f}
\]
After the possible conjugation at \(2\), this still contains
an ordinary principal congruence subgroup: if \(2^dP\)
and \(2^dP^{-1}\) are integral, then
\(P^{-1}(B_2(k+2d)-I)P\) has entries in \(2^k\Z_2\).
The compact image \(N\) is closed. Taking the closure of
the finite products of these independently supported groups
gives an open product subgroup of \(S\) contained in \(N\).
Also \(N\subset S\), since \(\Q^{\mathrm{cm}}\)
contains \(\Q^{\mathrm{ab}}\).

Let \(D\) be the closed commutator subgroup of
\(\Gamma_{\Q^{\mathrm{cm}}}\). For every CM field \(K\),
the extension \(K^{\mathrm{ab}}\Q^{\mathrm{cm}}/
\Q^{\mathrm{cm}}\) is abelian by base change. Therefore
\[
 D\subset\Gamma_{F_{\mathrm{cm}}}.
 \tag{E.2.23g}
\]
The compact surjection onto \(N\) maps \(D\) onto its
closed commutator subgroup \(N'\): products of commutators
map to the corresponding products, and compactness and
continuity give equality after closure.
For \(\ell\geq5\), a full supported
\(\operatorname{SL}_2(\Z_\ell)\) is perfect, as
\[
 [\operatorname{diag}(2,1/2),U(t)]=U(3t)
\]
and the lower-root identity give both full root groups.
At the finitely many remaining primes choose a supported
\(B_\ell(a)\subset N\), with \(a\geq1\), and
\(a\geq2\) when \(\ell=2\). Setting \(d=1+\ell^a\),
\[
 [\operatorname{diag}(d,d^{-1}),U(\ell^at)]
       =U\!\left(\ell^{2a}(2+\ell^a)t\right).
 \tag{E.2.23h}
\]
The analogous lower roots have the same depth. That depth
is \(2a\) at odd primes and \(2a+1\) at \(2\).
Equations (E.2.23e)--(E.2.23f) again give an open congruence
subgroup in each supported factor. Closure gives an actual
open product inside \(N'\). We have proved
\[
 \rho(\Gamma_{F_{\mathrm{cm}}})
       \text{ is open in }S.
 \tag{E.2.23i}
\]
Only the inclusion (E.2.23g), not an equality of Galois
commutator subgroups, was needed.

We next choose one element satisfying all integral constraints.
Enumerate pairs of integers \((A,B)\) by height and then
lexicographically, discard singular curves
\(y^2=x^3+Ax+B\) and CM curves, and retain the first model
in each geometric isogeny class. Denote this sequence by
\(E_1,E_2,\ldots\). It includes every required class.
Fix its integral Betti bases canonically as follows. For
\(\omega=dx/y\), take the smallest positive real period
\(\omega_1\), and the period \(\omega_2\) of smallest
positive imaginary part with
\(0\leq\operatorname{Re}\omega_2<\omega_1\).
The conjugation-stable lattice has real intersection of
rank one, and its quotient by that primitive subgroup is
infinite cyclic; these rules therefore specify a unique
oriented basis. Take the unique cycles with these periods.

By E.2.6, \(P\cap F_{\mathrm{cm}}=P^{V_4}\).
Thus the compact set
\[
 C_0=\{\sigma\in\Gamma_{F_{\mathrm{cm}}}:
                    \sigma|_P=(12)(34)\}
\]
is nonempty. Suppose that the nonempty compact set \(C_n\)
has already imposed
\(\rho_{E_i}(\sigma)=S_i\in\operatorname{SL}_2(\Z)\)
for \(1\leq i\leq n\). It is a coset of
\[
 H_n=\Gamma_{F_{\mathrm{cm}}}\cap\Gamma_P
                      \cap\bigcap_{i\leq n}\ker\rho_{E_i}.
 \tag{E.2.23j}
\]
By (E.2.23i), the image of
\(\Gamma_{F_{\mathrm{cm}}}\cap\Gamma_P\) on the first
\(n+1\) factors is open: it is a compact, finite-index
subgroup of the open image for \(\Gamma_{F_{\mathrm{cm}}}\).
It contains a product congruence subgroup. In particular the
image \(\rho_{E_{n+1}}(H_n)\) is open, since that product
contains every sufficiently deep last-factor element with
the first \(n\) components exactly the identity.
Hence \(\rho_{E_{n+1}}(C_n)\) is a nonempty open coset.

The subgroup \(\operatorname{SL}_2(\Z)\) is dense in
\(\operatorname{SL}_2(\widehat\Z)\). Explicitly,
\(\operatorname{SL}_2(\Z/m\Z)\) is generated by upper
and lower elementary matrices: over each local factor a
unimodular column has a unit entry, elementary operations
and a determinant-one interchange reduce it, and the
diagonal is also a product of elementary matrices.
Chinese remaindering gives the claim for \(m\), and the
elementary matrices lift integrally. Choose the first
\(S_{n+1}\in\operatorname{SL}_2(\Z)\), in entry-height
and lexicographic order, that belongs to this open coset.
Its inverse image \(C_{n+1}\subset C_n\) is nonempty
and compact. Compactness gives
\[
 \varnothing\ne C_\infty=\bigcap_{n\geq0}C_n.
 \tag{E.2.23k}
\]
Use the fixed finite Galois exhaustion and ordered finite
quotients of E.2.17; choose successively the first cylinder
meeting \(C_\infty\). The resulting extendible branch
specifies a unique \(\sigma\in C_\infty\).
This is a full mathematical specification; no computability
of its membership decisions is asserted.

It remains to check every geometric isogeny lattice, rather
than just the representatives. On a product \(B\) of these
representatives and a CM variety, let \(T_B\) be the direct
sum of the integral \(S_i\) and the integral CM comparison
with idele \(1\) from E.2.6. Its completion is exactly
\(\rho_B(\sigma)\). If \(q:B\to A\) is any geometric
isogeny, identify rational homology using \(q_*\) and put
\[
 L_A=q_*^{-1}\Lambda_A,
 \qquad \Lambda_B\subset L_A\subset m^{-1}\Lambda_B
 \tag{E.2.23l}
\]
for an integer \(m\) killing \(\ker q\).
The subgroup \(L_A/\Lambda_B\) is exactly \(\ker q\)
under the analytic torsion identification. Since the
completion of \(T_B\) is the actual Galois map, its
reduction modulo \(m\) carries this finite subgroup onto
\(\ker({}^{\sigma}q)\). Thus
\(T_B L_A=L_{{}^{\sigma}A}\), and the rational map
\[
 T_A=({}^{\sigma}q)_*\,T_B\,q_*^{-1}
       :H_1(A,\Q)\longrightarrow H_1({}^{\sigma}A,\Q)
 \tag{E.2.23m}
\]
restricts to an isomorphism of the actual integral lattices.
Its completion is \(\rho_A(\sigma)\). This forces
independence of the chosen isogeny, because an integral
homomorphism is determined by its completion.

Use inverse dual maps on cohomology and exterior powers in
all degrees. Every homomorphism, correspondence, tensor,
polarization and projector equation follows from its actual
\(\acute{e}\)tale equation and injectivity of the Betti
completion, as in E.2.16; the Tate object is fixed since
\(\sigma\in\Gamma_{F_{\mathrm{cm}}}\).
The class \(\mathcal A_{\mathrm{ratell}}\) is stable
under conjugation: its rational elliptic representatives
remain fixed and CM varieties remain CM. The finite-model
Artin construction of E.2.3 therefore applies to every
embedding block, with \(e_u\mapsto e_{\sigma u}\), and
respects all changes of finite model and all module maps.
It gives the promised single comparison on the entire
diagram and the stated comparison of the two \(k\)-places.
No assertion in this construction supplies the corresponding
integrality for arbitrary additional simple abelian factors
or an extension to every pure motive.
\end{proof}

\begin{result}{E.2.24}{One integral comparison for all elliptic and CM geometric factors}{Proved}
Let \(\mathcal A_{\mathrm{ell}}\) consist of all abelian
varieties over \(\Qb\) geometrically isogenous to products
of arbitrary elliptic curves and arbitrary CM abelian varieties.
Take their full integral diagram with every geometric isogeny
lattice, product, finite-model Artin module map, all Artin
objects, and the correspondence, tensor and specified projector
constructions of E.2.16.
With \(P\) and its root labels as in E.2.23, there is a
specified \(\sigma_*\in\Gamma_{F_{\mathrm{cm}}}\) with
\[
 \sigma_*|_P=(12)(34),\qquad
 \rho_A(\sigma_*)\Lambda_A=\Lambda_{{}^{\sigma_*}A}
            \quad(A\in\mathcal A_{\mathrm{ell}}).
 \tag{E.2.24a}
\]
These actual transports define one natural integral tensor
comparison on the entire stated diagram and, by E.2.3, a
comparison between the two specified real embeddings of
\(\Q(\alpha_1)\). The comparison concerns the plain integral
diagram; the labelled real Frobenius constraint of E.2.7 is
not imposed.
\end{result}
\begin{proof}
The additional issue relative to E.2.23 is that conjugation
may move an elliptic model and even its geometric isogeny
class. We first prove the precise finite-family openness
needed in this situation.

Start with any finite family of pairwise geometrically
nonisogenous non-CM elliptic curves over \(\Qb\).
Enlarge its classes to their finite \(\Gamma_\Q\)-orbits
and choose representatives \(E_i\) for the resulting finite
set. Complex conjugation permutes this set by an involution
\(\pi\). Choose actual isogenies
\(\phi_i:{}^cE_i\to E_{\pi(i)}\), and a finite Galois
field \(K/\Q\) defining all these models and isogenies.
The product representation on \(\Gamma_K\) has open image
\(G\) in the equal-determinant group \(\Delta\), by
\cite[Theorem 1.1, author PDF p.1]{GenLombardo} and the
one-factor theorem \cite[Theorem 3, p.299]{GenSerre1972}.
Set
\[
 N=\rho(\Gamma_K\cap\Gamma_{\Q^{\mathrm{cm}}}),
 \qquad S=\prod_i\operatorname{SL}_2(\widehat\Z).
\]
This is a closed normal subgroup of \(G\), contained in \(S\).

Write \(c_i\) for the actual antiholomorphic map on Betti
homology \(H_1(E_i)\to H_1({}^cE_i)\), and put
\(C_i=(\phi_i)_*c_i\). It is a rational linear isomorphism
between the indicated two homology spaces, and is integral
with nonzero determinant in their integral bases.
Naturality of the isogeny defined over \(K\) gives the exact
intertwining identity
\[
 C_i\rho_i(c\gamma c)C_i^{-1}
                 =\rho_{\pi(i)}(\gamma)
                       \quad(\gamma\in\Gamma_K).
 \tag{E.2.24b}
\]
Thus conjugation by \(c\) acts on \(G\) by
\(\theta(g)_i=C_i^{-1}g_{\pi(i)}C_i\).
Since \(K/\Q\) is Galois and \(c\) is central on
\(\Q^{\mathrm{cm}}\), the actual commutators satisfy
\[
 [\gamma,c]\in\Gamma_K\cap\Gamma_{\Q^{\mathrm{cm}}},
 \qquad g\theta(g)^{-1}\in N.
 \tag{E.2.24c}
\]
As in (E.2.23b), \(G\) contains independently supported
principal congruence subgroups at every factor and prime,
and full supported factors at all but finitely many primes.

For a two-element orbit \(\pi(i)=j\ne i\), choose such
an element \(g\) supported at \((i,\ell)\).
The element in (E.2.24c) has components \(g\) at \(i\)
and \(C_j^{-1}g^{-1}C_j\) at \(j\).
Its commutator with an independently supported \(h\) at
\((i,\ell)\) has sole component \([g,h]\) at \(i\).
Normality keeps this in \(N\). The supported-root and
derived-group identities (E.2.23e)--(E.2.23h) therefore
give an open supported subgroup at every prime, and the
full supported factor at all sufficiently large primes.
The same holds for \(j\).

For a fixed class \(\pi(i)=i\), complex uniformization
makes \(C_i\) an antilinear map \(z\mapsto a\bar z\).
Its square is positive scalar multiplication. Algebraically
that square is \(\phi_i{}^c\phi_i\in\End(E_i)=\Z\).
Comparing degrees gives
\[
 C_i^2=d_iI,\quad \operatorname{tr}C_i=0,\quad
 \det C_i=-d_i,\qquad d_i=\deg\phi_i>0.
 \tag{E.2.24d}
\]
Hence \(C_i\) is nonscalar over every \(\Q_\ell\), and
is a nonscalar invertible matrix modulo every prime
outside a finite set. In particular no local square root
of \(d_i\) is needed. Equation (E.2.24c) supplies the
supported elements \([g,C_i^{-1}]\) of \(N\).

We give both local steps explicitly. At an exceptional prime,
let \(C\in\GL_2(\Q_\ell)\) be nonscalar, and let \(R\)
be the supported subgroup of \(N\). It is normal under a
principal congruence group and contains \([g,C]\) for every
sufficiently deep supported \(g\). One of
\(X=E_{12},E_{21}\) has
\(Z=X-\operatorname{Ad}_C X\ne0\), since a matrix
commuting with both is scalar. The polynomial matrix curve
\[
 f(t)=[I+tX,C]
\]
lies in \(R\) for sufficiently small \(t\), and has
derivative \(Z\) at zero. The \(\Q_\ell\)-span of
\(\operatorname{Ad}_hZ\), with \(h\) in the normalizing
principal congruence group, is all of
\(\mathfrak{sl}_2(\Q_\ell)\). Indeed conjugation by
\(I+tE_{12}\) or \(I+tE_{21}\) is polynomial of degree
at most two in \(t\). Interpolation at three sufficiently
small distinct values makes this span stable under brackets
with both root matrices. It is a nonzero ideal, and
\([H,E_{12}]=2E_{12}\), \([H,E_{21}]=-2E_{21}\),
\([E_{12},E_{21}]=H\) show that the characteristic-zero
Lie algebra has no such proper ideal.

Choose \(h_1,h_2,h_3\) giving independent derivatives.
The polynomial map
\[
 F(t_1,t_2,t_3)=\prod_{j=1}^3h_jf(t_j)h_j^{-1}
 \tag{E.2.24e}
\]
takes a neighborhood of zero into \(R\) and has invertible
derivative. To justify the open-image step directly, use
coordinates \((a-1,b,c)\) on
\(\left(\begin{smallmatrix}a&b\\c&d\end{smallmatrix}\right)\)
near the identity, where \(d=(1+bc)/a\).
Write the coordinate map as \(Jx+P(x)\), with \(J\)
invertible and \(P\) having no constant or linear term.
On a sufficiently small complete \(\ell\)-adic ball,
\(J^{-1}P\) has Lipschitz constant less than one and
maps the ball into a smaller ball. For every \(y\) in
that ball, iteration of
\(x\mapsto y-J^{-1}P(x)\) converges and solves
\(F(x)=Jy\) in coordinates. Thus the actual image of
(E.2.24e) contains a neighborhood of the identity.
This proves openness at every exceptional prime, including
\(2\) and the nonsplit cases of (E.2.24d).

At every remaining \(\ell\geq5\), \(R\) is normal
under the full supported \(\operatorname{SL}_2(\Z_\ell)\).
Some \([g,C]\) is noncentral modulo \(\ell\).
Otherwise conjugation by \(C\) would fix each upper and
lower unipotent up to a sign. Their trace is two, which
excludes the negative sign; commuting with both roots
would make \(C\) scalar.
The required finite-group argument can be given directly.
The root commutator with \(\operatorname{diag}(2,1/2)\)
and its lower-root counterpart make
\(\operatorname{SL}_2(\mathbf F_\ell)\) perfect.
Its quotient \(\operatorname{PSL}_2(\mathbf F_\ell)\)
acts faithfully and doubly transitively on
\(\mathbf P^1(\mathbf F_\ell)\). A point stabilizer
has the normal abelian translation subgroup, whose
conjugates generate the group. Any nontrivial normal
subgroup is transitive, since its orbits are an invariant
partition for a primitive faithful action. Modulo that
subgroup every conjugate of the translation subgroup has
the same image: write an element as a product of a normal
subgroup element and a point stabilizer element.
The quotient is therefore abelian and hence trivial by
perfectness. This proves the required simplicity.
It follows that the normal reduction of \(R\) surjects
onto \(\operatorname{PSL}_2\); its remaining quotient
has order at most two and is again trivial by perfectness.
Thus \(R\) has full reduction modulo \(\ell\).

Here is the full lifting calculation. A closed subgroup
\(R\subset\operatorname{SL}_2(\Z_\ell)\) with full
reduction contains a \(g\equiv U(1)\pmod\ell\).
Write \(g=U(1)(I+\ell B)\pmod{\ell^2}\).
The first-order error in \(g^\ell\) is \(\ell\)
times the sum of \(\operatorname{Ad}_{U(-j)}B\) for
\(0\leq j<\ell\). Its entries are polynomials in
\(j\) of degree at most two, and the sums of
\(1,j,j^2\) vanish modulo \(\ell\). Hence
\[
 g^\ell\equiv I+\ell E_{12}\pmod{\ell^2}.
 \tag{E.2.24f}
\]
Conjugating by lifts in \(R\) of a lower root and an
interchange matrix gives the span of
\(E_{12},E_{21},H\) in the first congruence quotient.
For \(n\geq1\),
\[
 (I+\ell^nX)^\ell
       \equiv I+\ell^{n+1}X\pmod{\ell^{n+2}}
\]
then gives every subsequent layer. Successive approximation
and closedness prove \(R=\operatorname{SL}_2(\Z_\ell)\).

The independently supported groups from the swapped and
fixed classes therefore give an open product inside \(N\).
The closed derived subgroup of
\(\Gamma_K\cap\Gamma_{\Q^{\mathrm{cm}}}\) lies in
\(\Gamma_K\cap\Gamma_{F_{\mathrm{cm}}}\), by the
abelian base-change argument of (E.2.23g), and maps onto
the closed derived subgroup of \(N\).
The same supported perfect factors and small-prime
commutator identities used in E.2.23 make that image open.
Projection to the original finite family proves
\[
 \rho_{E_1,\ldots,E_r}
       (\Gamma_K\cap\Gamma_{F_{\mathrm{cm}}})
   \text{ is open in }
       \operatorname{SL}_2(\widehat\Z)^r.
 \tag{E.2.24g}
\]
Further finite enlargement of \(K\) preserves this:
the new compact image has finite index in the old one,
and is therefore open.

We now perform the countable construction while fixing
the changing targets. Enumerate one short Weierstrass
model \(E_n/\Qb\) from each non-CM geometric isogeny
class, using the fixed algebraic-number ordering from
E.2.17 and retaining the first model in each class.
Fix an oriented Betti basis for each model and each
conjugate by its period lattice for \(dx/y\): among all
oriented lattice bases \((\omega_1,\omega_2)\), minimize
\(|\omega_1|^2+|\omega_2|^2\), then the lexicographic
tuple of their four real coordinates. Each bounded
region contains finitely many lattice pairs, so this
specifies a unique basis. The orientation is
\(\operatorname{Im}(\bar\omega_1\omega_2)>0\).

Start with the nonempty compact quartic coset \(C_0\)
from E.2.23. Choose increasing finite Galois fields
\(K_n\) containing \(P\), the first \(n\) models,
the finite orbit enlargements and isogenies used in
(E.2.24g), and the \(n\)-th member of the fixed finite
Galois exhaustion of \(\Qb\).
These choices can be fixed by the first appropriate
field and isogeny in the countable algebraic-data
enumerations. Partition the preceding compact fibre
by restriction to \(K_n\), and select the first
nonempty finite cylinder \(C'_n\).
The models \({}^{\sigma}E_i\) and all target bases
are now fixed for \(\sigma\in C'_n\).

For any \(\sigma_n\in C'_n\), its stabilizer under
right multiplication is exactly
\[
 H_n=\Gamma_{F_{\mathrm{cm}}}\cap\Gamma_{K_n}
                  \cap\bigcap_{i<n}\ker\rho_{E_i}.
 \tag{E.2.24h}
\]
Indeed, for \(k\in\Gamma_{K_n}\), actual transport
satisfies
\(\rho_{E_i}(\sigma_nk)=
\rho_{E_i}(\sigma_n)\rho_{E_i}(k)\).
Thus equality of all previous matrices is precisely
the kernel condition in (E.2.24h).
Openness (E.2.24g) makes \(\rho_{E_n}(H_n)\) open,
because its product congruence neighborhood contains
each sufficiently deep last-factor element with the
previous components exactly the identity.
The possible new matrices form the open coset
\[
 \rho_{E_n}(\sigma_n)\rho_{E_n}(H_n)
                \subset\operatorname{SL}_2(\widehat\Z).
 \tag{E.2.24i}
\]
The determinant is one since \(\sigma_n\) fixes all
roots of unity and the bases are oriented.
By the proved elementary density in E.2.23 this coset
contains an integral \(\operatorname{SL}_2(\Z)\)
matrix. Choose the first in entry-height and
lexicographic order, and let \(C_n\subset C'_n\)
be its inverse image. These are nested nonempty
compact sets. Their intersection contains an element
\(\sigma_*\), and contains only one: the selected
\(K_n\)-restrictions exhaust \(\Qb\).
It retains the prescribed quartic action and has
integral actual transport on every \(E_n\).
This specifies an element without claiming an
effective decision procedure for the choices.

On products of these elliptic representatives and
CM factors, take the direct sum of their actual
integral transports and the idele-one CM comparison.
For every geometric isogeny \(q:B\to A\), the exact
superlattice (E.2.23l) and its finite-kernel
interpretation show that
\(({}^{\sigma_*}q)_*T_Bq_*^{-1}\) carries
\(\Lambda_A\) onto \(\Lambda_{{}^{\sigma_*}A}\).
Completion injectivity gives independence of \(q\)
and every correspondence, tensor and projector
equation. Inverse duals and exterior powers give
cohomology, and the Tate object is fixed.
All finite-model Artin blocks and their module maps
are compatible by E.2.3 and E.2.16. This proves
(E.2.24a) on the entire specified diagram.

The theorem supplies no integrality for additional
simple non-CM abelian factors of dimension at least
two and no extension to every pure motive.
Those are still part of the full original
place-separation problem.
\end{proof}

\begin{result}{E.2.25}{One integral comparison for all surface and CM geometric factors}{Proved}
Let \(\mathcal A_{\leq2,\mathrm{CM}}\) consist of all
abelian varieties over \(\Qb\) geometrically isogenous to
products of simple factors of dimension at most two and
arbitrary CM abelian varieties. Take the full integral
diagram of E.2.16, including every geometric isogeny
lattice, all products, algebraic correspondences, tensor
constructions and specified projectors, finite-model
Artin module maps, and all Artin objects.
With \(P\) and its root labels as in E.2.23, there is a
specified \(\sigma_*\in\Gamma_{F_{\mathrm{cm}}}\) such that
\[
 \sigma_*|_P=(12)(34),\qquad
 \rho_A(\sigma_*)\Lambda_A
       =\Lambda_{{}^{\sigma_*}A}
       \quad(A\in\mathcal A_{\leq2,\mathrm{CM}}).
 \tag{E.2.25a}
\]
These actual completed transports give one integral
tensor comparison on the entire stated diagram, with
Tate map \(1\). By E.2.3 it compares the two specified
real embeddings of \(\Q(\alpha_1)\). This is the plain
integral diagram: E.2.7 supplies the explicit obstruction
when labelled real Frobenius is additionally required.
\end{result}
\begin{proof}
We prove the exact finite-family openness and integral
coset intersection before performing the simultaneous
construction.

\paragraph{1. The actual algebraic groups.}
Let \(A_i/\Qb\) be a finite family of pairwise
geometrically nonisogenous, absolutely simple, non-CM
abelian varieties of dimension one or two. Write
\[
 V_i=H_1(A_i(\C),\Q),\qquad
 \Lambda_i=H_1(A_i(\C),\Z),
\]
and choose labelled algebraic polarizations \(\psi_i\).
Enlarge a finite field of definition to define all
geometric endomorphisms and the finite set of labelled
Hodge tensors needed below.

The simple-surface classification used here is explicit
in the proofs of Lombardo's Corollaries 4.4--4.5,
author p.~12: there is no absolutely simple surface of
type III in characteristic zero, and the endomorphism
algebra of an absolutely simple surface cannot be an
imaginary quadratic field. Thus a non-CM simple surface
has geometric endomorphism algebra \(\Q\), a real
quadratic field \(F\), or an indefinite quaternion
division algebra \(B\). The remaining simple-surface
case is CM and already has the comparison of E.2.6
\cite[Corollaries 4.4--4.5 and their proofs, p.~12]
{GenLombardoNonsimple2015}.

The derived and full Mumford--Tate groups, in their
actual actions on \(V_i\), are as follows:
\[
\begin{array}{c|c|c}
 \End^0(A_i)&H_i&M_i\\ \hline
 \Q&\operatorname{Sp}(V_i,\psi_i)
       &\operatorname{GSp}(V_i,\psi_i)\\
 F&\operatorname{Res}_{F/\Q}\operatorname{SL}_{F}(V_i)
       &\{g\in\GL_F(V_i):\det_F(g)\in\Q^\times\}\\
 B&\operatorname{SL}_1(B^{\mathrm{op}})
       &(B^{\mathrm{op}})^\times .
\end{array}
\tag{E.2.25b}
\]
For an elliptic non-CM factor the first row is
\(\operatorname{SL}_2\subset\GL_2\).

Here is the verification of the actions and multipliers
in (E.2.25b). In the real multiplication case \(V_i\)
has dimension two over \(F\), and Rosati fixes \(F\)
pointwise. Indeed the nontrivial involution of a real
quadratic field sends a nonzero trace-zero element
\(a\) to \(-a\), contradicting positivity of
\(\operatorname{Tr}(aa^*)\). The trace pairing defines
a unique \(F\)-valued form \(\Psi_i\) by
\[
 \operatorname{Tr}_{F/\Q}\bigl(a\Psi_i(v,w)\bigr)
                  =\psi_i(av,w)\qquad(a\in F).
\]
Self-adjointness of \(F\) gives \(F\)-bilinearity,
alternation and nondegeneracy. In dimension two its
similitude is its determinant. This proves the displayed
rational-multiplier condition.

In the quaternionic case \(V_i\) is a free left
\(B\)-module of rank one; its centralizer is right
multiplication by \(B^{\mathrm{op}}\). Lombardo's
Lemmas 5.3--5.4 explicitly identify these right Galois
and left endomorphism actions at a sufficiently large
split prime, and Lemma 5.6 identifies the right
determinant with the cyclotomic multiplier
\cite[Lemmas 5.3--5.6, pp.~19--20]
{GenLombardoSurfaces2015}. Faithful extension
\(\Q\to\Q_\ell\) then proves over \(\Q\) that the
multiplier of right multiplication by \(b\) is
\(\operatorname{Nrd}(b)\).

To verify that the groups are the full groups in
(E.2.25b), use the Mumford--Tate equality for each
factor from Lombardo's Corollary 4.5, and the respective
full large-prime images from Theorem 1.3, Corollary 1.5
and Theorem 1.7 of his surface paper
\cite[Theorem 1.3, p.~2; Corollary 1.5 and
Theorem 1.7, p.~3]{GenLombardoSurfaces2015}.
The Mumford--Tate group lies in the displayed
endomorphism/polarization stabilizer because these
are labelled Hodge tensors. At one sufficiently large
prime this inclusion is an equality by those two
source statements. Faithful scalar extension gives
the equality over \(\Q\).

For the product \(A=\prod_i A_i\), Lombardo's
Corollary 4.4 gives, in the actual block representation,
\[
 H(A)=H:=\prod_i H_i,\qquad
 M:=\operatorname{MT}(A)=\Gm H,
 \qquad M^{\mathrm{der}}=H.
\]
Its factor projections are the given maps on homology;
the cited theorem proves that their inclusion into
the product is onto. The passage from the stated
\(\ell\)-adic equality to the Hodge equality uses
the same corollary's Mumford--Tate assertion
\cite[Corollary 4.4, p.~12]{GenLombardoNonsimple2015}.

Each \(H_i\) is simply connected and semisimple and
contains \(-I\) in this representation. The map
\((t,h)\mapsto th\) identifies \(M\) with
\((\Gm\times H)/\mu_2\), where the kernel is the
diagonal pair \((-1,-I)\). Its common multiplier
\(\nu(th)=t^2\) therefore gives the exact sequence
of algebraic groups
\[
 1\longrightarrow H\longrightarrow M
       \xrightarrow{\nu}\Gm\longrightarrow1.
 \tag{E.2.25c}
\]
In particular its multiplier-one kernel is precisely
\(H\), with its scheme structure.

The algebraic fundamental group \(\pi_1(M)\) is
\(\Z\) through \(\nu\): the coroot lattice is the
full cocharacter lattice of a maximal torus of the
simply connected kernel \(H\), and the quotient
cocharacter lattice is that of \(\Gm\). The Hodge
cocharacter has multiplier \(z\) or \(z^{-1}\),
according to the homology convention, so its class
generates this \(\Z\). This proves Hodge maximality
by Cadoret--Moonen's precise criterion
\cite[Sections 2.2--2.3, author p.~4]
{GenCadoretMoonen2020}. Their Theorem A(ii), applied
to the established Mumford--Tate equality of this
finite product, now proves that after a finite
extension \(K\),
\[
 \rho_A(\Gamma_K)\text{ is compact open in }M(\mathbf A_f).
 \tag{E.2.25d}
\]
This is the full restricted-product topology
\cite[Theorem A(ii), author p.~2; Theorem 5.3 and
proof, author pp.~11--12]{GenCadoretMoonen2020}.
The source uses \(H^1\); inverse dual transfers its
actual representation and integral comparison to our
\(H_1\) convention. The multiplier exponent still has
absolute value one, as required above. Further finite
enlargement of \(K\) preserves (E.2.25d), because the
new compact image has finite index in the old one and
is therefore open.

\paragraph{2. Supported conjugation commutators.}
Enlarge the finite set of geometric classes to its
finite \(\Gamma_\Q\)-orbit, choose representatives
\(A_i\), and let actual complex conjugation \(c\)
permute them by \(\pi\). Choose actual isogenies
\(\phi_i:{}^c A_i\to A_{\pi(i)}\). Take \(K/\Q\)
finite Galois defining all models, endomorphisms,
polarizations, these isogenies, and the data of
(E.2.25d). Put
\[
 C_i=(\phi_i)_*c_i:V_i\longrightarrow V_{\pi(i)},
\]
where \(c_i\) is the actual anti-complex-linear
map on homology. Functoriality of Betti--\'etale
comparison gives the exact identity
\[
 C_i\rho_i(c\gamma c)C_i^{-1}
          =\rho_{\pi(i)}(\gamma)
                  \qquad(\gamma\in\Gamma_K).
 \tag{E.2.25e}
\]
These maps identify \(H_i\) with \(H_{\pi(i)}\).
Indeed \(\phi_i\) is a rational Hodge isomorphism,
and \(c_i\) changes the Hodge homomorphism \(h(z)\)
to \(h(\bar z)\); the rational algebraic group
generated by their images is the same. This also
proves the identification when the selected
polarizations are different: their change lies in
the labelled Rosati endomorphism space centralized
by \(H_i\).

Set \(G=\rho(\Gamma_K)\), and let \(\theta\) be
the automorphism induced by \(\gamma\mapsto c\gamma c\)
through (E.2.25e). The compact group
\[
 N=\rho(\Gamma_K\cap\Gamma_{\Q^{\mathrm{cm}}})
\]
is closed and normal in \(G\). It is contained in
\(H(\mathbf A_f)\), since
\(\Q^{\mathrm{ab}}\subset\Q^{\mathrm{cm}}\) and
the multiplier is cyclotomic. Actual \(c\) is
central on \(\Q^{\mathrm{cm}}\). Hence
\[
 [\gamma,c]\in\Gamma_K\cap\Gamma_{\Q^{\mathrm{cm}}},
 \qquad g\theta(g)^{-1}\in N
                \quad(g=\rho(\gamma)\in G).
 \tag{E.2.25f}
\]
By (E.2.25d), \(G\) contains independently supported
open subgroups in every local factor of \(H\), and
the full standard integral factors outside a finite
set of primes. We always use these supported elements.

If \(\theta\) swaps two classes or two local factors,
take \(g\) supported on the first. The element
\(g\theta(g)^{-1}\) has support in the two factors.
Its commutator with \(h\) supported on the first
is \([g,h]\) in that factor and the identity in
the other. Normality of \(N\) supplies this
independently supported derived subgroup.

For a fixed class, we verify nontriviality on every
fixed local simple factor. When \(\End^0(A_i)=\Q\),
the centralizer of the standard symplectic group
is scalar. If \(\operatorname{Ad}(C_i)\) were the
identity, \(C_i\) would be a scalar and thus
complex-linear, contrary to its anti-complex-linearity.
A noncentral rational matrix stays noncentral at
every completion and at all but finitely many
residue characteristics.

For real multiplication let \(\kappa\) be the
automorphism of \(F\) induced by \(C_i\). If it is
nontrivial, it swaps the two \(\operatorname{SL}_2\)
factors at a split prime. At an inert or ramified
prime it acts nontrivially on the centroid \(F_v\)
of the \(\Q_\ell\)-Lie algebra
\(\mathfrak{sl}_2(F_v)\), so the Lie automorphism
is nontrivial. This Lie algebra is \(\Q_\ell\)-simple:
over an algebraic closure its simple summands are
transitively permuted by the local Galois group,
and an ideal is a stable sum of such summands.

If \(\kappa=1\), the matrix \(C_i\) is \(F\)-linear.
On each real \(F\)-component it is an invertible
anti-complex-linear map of a real plane. Such a
matrix has trace zero and square a positive scalar.
Consequently
\(\operatorname{Tr}_F(C_i)=0\) and \(C_i^2=d_i\)
with \(d_i\) totally positive. It is nonscalar
in every \(F_v\), because each \(F\to F_v\) is
injective. Choose a nonzero off-diagonal entry or
diagonal difference in \(F\); excluding the primes
dividing its norm proves nonscalarity in every
residue-field factor outside a finite set.

For quaternionic multiplication, Skolem--Noether
gives \(u\in B^\times\) realizing the automorphism
of the labelled left algebra induced by \(C_i\).
Then \(D_i=u^{-1}C_i\) centralizes \(B\) and is
right multiplication by \(b_i\in(B^{\mathrm{op}})^\times\).
Left multiplication by \(u\) commutes with \(H_i\),
so \(\operatorname{Ad}(C_i)|_{H_i}
=\operatorname{Ad}(D_i)|_{H_i}\). It is a Hodge
endomorphism, so \(D_i\) remains anti-complex-linear.
Under \(B_\R=M_2(\R)\), its right action on the
standard real plane has trace zero and negative
determinant. Thus \(b_i\) is noncentral in every
completion, and in all but finitely many split
matrix reductions. The local Lie algebras of
\(\operatorname{SL}_1(B^{\mathrm{op}})\) are forms
of \(\mathfrak{sl}_2\) and are simple in
characteristic zero.

\paragraph{3. Actual local neighborhoods and the full prime product.}
Consider a fixed local simple factor \(J/\Q_\ell\)
and the supported closed subgroup
\(R=N\cap J(\Q_\ell)\), normal under an open
subgroup \(U\). The preceding argument makes
\(\theta\) nontrivial on \(\Lie J\). For
\(X\in\Lie J\), use the actual Cayley curve
\[
 q_X(t)=(1+tX)(1-tX)^{-1},\qquad q'_X(0)=2X.
 \tag{E.2.25g}
\]
For a symplectic \(X\), the matrices \(1+tX\)
and \(1-tX\) have the same pulled-back alternating
form. For trace-zero \(2\) by \(2\) matrices
their determinants agree, and for a reduced-trace-zero
quaternion their reduced norms agree. Thus (E.2.25g)
is a \(J\)-valued rational analytic curve also at
division primes, and lies in \(U\) on a small ball.

The curve \(q_X(t)\theta(q_X(t))^{-1}\) lies in
\(R\) and has derivative \(2(X-\theta X)\).
The derivatives of its conjugates by \(U\) span
\(\Lie J\). Their span is stable under
\(\operatorname{Ad}U\), hence under the Lie
algebra by differentiating \(\operatorname{Ad}q_Y(t)\);
it is a nonzero ideal in a simple Lie algebra.
Choose finitely many curves with basis derivatives.
Their product has invertible Jacobian at zero.
In local coordinates it is \(J_0x+P(x)\), with
\(P\) of order at least two. On a small ball
\(J_0^{-1}P\) is strictly contracting; iterating
\(x\mapsto y-J_0^{-1}P(x)\) solves
\(J_0x+P(x)=J_0y\) for every sufficiently small
\(y\). Its image contains an actual open
neighborhood of the identity, entirely in \(R\).

The same argument proves that the closed derived
subgroup of an open local semisimple group is open.
Use the curves
\(gq_X(t)g^{-1}q_X(t)^{-1}\), whose derivatives are
\(2(\operatorname{Ad}_g-1)X\). A linear functional
annihilating them annihilates every \([Y,X]\),
by differentiating \(g=q_Y(s)\); semisimplicity
therefore makes their span the entire Lie algebra.
The same product and contraction give an open
neighborhood. For a product, use the supported
curves in its separate simple factors.

Now exclude the finite set containing \(2,3\),
field and quaternion discriminants, denominators
of the fixed maps, noncentrality exceptions, and
the exceptions in (E.2.25d). The supported integral
factors are
\(\operatorname{Sp}_4(\Z_\ell)\),
\(\operatorname{SL}_2(\Z_\ell)\), and
\(\operatorname{SL}_2(\mathcal O_v)\), with
\(\mathcal O_v/\Z_\ell\) unramified quadratic.

For \(\operatorname{Sp}_4\), the reduction of
\(R\) is normal and contains a noncentral
\(g\theta(g)^{-1}\). Otherwise \(\theta\) differs
from the identity by a central character of the
perfect finite symplectic group. That character
is trivial, forcing \(\theta=1\); its standard
centralizer is scalar, contradicting the chosen
noncentral reduction of \(C_i\). The normal
subgroup is therefore the whole finite group
\cite[Lemma 2.3(4), p.~6]{GenMayleWang2024}.
The symplectic lifting lemma then gives
\(R=\operatorname{Sp}_4(\Z_\ell)\)
\cite[Lemma 2.4, p.~6]{GenLombardoSurfaces2015}.

For \(\operatorname{SL}_2(k_v)\), the elementary
simplicity proof in E.2.24 works over every finite
field of characteristic at least five: the faithful
doubly transitive projective-line action, the normal
abelian root group in a point stabilizer, and
perfectness prove simplicity of the projective
quotient. Perfectness rules out a remaining
quotient of order two. A nontrivial semilinear
field automorphism cannot act trivially: traces
of \(\operatorname{SL}_2(k_v)\) cover \(k_v\),
so it would fix every field element; the remaining
matrix centralizer is scalar. Thus the normal
subgroup generated by the \(\theta\)-commutators
has full reduction. Swapped factors have already
been isolated by the supported commutators.

Here is the complete lifting step for the unramified
quadratic ring. Lift each upper unipotent \(U(a)\)
to \(g\in R\), and write
\(g=U(\widetilde a)(1+\ell B)\pmod{\ell^2}\).
The first error in \(g^\ell\) is the sum of
\(\operatorname{Ad}(U(-j\widetilde a))B\) for
\(0\leq j<\ell\). It has degree at most two in
\(j\), and the sums of \(1,j,j^2\) vanish in
\(k_v\). Hence
\[
 g^\ell=1+\ell aE_{12}\pmod{\ell^2}.
\]
All \(a\), the lower roots and their conjugates
span \(\mathfrak{sl}_2(k_v)\) as an
\(\mathbf F_\ell\)-space. The identity
\[
 (1+\ell^nX)^\ell
       =1+\ell^{n+1}X\pmod{\ell^{n+2}}
            \qquad(n\geq1)
\]
supplies every successive congruence layer.
Closedness gives the full integral group.

These integral groups are perfect. For
\(\operatorname{SL}_2\), use
\([
\operatorname{diag}(2,1/2),U(t)]=U(3t)\)
and the lower-root version, with \(3\) a unit.
For \(\operatorname{Sp}_4\), choose a torus
element with root value \(2\) on a short root,
and value \(4\) on a long root. The commutator
coefficient is respectively \(1\) or \(3\).
The root groups generate the integral symplectic
group by elimination: move a unit coordinate of
a primitive first column to its first position,
clear the others by symplectic elementary operations,
clear its paired column, and reduce to
\(\operatorname{Sp}_2\). The same argument proves
the finite-field perfectness used above.

We have therefore obtained independently supported
open subgroups at the exceptional primes and full
integral factors elsewhere. Their finite supported
products lie in \(N\); closedness puts their full
product in \(N\). This proves actual adelic openness.
Finally, for
\(U_0=\Gamma_K\cap\Gamma_{\Q^{\mathrm{cm}}}\),
\[
 \overline{[U_0,U_0]}
       \subset\Gamma_K\cap\Gamma_{F_{\mathrm{cm}}},
 \qquad
 \rho\bigl(\overline{[U_0,U_0]}\bigr)
       =\overline{[N,N]}.
\]
The inclusion holds because
\(L^{\mathrm{ab}}\Q^{\mathrm{cm}}/
\Q^{\mathrm{cm}}\) is abelian for every CM field
\(L\); the equality follows by compactness and
surjectivity onto \(N\). The local open-derived
argument at the exceptional primes and perfectness
elsewhere prove, with the same supported products,
\[
 \rho\bigl(\Gamma_K\cap\Gamma_{F_{\mathrm{cm}}}\bigr)
             \text{ is open in }H(\mathbf A_f).
 \tag{E.2.25h}
\]
Further finite enlargement of \(K\) preserves
this conclusion. Only the displayed derived-subgroup
inclusion was used.

\paragraph{4. The exact rational point and integral coset.}
Every \(H_i\) satisfies the proved criterion D.4.1:
it is simply connected semisimple and has noncompact
real points, respectively \(\operatorname{Sp}_4(\R)\),
\(\operatorname{SL}_2(\R)^2\), or
\(\operatorname{SL}_2(\R)\) in the indefinite
quaternion case. Thus
\[
 H_i(\Q)\text{ is dense in }H_i(\mathbf A_f).
 \tag{E.2.25i}
\]
Fix \(\sigma\in\Gamma_{F_{\mathrm{cm}}}\) and a
conjugate target \({}^{\sigma}A_i\). E.2.11 supplies
an actual rational tensor lift of the prescribed
all-CM and Artin comparison. Its block selected by
the finite-model Artin idempotent is
\[
 T_0:V_i\longrightarrow V_{{}^{\sigma}A_i}.
\]
It intertwines every labelled endomorphism and
polarization and has Tate map \(1\). The actual
completed transport \(\rho_i(\sigma)\) has these
same properties. The centralizers and multiplier-one
kernels in (E.2.25b)--(E.2.25c) therefore give
\[
 a=T_0^{-1}\rho_i(\sigma)\in H_i(\mathbf A_f).
 \tag{E.2.25j}
\]
The underlying group map also has this exact image:
the all-CM kernel of E.2.11 has connected semisimple
image in the finite geometric Mumford--Tate group;
its quotient is a quotient of the geometric CM
protorus. It therefore contains the derived group,
and semisimplicity gives the reverse inclusion.
The image is precisely \(H_i\).

Let \(U\) be an open subgroup of \(H_i(\mathbf A_f)\)
contained in the integral lattice stabilizer.
Choose \(h\in H_i(\Q)\cap aU\) by (E.2.25i).
Then
\[
 T=T_0h\in\rho_i(\sigma)U.
 \tag{E.2.25k}
\]
It is rational, and its completion is an integral
isomorphism between the specified source and target
lattices. Its entries and inverse entries belong to
\(\Q\cap\widehat\Z=\Z\). Thus it is an actual
integral unimodular isomorphism. This is the required
integral intersection of the frozen-target coset;
the rational point used is the established point
\(T_0\) of E.2.11.

\paragraph{5. One simultaneous quartic swap.}
Enumerate one \(A_n\) in every geometric isogeny
class of non-CM simple abelian varieties of dimension
at most two, using the fixed algebraic-model ordering
from E.2.17. Fix bases without an unselected
topological choice as follows. Choose the first
algebraic basis of \(H^0(A_n,\Omega^1)\) in the
same data ordering, also for each conjugate model.
Integration identifies \(\Lambda_{A_n}\) with a
full lattice in \(\C^{\dim A_n}\). Among its
ordered integral bases minimize the sum of the
squared Euclidean norms of their period vectors,
then the lexicographic tuple of all their real
coordinates. Any first available basis supplies
a finite bound for this sum, and only finitely
many ordered lattice bases lie below that bound.
The minimum and its lexicographic tie-break
therefore exist and give a unique basis. No computational
decidability is asserted.

Start with the nonempty compact fibre
\[
 C_0=\{\sigma\in\Gamma_{F_{\mathrm{cm}}}:
                         \sigma|_P=(12)(34)\}
\]
of E.2.6. Choose increasing finite Galois fields
\(K_n\) containing \(P\), the first \(n\) model
data, their finite class orbits and all endomorphisms
and isogenies required in (E.2.25h), and the \(n\)-th
field in the fixed Galois exhaustion of \(\Qb\).
Make each choice by the first suitable item in the
fixed countable data ordering. Partition the
preceding compact fibre by its restriction to
\(K_n\), and choose the first nonempty cylinder
\(C'_n\). All conjugate targets and their bases
are now fixed.

For any \(\sigma_n\in C'_n\), the stabilizer under
right multiplication is exactly
\[
 \mathcal K_n=\Gamma_{F_{\mathrm{cm}}}\cap\Gamma_{K_n}
                   \cap\bigcap_{i<n}\ker\rho_{A_i},
 \qquad C'_n=\sigma_n\mathcal K_n.
 \tag{E.2.25l}
\]
Here \(\mathcal K_n\) is merely notation for this
Galois subgroup. The composition identity
\(\rho_{A_i}(\sigma_n k)
=\rho_{A_i}(\sigma_n)\rho_{A_i}(k)\), valid for
\(k\in\Gamma_{K_n}\), proves the equality of
the stabilizer and the fibre in (E.2.25l).
Write \(J_n\) for the algebraic group \(H_i\)
attached to \(A_n\), to avoid confusing it with
a Galois subgroup. Openness (E.2.25h) for the first
\(n\) factors implies that
\(U_n=\rho_{A_n}(\mathcal K_n)\) is open in
\(J_n(\mathbf A_f)\): its product congruence
neighborhood contains all sufficiently deep
last-factor elements with the preceding components
exactly the identity. The possible next matrices are
\[
 \rho_{A_n}(\sigma_n)U_n.
 \tag{E.2.25m}
\]
By (E.2.25k) this coset contains an actual integral
lattice isomorphism. Choose the first such integer
matrix in entry-height and lexicographic order,
and impose it to obtain \(C_n\subset C'_n\).
These compact nonempty sets are nested. Their
intersection is nonempty and contains only one
\(\sigma_*\), because the prescribed restrictions
to the \(K_n\) exhaust \(\Qb\). This is a
specified Galois element retaining the quartic
action, with integral actual transport on every
\(A_n\).

For this same \(\sigma_*\), every CM factor has
the prescribed integral algebraic comparison of
E.2.6. On a product \(B\) of representatives and
CM factors take the direct sum map. For any
geometric isogeny \(q:B\to A\), put
\[
 L_A=q_*^{-1}\Lambda_A,
 \qquad \Lambda_B\subset L_A\subset\tfrac1m\Lambda_B,
 \qquad L_A/\Lambda_B=\ker(q)\subset B[m](\C)
\]
for a suitable \(m\). Actual completed transport
sends this finite kernel to \(\ker({}^{\sigma_*}q)\).
Its integral source map therefore carries \(L_A\)
onto \(L_{{}^{\sigma_*}A}\), giving the actual
integral map
\[
 T_A=({}^{\sigma_*}q)_*T_Bq_*^{-1}:
       \Lambda_A\xrightarrow{\sim}\Lambda_{{}^{\sigma_*}A}.
\]
Completion injectivity as proved in E.2.16 gives
independence of \(q\) and of the decomposition,
and every correspondence, tensor, exterior-power
and projector equation. Cohomology uses inverse
duals, and the common Tate map is \(1\). The
finite-model Artin blocks have the actual maps
\(e_u\mapsto e_{\sigma_*u}\), so their module
maps give precisely the descent of E.2.3.
This proves (E.2.25a) on the whole specified diagram.

The comparison scope includes every abelian surface
over \(\Qb\). It does not supply integrality for
the additional simple non-CM factors of dimension
at least three or extend the comparison to every
pure motive. E.2.22 remains the exact cover-tower
test for extension of this same \(\sigma_*\) to
the entire abelian diagram.

The exact finite certificate
\path{code/verify_surface_groups.py} checks the
Cayley identities, symplectic root commutators,
the full Lie-bracket span and scalar centralizer,
the real multiplication trace-zero identities,
the quaternion reduced-norm identity, and the
anti-Hodge real-plane matrix. The neighborhood,
prime-product, strong-approximation and simultaneous
comparison assertions are established by the proof
above.
\end{proof}

\begin{result}{E.2.26}{One integral comparison for all threefold and CM geometric factors}{Proved}
Let \(\mathcal A_{\leq3,\mathrm{CM}}\) consist of all abelian varieties
over \(\Qb\) geometrically isogenous to products of simple factors of
dimension at most three and arbitrary CM abelian varieties. Take the
full integral diagram of E.2.16: every geometric isogeny lattice,
products, algebraic correspondences, tensor and exterior-power
constructions, specified projectors, finite-model Artin module maps,
and all Artin objects. With \(P\), its labelled roots, and the quartic
pair as in E.2.17 and E.2.23, there is a specified
\(\sigma_*\in\Gamma_{F_{\mathrm{cm}}}\) for which
\[
 \sigma_*|_P=(12)(34),\qquad
 \rho_A(\sigma_*)\Lambda_A=\Lambda_{{}^{\sigma_*}A}
       \quad(A\in\mathcal A_{\leq3,\mathrm{CM}}).
 \tag{E.2.26a}
\]
The completed maps in (E.2.26a) are the completions of one integral
tensor comparison on this entire diagram, with Tate map \(1\).
Its Artin module maps give the comparison of the two specified real
embeddings by E.2.3. This statement concerns the plain integral
diagram; the labelled real-Frobenius witness remains E.2.7.
\end{result}
\begin{proof}
Throughout, \(\Lambda_A=H_1(A(\C),\Z)\) and the Tate modules
are covariant homology Tate modules. Cohomology uses inverse duals.
We prove the actual finite-family image and then its integral
intersection while retaining every previously selected matrix.

\paragraph{1. Actual groups and arithmetic Mumford--Tate equality.}
For a geometrically simple non-CM threefold \(A\), put
\(V=\Lambda_A\otimes\Q\), choose a labelled polarization \(\psi\),
and write \(H_A=\operatorname{MT}(A)^{\mathrm{der}}\).
The possibilities, as actual groups acting on \(V\), are
\[
\begin{array}{c|c|c}
 \End^0(A)&\text{coefficient dimension}&H_A\\ \hline
 \Q&\dim_{\Q}V=6&\operatorname{Sp}(V,\psi)\\
 F\text{ totally real cubic}&\dim_FV=2
      &\operatorname{Res}_{F/\Q}\operatorname{SL}_F(V)\\
 E\text{ imaginary quadratic}&\dim_EV=3
      &\operatorname{SU}_E(V,h).
\end{array}
\tag{E.2.26b}
\]
In the second row, the unique \(F\)-alternating form \(\Psi\)
satisfies \(\psi=\operatorname{Tr}_{F/\Q}\Psi\). In the third,
fix \(a\in E^\times\) with \(\bar a=-a\); the Hermitian form
\(h\) is determined by \(\psi=\operatorname{Tr}_{E/\Q}(a h)\).
It has signature \((2,1)\), up to interchange. The remaining simple
threefold case has sextic CM endomorphism field. Moonen--Zarhin's
classification and trace-form identifications give exactly these
groups, and the imaginary-quadratic Hodge group is the full unitary
group \(\operatorname U_E(V,h)\)
\cite[Section 2.3 and Proposition 2.4, author p.~9;
Theorem 2.7, author p.~10]{GenMoonenZarhin1999}.
Consequently its full Mumford--Tate group is
\(\operatorname{GU}_E(V,h)\): the product of the Hodge group and
the rational scalar torus is connected of dimension ten, the
dimension of \(\operatorname{GU}\). The lower-dimensional factors
retain the exact groups and forms of E.2.25.

Every group \(H_A\) in (E.2.26b) is semisimple simply connected.
Its \(\Q\)-simple factors have noncompact real points:
\(\operatorname{Sp}_6(\R)\), \(\operatorname{SL}_2(\R)^3\),
or \(\operatorname{SU}(2,1)\), respectively.

We verify the arithmetic equality before applying an adelic image
theorem. Under the actual Betti--\'{e}tale comparison it is
\[
 G_\ell(A)^0=\operatorname{MT}(A)_{\Q_\ell}.
 \tag{E.2.26c}
\]
Tankeev's original proof in prime dimension different from seven
gives the semisimple Lie dimensions needed here: eight for the
imaginary-quadratic threefold, nine for the real-cubic case, and
twenty-one for \(\End^0(A)=\Q\). The natural inclusion into the
Mumford--Tate Lie algebra is used explicitly in that proof. Its full
unitary equality includes the hypothesis \(c_\ell^0\ne0\), which
we now verify
\cite[Lemma 3 and proof, pp.~532--535; type-I cases,
pp.~535--539]{GenTankeev1988}.

Commelin's exact center theorem, under the same comparison, is
\[
 Z(G_\ell(A)^0)^0
     = Z(\operatorname{MT}(A))^0_{\Q_\ell}.
 \tag{E.2.26d}
\]
Its proof has the following maps. Faltings' endomorphism theorem
identifies the centralizer of the connected arithmetic image with
the geometric endomorphism algebra; Commelin's Lemma 4.15 then
places its connected center inside the Mumford--Tate center. A
faithful representation \(N\) of
\(\operatorname{MT}(A)^{\mathrm{ab}}\) is an abelian CM motive,
by the full Hodge realization theorem and Section 4.16. The CM
Mumford--Tate theorem gives its connected arithmetic image equal
to that torus. Thus the actual projection
\(G_\ell(A)^0\to\operatorname{MT}(A)^{\mathrm{ab}}_{\Q_\ell}\)
is surjective and proves (E.2.26d)
\cite[Theorem 4.13, Lemmas 4.15--4.18 and proof,
pp.~658--659]{GenCommelin2019}.
The source takes a finitely generated definition extension for
\(N\). Its algebraic part over the number field is finite; the
extension is regular over that algebraic part, so restriction of
absolute Galois groups is surjective onto its Galois group.
Consequently the image in the original arithmetic group changes
by finite index and its identity component is unchanged.
In the unitary threefold case (E.2.26d) gives center dimension two,
so \(c_\ell^0\) has dimension one. The qualification in Tankeev's
unitary argument is satisfied. In all three rows of (E.2.26b),
the natural inclusion has equal Lie dimension and hence is equality
of connected groups. This proves (E.2.26c); Lombardo's explicit
prime-dimensional statement supplies a separate source check
\cite[Theorem 2.18, p.~6]{GenLombardoNonsimple2015}.

There is also a direct arithmetic check of the unitary center.
After a model extension defining the \(E\)-action and containing
\(E\), Fit\'e--Goodman construct a CM elliptic curve \(B\) with
\[
 \left(\bigwedge_{E_\lambda}^3 V_\lambda(A)
       \oplus\bigwedge_{E_\lambda}^3\overline{V_\lambda(A)}\right)(-1)
       \simeq V_\ell(B)\otimes_{\Q_\ell}E_\lambda.
 \tag{E.2.26e}
\]
Here the wedge is the determinant over the coefficient field
\(E_\lambda\); the bar denotes the coefficient-conjugate
representation defined in that source, including the isomorphism
between the conjugate coefficient completions
\cite[Theorem 1.1(i), p.~2, and Theorem 3.1 with proof,
pp.~9--10]{GenFiteGoodman2026}.
The determinant characters are therefore the two CM elliptic
characters multiplied by the common cyclotomic character.
The torus map is
\(t\mapsto t\operatorname{Norm}_{E/\Q}(t)\); on split
characters its matrix is
\(\left(\begin{smallmatrix}2&1\\1&2\end{smallmatrix}\right)\),
of determinant three. It is an isogeny of a two-dimensional torus,
which verifies the center dimension directly. This use of
(E.2.26e) is an equality of the stated arithmetic representations.

For every finite family of these factors, arithmetic
Mumford--Tate equality holds for their product by Commelin's
product theorem
\cite[Theorem 10.5, p.~676]{GenCommelin2019}.
We use its actual joint group, whose shared factors are handled
next.

\paragraph{2. The joint derived group and an exact rational splitting.}
For a finite family let
\(M=\operatorname{MT}(\bigoplus_iV_i)\) and \(H=M^{\mathrm{der}}\).
The projections \(M\to M_i\) and \(H\to H_i\) are surjective,
and the faithful sum representation embeds \(H\) as a closed
connected subgroup of \(\prod_iH_i\).

Here is an explicit proof that \(H\) is simply connected.
Over \(\Qb\), write its universal cover as
\(H^{\mathrm{sc}}=\prod_jS_j\). For an absolutely simple
factor \(T\) of a target \(H_i\), the images of the \(S_j\)
in \(T\) are commuting connected normal subgroups. Surjectivity
implies that exactly one is nontrivial. Its map \(S_j\to T\)
is a central isogeny, hence an isomorphism because \(T\) is
simply connected. Every \(S_j\) reaches at least one such
target, since otherwise the faithful embedding of \(H\) would
kill its positive-dimensional image. An element of the kernel
of \(H^{\mathrm{sc}}\to H\) therefore has trivial coordinate
in every \(S_j\), using its isomorphic target projection.
The kernel is trivial, and this descends to \(\Q\).

Every \(\Q\)-simple factor \(S\) of \(H\) maps nontrivially
to some \(\Q\)-simple \(H_i\). The image is connected normal,
so it is all of \(H_i\); \(\Q\)-simplicity gives finite
kernel, and the simply connected target makes the map an
isomorphism. Every such \(S\) therefore has noncompact real
points and is one of the forms in (E.2.26b) or E.2.25.

For the joint derived groups of the first \(n\) generators,
the forgetting map \(p_n:H_n\to H_{n-1}\) is surjective.
Over \(\Qb\) it discards a subset of simple factors and maps
the remaining ones isomorphically onto the target factors.
The kernel subset and its complement are Galois stable, giving
the actual \(\Q\)-decomposition
\[
 H_n=L_n\times C_n,\qquad L_n=\ker p_n,
 \qquad p_n|_{C_n}:C_n\xrightarrow{\sim}H_{n-1}.
 \tag{E.2.26f}
\]
Let \(j_n:H_{n-1}\to H_n\) be its inverse section.
The group \(L_n\) is connected semisimple simply connected,
and all its \(\Q\)-simple factors have noncompact real points.
The possibility \(L_n=1\) is included. In particular (E.2.26f)
retains all shared diagonal factors in the precise way needed
to keep earlier matrices fixed.

\paragraph{3. Actual prime-supported derived subgroups.}
For a finite product choose a number field \(k\) defining
the models and endomorphisms, enlarged so that every
\(G_\ell\) is connected. The common connectedness field is
the one recalled in Cadoret--Moonen's Section 1.1.
Their Theorem A(i), with Section 1.3, applies to (E.2.26c):
each local image is open in \(M(\Z_\ell)\), and for all
but finitely many \(\ell\) it contains
\([M(\Z_\ell),M(\Z_\ell)]\)
\cite[Theorem A(i), Sections 1.1--1.3]{GenCadoretMoonen2020}.
The integral models are the closures in the actual Betti
lattice; outside a finite set they are the smooth reductive
models of the groups above.

Serre's independence theorem, applied to the full family of
Tate modules of this product, gives a finite extension \(k'/k\)
with the actual equality
\[
 \rho(\Gamma_{k'})=\prod_\ell\rho_\ell(\Gamma_{k'}).
 \tag{E.2.26g}
\]
This is the independence defined on p.~542 and proved for
these Tate modules in Theorem 1 and Section 3.1
\cite[pp.~542--545]{GenSerreIndependence2013}.
Apply the integral Mumford--Tate statement again over \(k'\).
For a further finite Galois field \(K/\Q\), the image of
\(\Gamma_K\) is a closed finite-index subgroup of the compact
product (E.2.26g), hence is open. It contains a product of open
local subgroups equal to the full local image outside a finite
set. At a good prime \(\ell\geq7\), the perfectness proved
below gives
\[
 H(\Z_\ell)=[H(\Z_\ell),H(\Z_\ell)]
       \subset[M(\Z_\ell),M(\Z_\ell)].
\]
At any remaining prime, intersect a local open subgroup with
\(H(\Q_\ell)\). We obtain the actual containment
\[
 W\subset\rho(\Gamma_K),\qquad
 W\subset H(\mathbf A_f)\text{ compact open},
 \quad W_\ell=H(\Z_\ell)\text{ for almost all }\ell.
 \tag{E.2.26h}
\]
Every element of \(W\) supported at one prime and one local
simple factor has an actual Galois preimage.

\paragraph{4. Finite groups and the full integral unitary lifting.}
The good reductions are
\(\operatorname{SL}_2(\mathbf F_q)\),
\(\operatorname{Sp}_4(\mathbf F_\ell)\),
\(\operatorname{Sp}_6(\mathbf F_\ell)\),
\(\operatorname{SL}_3(\mathbf F_\ell)\), or
\(\operatorname{SU}_3(\mathbf F_\ell)\).
For \(\ell\geq7\) they are perfect and their projective
quotients are simple. The opened matrix proofs are Taylor's
Theorems 4.4--4.5, 8.7--8.8, and 10.15; the last explicitly
proves unitary-transvection generation and perfectness
\cite[pp.~22--23, 72--74, 122--123]{GenTaylorClassical2021}.
A normal subgroup containing a noncentral element surjects
onto the simple projective quotient. The remaining quotient
of the full group is then a quotient of its center, hence
abelian, and perfectness makes it trivial.

We prove the integral lifting statement for \(\operatorname{SU}_3\).
Let \(\mathcal O/\Z_\ell\) be unramified quadratic, choose
a trace-zero unit \(\vartheta\), and use
\[
 J=\begin{pmatrix}0&0&1\\0&1&0\\1&0&0\end{pmatrix},\qquad
 Q(a,b)=\begin{pmatrix}1&-\bar a&b\\0&1&a\\0&0&1\end{pmatrix},
 \quad a=x+y\vartheta,\quad
 b=-a\bar a/2+z\vartheta.
 \tag{E.2.26i}
\]
The exact identities are
\(\det Q=1\), \(\bar Q^{\,t}JQ=J\), and
\(Q(a,b)^{-1}=Q(-a,\bar b)\). Thus both \(Q\) and
its inverse have polynomial degree at most two, and
\(\operatorname{Ad}Q\) has degree at most four.
The matrix \(u=1+N\), \(N=\vartheta E_{13}\), is an
integral unitary transvection with \(N^2=0\).

Suppose \(R\) is a closed subgroup of
\(\operatorname{SU}_3(\Z_\ell)\) surjecting modulo \(\ell\).
Choose \(g\in R\) lifting \(u\bmod\ell\), and write
\(g=u(1+\ell B)\bmod\ell^2\). Then
\[
 g^\ell=u^\ell\left(1+\ell
       \sum_{j=0}^{\ell-1}\operatorname{Ad}(u^{-j})B\right)
       \pmod{\ell^2}.
\]
Here
\(\operatorname{Ad}(u^{-j})B
 =B+j(BN-NB)-j^2NBN\).
The sums of \(1,j,j^2\) vanish modulo \(\ell\geq5\),
so \(g^\ell=1+\ell N\bmod\ell^2\).
The first congruence layer of \(R\) contains \(N\), and
the span of its conjugates is stable under the finite unitary
group. For \(\ell\geq7\), polynomial interpolation in
\(x,y,z\) in (E.2.26i) makes this subspace stable under
bracketing with the six root derivatives
\[
\begin{aligned}
 P_x&=-E_{12}+E_{23},&
 P_y&=\vartheta(E_{12}+E_{23}),&P_z&=\vartheta E_{13},\\
 M_x&=E_{21}-E_{32},&
 M_y&=\vartheta(E_{21}+E_{32}),&M_z&=\vartheta E_{31}.
\end{aligned}
\]
Their brackets
\[
 [P_x,M_x]=\operatorname{diag}(-1,0,1),\qquad
 [P_x,M_y]=\vartheta\operatorname{diag}(-1,2,-1)
\]
together with the six root matrices form a basis of
\(\mathfrak{su}_3(\mathbf F_\ell)\).
The conjugate span is therefore a nonzero Lie ideal.
After quadratic scalar extension this Lie algebra is
\(\mathfrak{sl}_3\), simple in characteristic greater
than three, so the first layer is full. Smoothness identifies
all successive congruence layers with that Lie algebra, and
\[
 (1+\ell^mX)^\ell
      =1+\ell^{m+1}X\pmod{\ell^{m+2}},\qquad m\geq1,
 \tag{E.2.26j}
\]
gives each next layer. Hence \(R\) surjects modulo every
\(\ell^m\), and closedness gives the full integral group.

For \(\operatorname{SL}_3\), use the elementary roots
\(1+tE_{ij}\); their conjugate logarithms span
\(\mathfrak{sl}_3\), and the same calculation applies.
For \(\operatorname{SL}_2\) over an unramified local
field \(F_v\), use transvections with every parameter in
its residue field \(\mathbf F_q\). The first layer
contains all residue multiples of \(E_{12}\) and their
conjugates, hence \(\mathfrak{sl}_2(\mathbf F_q)\) as
an \(\mathbf F_\ell\)-space; (E.2.26j) gives all layers.
For \(\operatorname{Sp}_6\), the logarithms
\(v\mapsto\psi(v,u)u\) span the Lie algebra: under
\(\mathfrak{sp}(V)\simeq\operatorname{Sym}^2(V)\)
they are the squares \(u^2\), which span because
\(2uv=(u+v)^2-u^2-v^2\). The same power calculation
gives all layers, including the already used
\(\operatorname{Sp}_4\) case.
The closed commutator subgroup of each good integral group
surjects modulo \(\ell\) by finite perfectness; the lifting
proof just given makes it the whole integral group. This
establishes the integral perfectness used in (E.2.26h).

\paragraph{5. Complex conjugation on every actual factor.}
Enlarge the finite set of geometric isogeny classes to its
finite \(\Gamma_\Q\)-orbit. For the fixed actual complex
conjugation \(c\), choose isogenies
\(\phi_i:{}^c A_i\to A_{\pi(i)}\), and let
\(C_i=\phi_{i,*}c_i\) on rational Betti homology.
Their block sum \(C\) gives the actual intertwining map;
choose a finite Galois \(K/\Q\) defining every map.
For \(\gamma\in\Gamma_K\), naturality gives
\[
 C_i\rho_i(c\gamma c)C_i^{-1}
         =\rho_{\pi(i)}(\gamma),\qquad
 C h(z)C^{-1}=h(\bar z).
 \tag{E.2.26k}
\]
Thus \(\theta=\operatorname{Ad}(C)\) acts on \(M\)
and \(H\). Its square is trivial there because \(C^2\)
is block diagonal with algebraic endomorphisms, which
centralize \(M\).

On a \(\theta\)-fixed \(\Q\)-simple factor \(S\)
of \(H\), this automorphism cannot be the identity.
Indeed the actual quotient \(M\to S^{\mathrm{ad}}\)
would turn (E.2.26k) into
\(h_{\mathrm{ad}}(z)=h_{\mathrm{ad}}(\bar z)\).
On the unit circle this gives
\(h_{\mathrm{ad}}(z)=h_{\mathrm{ad}}(z)^{-1}\), so
connectedness makes the homomorphism trivial there.
It is also trivial on the real weight torus, since
weight is central. The resulting trivial adjoint
Hodge homomorphism contradicts the fact that this
nontrivial quotient of \(M\) is generated as a
\(\Q\)-group by its image.

Write \(S=\operatorname{Res}_{F/\Q}T\), with \(T\)
absolutely simple over its centroid \(F\). The action
permutes the places of \(F\) over \(\ell\). On a
fixed place, a nontrivial action on \(F\) stays
nontrivial on the centroid \(F_v\). When it fixes
\(F\), identity after the faithful scalar extension
\(F\to F_v\) would force identity over \(F\).
Thus each fixed local simple factor has nontrivial
automorphism; the other local factors occur in exchanged
pairs. Away from finitely many primes all these maps
and models are integral smooth and the fixed-factor
Lie automorphisms remain nontrivial after reduction.

For such a finite factor some
\(g\theta(g)^{-1}\) is noncentral. If every one were
central, this assignment would be a homomorphism from
a perfect group to its center, hence trivial. Then
\(\theta\) would fix every finite point. Interpolation
in the root matrices above, of degree at most two,
would make its differential fix their Lie generators,
contradicting the nontrivial reduced differential.
When the centroid automorphism is nontrivial, using
every residue parameter in those root matrices
would also force that residue-field automorphism
to be trivial. This proves the required noncentral
commutator in each case.

\paragraph{6. Actual supported commutators over \(F_{\mathrm{cm}}\).}
Put \(U_0=\Gamma_K\cap\Gamma_{\Q^{\mathrm{cm}}}\) and
\(N=\rho(U_0)\). This is a compact normal subgroup
of \(\rho(\Gamma_K)\). Actual \(c\) is central on
\(\Q^{\mathrm{cm}}\), so
\[
 [\gamma,c]\in U_0,
 \qquad g\theta(g)^{-1}\in N
       \quad(\gamma\in\Gamma_K, g=\rho(\gamma)).
 \tag{E.2.26l}
\]
We construct subgroups of \(H(\mathbf A_f)\)
inside \(N\); this accommodates any finite action
on the joint central CM torus.

At a good prime, the independently supported full
factors in (E.2.26h), the noncentral commutator in
(E.2.26l), normality, finite projective simplicity,
and integral lifting prove that \(N\) contains
the whole supported fixed factor. For a pair of
factors exchanged by \(\theta\), the first
commutator has two coordinates. A second commutator
with an element supported in the first factor has
only that coordinate. Perfectness makes the closed
group generated by these second commutators the
full factor, and the second factor follows likewise.

At an exceptional prime, use a supported open
subgroup supplied by (E.2.26h). On a fixed local
factor the actual curves
\[
 t\longmapsto\exp(tX)\theta(\exp(tX))^{-1}
\]
have tangent vectors \((1-\theta)X\). For small
\(t\) the matrix exponentials converge and lie
in the actual group: they preserve its defining
bilinear or Hermitian form and have determinant
one for trace-zero \(X\). For a quaternion factor
the same assertions follow from reduced trace
and norm after scalar extension to a splitting
field. A nonzero such tangent and its conjugates
under a local open subgroup span a nonzero ideal
of the local simple Lie algebra, hence all of it.
Choose finitely many resulting actual curves
whose tangents form a basis. Their product map
has invertible differential; the explicit
\(p\)-adic contraction argument of E.2.25 gives
an actual neighborhood of the identity in \(N\)
on this supported factor. For exchanged factors,
second commutators have tangents
\((1-\operatorname{Ad}h)X\). The zero center of
the simple Lie algebra makes their conjugate
span full, giving the same conclusion.

All these constructions have individual prime
and factor support. Finite products and closedness
therefore place a compact open subgroup of
\(H(\mathbf A_f)\) inside \(N\).
The closed derived subgroup of \(U_0\) fixes
every \(E^{\mathrm{ab}}\) for CM \(E\), because
\(U_0\) fixes \(E\) and its restriction to
the abelian Galois group kills commutators.
Compactness gives the exact image equality
\[
 \rho\bigl(\overline{[U_0,U_0]}\bigr)
          =\overline{[N,N]}.
 \tag{E.2.26m}
\]
The right side contains all good integral factors
by their proved perfectness. At an exceptional
prime the closed derived subgroup of a local
open subgroup is open: its Lie algebra is
perfect, and a finite product of the actual
commutator curves with spanning tangents gives
the same local contraction proof. Consequently
\[
 \rho(\Gamma_K\cap\Gamma_{F_{\mathrm{cm}}})
       \text{ contains a compact open subgroup of }
       H(\mathbf A_f).
 \tag{E.2.26n}
\]
When the finite family was enlarged to its Galois
orbit, project back to the original joint derived
group. The projection splits over \(\Q\) by
(E.2.26f), so on finite adeles it is an open
direct-product projection. Its compact open
image still has actual preimages in the same
Galois group. Thus (E.2.26n) holds for every
original finite family as well.

\paragraph{7. Integral approximation with every earlier matrix fixed.}
Let \(\sigma\in\Gamma_{F_{\mathrm{cm}}}\) have
fixed conjugate targets for the first \(n\)
generators. E.2.11 supplies a rational tensor
comparison \(T_0\) on the full rational abelian
category with prescribed all-CM comparison
\(\alpha(\sigma)\), Artin action, and Tate map
one. On the joint representation put
\[
 a=T_0^{-1}\rho(\sigma)\in H_n(\mathbf A_f).
 \tag{E.2.26o}
\]
Here the group membership follows from an actual
restriction map. Let \(D\) be the kernel of
the full abelian-to-CM Tannaka restriction.
Its image in \(M_n\) is connected semisimple,
by the exact kernel theorem used in E.2.11,
and the quotient of \(M_n\) by this image is
a quotient of the geometric CM protorus.
It is a torus, so the image contains
\(M_n^{\mathrm{der}}\); semisimplicity gives
the reverse inclusion. Its image is exactly
\(H_n\). The difference fixes the CM point
at every prime, hence its finite projection
belongs to \(H_n(\Q_\ell)\). Its entries
and inverse are integral outside the finitely
many denominators of \(T_0\), proving
(E.2.26o).

Suppose the old integral matrices are fixed as
\(S\). Since \(\rho(\sigma)\) has exactly
those matrices, \(p_n(a)=T_{0,\mathrm{old}}^{-1}S\)
is rational. It satisfies the defining
\(\Q\)-equations of \(H_{n-1}\) in every
completion, so lies in \(H_{n-1}(\Q)\).
Set
\[
 b=j_n(p_n(a))\in H_n(\Q),\qquad
 b^{-1}a\in L_n(\mathbf A_f).
 \tag{E.2.26p}
\]
Freeze the current restriction to a finite
Galois model field \(K_n\). The exact
stabilizer of this cylinder and all old
matrices is
\[
 \mathcal K_n=\Gamma_{F_{\mathrm{cm}}}
       \cap\Gamma_{K_n}\cap\ker(\rho_{\mathrm{old}}).
 \tag{E.2.26q}
\]
Intersect the compact open subgroup in
(E.2.26n) with \(\ker p_n=L_n(\mathbf A_f)\).
Every element in that intersection has an
actual preimage fixing all old Tate modules,
hence a preimage in \(\mathcal K_n\).
It supplies a compact open subgroup
\(U\subset L_n(\mathbf A_f)\) of its image;
shrink it into the integral lattice stabilizer.

D.4.1 applies to \(L_n\), whose simple factors
are simply connected and have noncompact real
points. Thus \(L_n(\Q)\) is dense in
\(L_n(\mathbf A_f)\). Choose
\[
 l\in L_n(\Q)\cap b^{-1}aU,
 \qquad h=bl\in H_n(\Q).
\]
Then \(p_n(h)=p_n(a)\) and \(h=au\)
for some \(u\in U\). The rational comparison
\(T=T_0h\) retains all old matrices and
its completion is \(\rho(\sigma)u\).
Since \(u\) has an actual preimage
\(k\in\mathcal K_n\), that completion
is \(\rho(\sigma k)\). It is an integral
isomorphism at every prime from the source
lattice onto the fixed target lattice.
Both \(T\) and \(T^{-1}\) have rational
entries in \(\widehat\Z\), and
\(\Q\cap\widehat\Z=\Z\) makes them
integer matrices. For \(L_n=1\),
(E.2.26p) already makes \(a\) rational
and the same conclusion holds with \(u=1\).
This proves the exact induction step.

\paragraph{8. A specified simultaneous element and every geometric lattice.}
Enumerate the geometric isogeny classes of
simple non-CM abelian varieties of dimension
at most three, using the fixed algebraic-data
ordering of E.2.25 to choose their models
and polarizations. Use its integral-basis
rule on each model and conjugate: a chosen
algebraic basis of regular differentials
embeds homology as a full lattice in
\(\C^g\); minimize the sum of squared
period norms of an ordered integral basis
and then its real-coordinate lexicographic
tuple. A first available basis supplies a
finite norm bound, and the full lattice
has finitely many vectors below that bound.
Thus the minimizing basis and finite
tie-break exist. This specifies choices
without a computability assertion.

Start from the nonempty compact fibre
\[
 C_0=\{\sigma\in\Gamma_{F_{\mathrm{cm}}}:
                  \sigma|_P=(12)(34)\}.
\]
Let \(K_n\) be the fixed increasing finite
Galois exhaustion of \(\Qb\), enlarged
at step \(n\) to define the selected
models, conjugate targets, and maps above.
In the current fibre choose the first
finite \(K_n\)-cylinder that meets it;
this freezes the conjugate targets. The
preceding step gives at least one integral
next matrix in that cylinder with all old
matrices unchanged. Choose the first such
integer matrix in entry-height and
lexicographic order. Its fibre \(C_n\)
is compact and nonempty. These fibres are
nested. Their intersection is nonempty,
and contains exactly one \(\sigma_*\),
because its specified restrictions to
the \(K_n\) exhaust \(\Qb\).
This is a specified actual Galois element
with the required quartic action and
integral transports on every enumerated
generator.

Every CM abelian variety has the actual
algebraic isomorphism \(\alpha_A\) of
E.2.6 with completion \(\rho_A(\sigma_*)\).
Products carry direct sums. For any
geometric isogeny \(q:B\to A\) from
such a product, put
\[
 L_A=q_*^{-1}\Lambda_A\subset H_1(B(\C),\Q),
 \qquad \Lambda_B\subset L_A\subset\tfrac1m\Lambda_B,
 \quad L_A/\Lambda_B=\ker(q)\subset B[m]
\]
for a suitable \(m\). Actual Galois
transport on the finite kernel agrees
with the selected integral map on \(B\)
modulo \(m\). That rational map
therefore carries \(L_A\) exactly onto
the lattice for \({}^{\sigma_*}q\).
It induces
\[
 T_A=({}^{\sigma_*}q)_*T_Bq_*^{-1}:
     \Lambda_A\xrightarrow{\sim}
     \Lambda_{{}^{\sigma_*}A},
 \tag{E.2.26r}
\]
whose completion is the actual
\'{e}tale transport. This includes every
specified geometric isogeny lattice.

Compare the actual completed equations
for two choices of \(q\), every product,
isogeny, algebraic correspondence, tensor,
exterior power, and specified projector.
Injectivity of rational matrices into
their completions gives independence
and each required compatibility over
\(\Z\), as in E.2.16. The same argument
preserves the labelled polarizations
with the one Tate line fixed. The
Artin maps are the actual permutations
\(x\mapsto\sigma_*x\). On every finite
model the primitive Artin idempotents
select its conjugate blocks, and their
module maps give exactly the descent
of E.2.3. This proves (E.2.26a) on
the full stated diagram.

The remaining extension requires actual
integral transports for simple non-CM
factors of dimension at least four and
then the further objects and morphisms
of the full pure category. The original
full-category place-separation gate
remains open. E.2.22 is still its exact
all-abelian cover-tower test for this
same \(\sigma_*\).

The finite certificate
\path{code/verify_threefold_groups.py}
checks (E.2.26i), the degree-four
adjoint polynomial, the six root and
two Cartan matrices with an exact
minor \(-1\), the degree-two
transvection power calculation, the
rank-twenty-one symplectic transvection
span, and the determinant-three CM
torus map. The arithmetic image,
supported products, approximation,
and countable construction have the
proofs above. Source access is recorded
in the bibliography and work proof:
Tankeev's case proof was read in the
original indexed PDF, while the missing
Theorem 1 display was not claimed as
visually verified.
\end{proof}

